%-%-%-%-%-%-%-%-%-%-%-%-%-%-%-%-%-%-%-%-%-%-%-%-%-%-%-%-%-%-%-%-%-%-%-%-%-%-%
%%% MAIN DOCUMENT %%%

%-%-%-%-%-%-%-%-%-%-%-%-%-%-%-%-%-%-%-%-%-%-%-%-%-%-%-%-%-%-%-%-%-%-%-%-%-%-%

\documentclass[12pt, oneside]{book}
\usepackage{draculatheme}

%%% AESTHETICS %%%
%-%-%-%-%-%-%-%-%-%-%-%-%-%-%-%-%-%-%-%-%-%-%-%-%-%-%-%-%-%-%-%-%-%-%-%-%-%-%


%%% Dimensions and Spacing %%%
\usepackage[margin=1in]{geometry}
\usepackage{setspace}
\linespread{1}
\usepackage{listings}

%%% Define new colors %%%
\usepackage{xcolor}
\definecolor{orangehdx}{rgb}{0.96, 0.51, 0.16}

% Normal colors
\definecolor{xred}{HTML}{BD4242}
\definecolor{xblue}{HTML}{4268BD}
\definecolor{xgreen}{HTML}{52B256}
\definecolor{xpurple}{HTML}{7F52B2}
\definecolor{xorange}{HTML}{FD9337}
\definecolor{xdotted}{HTML}{999999}
\definecolor{xgray}{HTML}{777777}
\definecolor{xcyan}{HTML}{80F5DC}
\definecolor{xpink}{HTML}{F690EA}
\definecolor{xgrayblue}{HTML}{49B095}
\definecolor{xgraycyan}{HTML}{5AA1B9}

% Dark colors
\colorlet{xdarkred}{red!85!black}
\colorlet{xdarkblue}{xblue!85!black}
\colorlet{xdarkgreen}{xgreen!85!black}
\colorlet{xdarkpurple}{xpurple!85!black}
\colorlet{xdarkorange}{xorange!85!black}
\definecolor{xdarkcyan}{HTML}{008B8B}
\colorlet{xdarkgray}{xgray!85!black}

% Very dark colors
\colorlet{xverydarkblue}{xblue!50!black}

% Document-specific colors
\colorlet{normaltextcolor}{black}
\colorlet{figtextcolor}{xblue}

% Enumerated colors
\colorlet{xcol0}{black}
\colorlet{xcol1}{xred}
\colorlet{xcol2}{xblue}
\colorlet{xcol3}{xgreen}
\colorlet{xcol4}{xpurple}
\colorlet{xcol5}{xorange}
\colorlet{xcol6}{xcyan}
\colorlet{xcol7}{xpink!75!black}

% Blue-Purple (should just used colorbrewer...)
\definecolor{xrainbow0}{HTML}{e41a1c}
\definecolor{xrainbow1}{HTML}{a24057}
\definecolor{xrainbow2}{HTML}{606692}
\definecolor{xrainbow3}{HTML}{3a85a8}
\definecolor{xrainbow4}{HTML}{42977e}
\definecolor{xrainbow5}{HTML}{4aaa54}
\definecolor{xrainbow6}{HTML}{629363}
\definecolor{xrainbow7}{HTML}{7e6e85}
\definecolor{xrainbow8}{HTML}{9c509b}
\definecolor{xrainbow9}{HTML}{c4625d}
\definecolor{xrainbow10}{HTML}{eb751f}
\definecolor{xrainbow11}{HTML}{ff9709}

%------- %
% XHFILL %
%------- %


%%% Chapter Headings %%%

% \usepackage[explicit]{titlesec}
% \newcommand*\chapterlabel{}\pmod
% \titleformat{\chapter}
%   {\gdef\chapterlabel{}
%    \normalfont\sffamily\Huge\bfseries\scshape}
%   {\gdef\chapterlabel{\thechapter\ }}{0pt}
%   {\begin{tikzpicture}[remember picture,overlay]
%     \node[yshift=-3cm] at (current page.north west)
%       {\begin{tikzpicture}[remember picture, overlay]
%         \draw[fill=LightSkyBlue] (0,0) rectangle
%           (\paperwidth,3cm);
%         \node[anchor=east,xshift=.9\paperwidth,rectangle,
%               sharp corners=downhill=20pt,inner sep=11pt,
%               fill=MidnightBlue]
%               {\color{white}\chapterlabel#1};
%        \end{tikzpicture}
%       };
%    \end{tikzpicture}
%   }
% \titlespacing*{\chapter}{0pt}{50pt}{-60pt}

% \usepackage[Conny]{fncychap}
% \usepackage[Rejne]{fncychap}
% \ChNameVar{\bfseries}  % Makes the chapter "Chapter #" bold
% \ChNumVar{\bfseries}   % Makes the chapter number bold
% \ChTitleVar{\bfseries} % Makes the chapter title bold
% % Define a custom color, for example:
% \definecolor{mycolor}{RGB}{0,128,255}

% % Redefine the chapter style in Rejne to change the line colors
% \makeatletter
% \ChRuleWidth{2pt}   % Change the thickness of the lines
% \renewcommand{\DOCH}{%
%   \vspace*{-50\p@}% Moves the chapter title up/down if needed
%   {\color{mycolor} \hrule \@chapapp{} \space \thechapter \hrule}% Customizes the chapter header with color
% }
% \renewcommand{\DOTI}[1]{%
%   \vskip 20\p@ % Adjusts the space above the title
%   \bfseries #1\par % Embolden the title
%   \vskip 20\p@ % Adjusts the space below the title
% }
% \makeatother

% %% Change Chapter Heading Placement %%
% \usepackage{etoolbox}
% \makeatletter
% \patchcmd{\@makechapterhead}{\vspace*{50\p@}}{\vspace*{-20\p@}}{}{}
% \patchcmd{\@makeschapterhead}{\vspace*{50\p@}}{\vspace*{-20\p@}}{}{}
% \patchcmd{\DOTI}{\vskip 80\p@}{\vskip 40\p@}{}{}
% \patchcmd{\DOTIS}{\vskip 40\p@}{\vskip 0\p@}{}{}
% \makeatother

\renewcommand{\thesection}{\thechapter.\arabic{section}} %% Chapter.Section Numbering


%%% FIGURES %%%
\usepackage{graphicx}  
\graphicspath{ {images/} }  
% \numberwithin{figure}{section}
\usepackage{float}
\usepackage{caption}

%%% Hyperlinks %%%
\usepackage{hyperref}
\definecolor{horange}{HTML}{f58026}
\hypersetup{
	colorlinks=true,
	linkcolor=horange,
	filecolor=horange,      
	urlcolor=horange,
}


%% Headers and Footers %%
\usepackage{fancyhdr} % This should be set AFTER setting up the page geometry
\pagestyle{fancy} % options: empty , plain , fancy
\renewcommand{\chaptermark}[1]{\markboth{\thechapter.\ #1}{}}
\fancyhead[R]{\leftmark}
\fancyhead[L]{Hendrix College}
\fancyhead[C]{\includegraphics[height=.50cm]{images/small logo.png}}
\usepackage{xpatch}
\xpretocmd\headrule{\color{orangehdx}}{}{\PatchFailed}
\setlength{\footskip}{0.5in}
\setlength{\headheight}{18.5764pt}


%%%%% Colored Boxes %%%%%
\usepackage{tcolorbox}
\tcbuselibrary{skins}
\tcbuselibrary{theorems}
\tcbuselibrary{minted}
\newcounter{BoxCounter}


\newtcolorbox{note}{colback=black!15!white, colframe=black, boxrule=0.5mm, arc=4mm, left=4mm, right=4mm, top=2mm, bottom=2mm}

%% Numbered Theorems %%
\newtcolorbox[use counter=BoxCounter, number within=section]{theorem}[1][]{%
enhanced,
sharp corners=uphill,
    colframe=xblue, %color of frame
    colback=xblue!15, %color of background
    coltitle=white, %color of title text
    label={thm:\theBoxCounter},
    title={\large \textbf{Theorem \theBoxCounter}},
    attach boxed title to top left={xshift=0.5cm},
    boxed title style={colback=xblue},
    #1
}


%% Named Theorems %%
\newtcolorbox[use counter=BoxCounter, number within=section]{ntheorem}[2][]{%
enhanced,
sharp corners=south,
    colframe=xblue, %color of frame
    colback=xblue!15, %color of background
    coltitle=white, %color of title text
    label={thm:#2},
    title={\large \textbf{Theorem \theBoxCounter: #2}},
    attach boxed title to top center,
    boxed title style={colback=xblue},
    #1
}


%% Definitions %%
\newtcolorbox[use counter=BoxCounter, number within=section]{definition}[1][]{%
enhanced,
sharp corners=downhill,
    colframe=xred,
    colback=xred!15,
    coltitle=white,
    label={def:\theBoxCounter},
    title={\large \textbf{Definition \theBoxCounter}},
    attach boxed title to top right={xshift=-0.5cm},
    boxed title style={colback=xred},
    #1
}


%% Axioms %%
\newtcolorbox[]{axiom}[2][]{%
enhanced,
sharp corners=downhill,
    colframe=xdarkgreen,
    colback=xdarkgreen!15,
    coltitle=white,
    label={ax:#2},
    title={\large \textbf{Axiom of #2}},
    attach boxed title to top left={xshift=0.5cm},
    boxed title style={colback=xdarkgreen},
    #1
}


%% Exercises %%
\newtcolorbox[]{exercise}[2][]{%
    enhanced,
    sharp corners=uphill,
    colframe=orangehdx!75,
    colback=orangehdx!10,
    coltitle=white,
    label={ex:#2},
    title={\large \textbf{Exercise: #2}},
    attach boxed title to top left={xshift=0.5cm},
    boxed title style={colback=orangehdx!75},
    #1,
}


%% Examples %%
\newcounter{ExampleCounter}
\newtcolorbox[use counter=ExampleCounter, number within=chapter]{example}[2][]{%
enhanced,
sharp corners=south,
    colframe=teal,
    colback=teal!15,
    coltitle=white,
    label={ex:\theExampleCounter},
    title={\large \textbf{Example \theExampleCounter: #2}},
    attach boxed title to top center,
    boxed title style={colback=teal},
    #1,
}


%% Corollaries %%
\newtcolorbox[]{corollary}[3][]{%
enhanced,
sharp corners=downhill,
    colframe=xgrayblue,  
    colback=xgrayblue!15,        
    coltitle=white,
    title={\large \textbf{Corollary to #3: #2}},
    attach boxed title to top left={xshift=0.5cm},
    boxed title style={colback=xgrayblue},
    #1
}


%% Lemmas %%
\newtcolorbox[use counter=BoxCounter, number within=chapter]{lemma}[1][]{%
enhanced,
sharp corners=uphill,
    colframe=xpurple,
    colback=xpurple!15,
    coltitle=white,
    label={lm:\theBoxCounter},
    title={\large \textbf{Lemma \theBoxCounter}},
    attach boxed title to top left={xshift=0.5cm},
    boxed title style={colback=xpurple},
    #1,
}

%% Referencing Macros %%

% Exercises
\newcommand{\defex}[1]{\hyperref[ex:#1]{Exercise #1}}

% Definitions
\newcommand{\defword}[2]{\hyperref[def:#2]{#1}}
\newcommand{\defref}[1]{\hyperref[def:#1]{Definition #1}}

% Theorems
\newcommand{\numrefthm}[1]{\hyperref[thm:#1]{Theorem #1}}
\newcommand{\namrefthm}[1]{\hyperref[thm:#1]{#1}}



%-%-%-%-%-%-%-%-%-%-%-%-%-%-%-%-%-%-%-%-%-%-%-%-%-%-%-%-%-%-%-%-%-%-%-%-%-%-%

%% MATH PACKAGES, ENVIRONMENTS, COMMANDS %%
%-%-%-%-%-%-%-%-%-%-%-%-%-%-%-%-%-%-%-%-%-%-%-%-%-%-%-%-%-%-%-%-%-%-%-%-%-%-%
%You'll need your own packages, theorem types, and commands.

\usepackage{fix-cm}
\usepackage{tikz}
\usetikzlibrary{shapes,backgrounds,calc,patterns}
\usepackage{amsmath,amsthm} 
\usepackage{mathtools}
\usepackage{amssymb}
\usepackage[framemethod=tikz]{mdframed}
\usepackage{mathrsfs}
\usepackage{changepage}
\usepackage{multicol}
\usepackage{slashed}
\usepackage{enumerate}
\usepackage{booktabs}
\usepackage{enumitem}
\usepackage{kantlipsum}  %This package lets us generate random text for example purposes.
\usepackage{pgfplots}


%%% Custom Comands %%%
% Natural Numbers 
\newcommand{\N}{\ensuremath{\mathbb{N}}}

% Whole Numbers
\newcommand{\W}{\ensuremath{\mathbb{W}}}

% Integers
\newcommand{\Z}{\ensuremath{\mathbb{Z}}}

% Rational Numbers
\newcommand{\Q}{\ensuremath{\mathbb{Q}}}

% Real Numbers
\newcommand{\R}{\ensuremath{\mathbb{R}}}

% Complex Numbers
\newcommand{\C}{\ensuremath{\mathbb{C}}}

\newcommand{\I}{\ensuremath{\mathbb{I}}}

\newcommand{\pfs}{\noindent\makebox[\linewidth]{\rule{\textwidth}{0.4pt}}\vspace{0.5cm}}

\newcommand{\caution}{\marginnote{\includegraphics[width=2em]{caution.png}}}

\newmdenv[
  topline=false,
  bottomline=true,
  rightline=false,
  leftline=true,
  linewidth=1.5pt,
  linecolor=black, % default color, will be overridden in custom commands
  backgroundcolor=draculabg, % Needed for Dracula theme
  fontcolor=draculafg, % Needed for Dracula theme
  innertopmargin=0pt,
  innerbottommargin=5pt,
  innerrightmargin=10pt,
  innerleftmargin=10pt,
  leftmargin=0pt,
  rightmargin=0pt,
  skipabove=\topsep,
  skipbelow=\topsep,
]{customframedproof}

\newenvironment{proofpart}[2][black]{
    \begin{mdframed}[
        topline=false,
        bottomline=false,
        rightline=false,
        leftline=true,
        linewidth=1pt,
        linecolor=#1!40, % Custom color
        innertopmargin=10pt,
        innerbottommargin=10pt,
        innerleftmargin=10pt,
        innerrightmargin=10pt,
        leftmargin=0pt,
        rightmargin=0pt,
        skipabove=\topsep,
        skipbelow=\topsep%
    ]
    \noindent
    \begin{minipage}[t]{0.08\textwidth}%
        \textbf{#2}%
    \end{minipage}%
    \begin{minipage}[t]{0.90\textwidth}%
        \begin{adjustwidth}{0pt}{0pt}%
}{
    \end{adjustwidth}
    \end{minipage}
    \end{mdframed}
}

% Define a new environment 'proofscratch' for Proof and Scratch Paper
\newenvironment{proofscratch}[2]{%
    \begin{mdframed}[
        topline=false,
        bottomline=false,
        rightline=false,
        leftline=false,
        backgroundcolor=draculabg, % Background color for Dracula theme
        fontcolor=draculafg,       % Font color for Dracula theme
        innerleftmargin=-6pt,
        innertopmargin=0pt,
        leftmargin=0pt,
        rightmargin=0pt,
        skipabove=\topsep,
        skipbelow=\topsep,
    ]
    % % Begin TikZ picture for the vertical dotted line
    % \begin{tikzpicture}[overlay, remember picture]
    %     % Draw a vertical dotted line at approximately the center of the mdframed width
    %     \draw[dotted, thick] ([xshift=0.005\linewidth]current bounding box.north west) -- ([xshift=0.005\linewidth]current bounding box.south west);
    % \end{tikzpicture}%
    % % Create the left minipage for Proof
    \begin{minipage}[t]{0.47\textwidth}%
        \noindent\textit{Proof.} #1 \hfill $\qed$
    \end{minipage}%
    \hfill
    % Create the right minipage for Scratch Paper
    \begin{minipage}[t]{0.47\textwidth}%
        \textit{Scratch Paper.} #2
    \end{minipage}%
}{
    \end{mdframed}
}

\newcommand{\pfscratch}[3]{%
    \begin{customframedproof}[linecolor=#3]
        \begin{proofscratch}
            {#1}{#2}
        \end{proofscratch}
    \end{customframedproof}
}

\newcommand{\lepfscratch}[2]{\pfscratch{#1}{#2}{teal}}

\newcommand{\sbseteqpf}[5][We need to show these expressions are subsets of each other in order to prove they are equivalent.]{%
    \begin{customframedproof}[linecolor=#5]
        \begin{proof}
                #1
            \begin{proofpart}[#5]{(\(\subseteq\))}
                #2
            \end{proofpart}
    \vspace{2mm}
            \begin{proofpart}[#5]{(\(\supseteq\))}
                #3
            \end{proofpart}
    \vspace{2mm}
            #4
        \end{proof}
    \end{customframedproof}
}

% \newcommand{\iffpf}[4]{
%     \begin{customframedproof}[linecolor=#4]
%         \begin{proof}
%             We show this by proving both implications:
%             \begin{proofpart}[#4]{(\(\Rightarrow\))}
%                 #1
%             \end{proofpart}
%     \vspace{2mm}
%             \begin{proofpart}[#4]{(\(\Leftarrow\))}
%                 #2
%             \end{proofpart}
%     \vspace{2mm}
%             #3
%         \end{proof}
%     \end{customframedproof}
% }

\newcommand{\iffpf}[5][We show this by proving both implications:]{
    \begin{customframedproof}[linecolor=#5]
        \begin{proof}
            #1 % This is the optional text, with a default value
            \begin{proofpart}[#5]{(\(\Rightarrow\))}
                #2
            \end{proofpart}
    \vspace{2mm}
            \begin{proofpart}[#5]{(\(\Leftarrow\))}
                #3
            \end{proofpart}
    \vspace{2mm}
            #4
        \end{proof}
    \end{customframedproof}
}



% Command for lemma proofs with a predefined color
\newcommand{\expf}[1]{
    \begin{customframedproof}[linecolor=orangehdx!75,]
        \begin{proof}
        #1
        \end{proof}
    \end{customframedproof}
}

% Command for proofs that are contained within solution environments. 
\newcommand{\innerpf}[1]{%
    \begin{customframedproof}[linecolor=xgray]
        \begin{proof}
            #1
        \end{proof}
    \end{customframedproof}
}


% Command for theorems proofs with a predefined color
\newcommand{\tpf}[1]{
    \begin{customframedproof}[linecolor=xblue]
        \begin{proof}
        #1
        \end{proof}
    \end{customframedproof}
}
% Command for corollaries proofs with a predefined color
\newcommand{\cpf}[1]{
    \begin{customframedproof}[linecolor=purplehdx]
        \begin{proof}
        #1
        \end{proof}
    \end{customframedproof}
}

\newcommand{\sol}[1]{
    \begin{customframedproof}[linecolor=orangehdx!75,]
        \begin{solution}
        #1
        \end{solution}
    \end{customframedproof}
}

\newcommand{\lesol}[1]{
    \begin{customframedproof}[linecolor=teal]
        \begin{solution}
        #1
        \end{solution}
    \end{customframedproof}
}

\newcommand{\lepf}[1]{
    \begin{customframedproof}[linecolor=teal]
        \begin{proof}
        #1
        \end{proof}
    \end{customframedproof}
}

\newcommand{\pb}[2]{
    \begin{exercise}{#1}[label=ex:#1]
    #2
    \end{exercise}
}

\def \proofDistance {10pt}
\newenvironment{solution}
  {\textit{Solution.}}

% Command for theorems proofs with a predefined color
\newcommand{\epf}[1]{
    \begin{customframedproof}[linecolor=teal]
        \begin{proof}
        #1
        \end{proof}
    \end{customframedproof}
}

\newcommand{\disabs}[1]{\ensuremath{\left|#1\right|}}
\newcommand{\limn}[1]{\ensuremath{\lim_{n \rightarrow \infty}}#1}
\newcommand{\limx}[2]{\lim_{x \rightarrow #1}#2}

\newcommand{\limc}[1]{\lim_{x \rightarrow c}#1}

\newcommand{\ninn}{n \in \N}


\newcommand{\limsupn}[1]{\limsup_{n\rightarrow \infty}#1}
\newcommand{\liminfn}[1]{\liminf_{n\rightarrow \infty}#1}

\renewcommand{\theenumi}{\alph{enumi}} 
\renewcommand{\labelenumi}{(\theenumi)}


% end of preamble
%-%-%-%-%-%-%-%-%-%-%-%-%-%-%-%-%-%-%-%-%-%-%-%-%-%-%-%-%-%-%-%-%-%-%-%-%-%-%

\begin{document}

\numberwithin{BoxCounter}{section}

	
%-%-%-%-%-%-%-%-%-%-%-%-%-%-%-%-%-%-%-%-%-%-%-%-%-%-%-%-%-%-%-%-%-%-%-%-%-%-%
 %%% COVER PAGE %%%
 %-%-%-%-%-%-%-%-%-%-%-%-%-%-%-%-%-%-%-%-%-%-%-%-%-%-%-%-%-%-%-%-%-%-%-%-%-%-%

% Do not use all caps.
\newcommand{\titlestandin}[0]{Real Analysis}
\newcommand{\cussubtitle}[0]{MATH 350}
% Date format should be like: "January 1, 2001"
\newcommand{\startdate}[0]{August 26, 2024}
\newcommand{\customenddate}[0]{December 2, 2024}
% Be sure to include degree recognition (e.g., B.S., M.S., Ph.D)
\newcommand{\professor}[0]{Prof. Christopher Camfield, Ph.D.}

%-%-%-%-%-%-%-%-%-%-%-%-%-%-%-%-%-%-%-%-%-%-%-%-%-%-%-%-%



\begin{titlepage}
\begin{center}

\vspace*{-2cm}
\includegraphics[width=0.8\textwidth]{images/Hendrix Logo.png}\\
\vfill

% Horizontal line above the title in 'horange' color
\textcolor{horange}{\rule{\textwidth}{1.0pt}}

\vspace{2em}

{\huge \textbf{\titlestandin}}

\vspace{1em} % Space between the title and the bottom line

\textcolor{horange}{\rule{\textwidth}{1.0pt}}

\vspace*{1\baselineskip}

{\LARGE \textbf{\cussubtitle}}

\begin{large}
\vspace*{2\baselineskip}

\textit{Start}  \\[1ex]
{\scshape \startdate} \\[0.3\baselineskip] % Year published

\vspace*{1\baselineskip}

\emph{Author} \\[1ex]
%Submitted by \\[\baselineskip]
{\Large Paul Beggs \\ \par} % Editor list
{\href{mailto:BeggsPA@Hendrix.edu}{{BeggsPA@Hendrix.edu}}}\\ % Editor affiliation

\vspace*{1\baselineskip}

\textit{Instructor} \\[1ex] % Tagline(s) or further description
\professor

\vspace*{1\baselineskip}

\textit{End}\\[1ex]
{\scshape  \customenddate} \\[0.3\baselineskip] % Year published

\thispagestyle{empty} 

\end{large}
\end{center}
\end{titlepage}
\pagebreak


%-%-%-%-%-%-%-%-%-%-%-%-%-%-%-%-%-%-%-%-%-%-%-%-%-%-%-%-%-%-%-%-%-%-%-%-%-%-%

%%% Table of Contents %%%

% No need to edit.

% -%-%-%-%-%-%-%-%-%-%-%-%-%-%-%-%-%-%-%-%-%-%-%-%-%-%-%-%-%-%-%-%-%-%-%-%-%-%
	
	\thispagestyle{plain}
	% Table of Contents will automatically update based on included chapters
	% You may need to compile twice for the toc to update
	\begin{spacing}{1} %Can change spacing between entries in toc
		\renewcommand{\contentsname}{\Large\textbf{Table of Contents}} %Can rename here
		\tableofcontents
		\addtocontents{toc}{\protect\enlargethispage{\baselineskip}} 
	\end{spacing}


% -%-%-%-%-%-%-%-%-%-%-%-%-%-%-%-%-%-%-%-%-%-%-%-%-%-%-%-%-%-%-%-%-%-%-%-%-%-%

\chapter{The Real Numbers}\label{Real nums}
\vspace*{-0.25in}
\renewcommand{\theenumi}{\arabic{enumi}}
\renewcommand{\labelenumi}{\theenumi.}
\section{Types of Numbers}

\begin{definition}
    
The \textit{natural numbers} contain all positive, non-zero, and non-fractional numbers. Expressed as \(\N = \{1,2,3,4\dots\}\). They do not have an additive inverse, but you can add and multiply them.
\end{definition}

\begin{definition}
    
The \textit{integers} contains all non-fractional numbers. Expressed as: \(\Z = \{\dots-2,-1,0,1,2\dots\}\)---are known as a Group (more specifically, a ``ring''). You can add, multiply, and subtract these numbers.
\end{definition}

\begin{definition}
    
The \textit{rational numbers} contain all numbers, except irrational numbers. Expressed as: \(\Q = \{\frac{p}{q} \ | \ p,q \in \Z, q \ne 0\}\)---are known as a ``Field.'' You can add, subtract, multiply, and divide these numbers.

\end{definition}

A problem that rational numbers could not explain: The \(45, 45, 90\) triangle had a hypotenuse of \(\sqrt{2}\). This did not exist at the time, so it was simply \(c^2 = 2\). Therefore, new numbers needed to be invented.

\begin{theorem}
    There does not exist a rational number \(r\) such that \(r^2 = 2\).
\end{theorem}

\tpf{
    Suppose there exists a rational number \(r\) such that \(r^2 = 2\). Since \(r\) is rational, there exists \(p,q \in \Z\) such that \(r = \frac{p}{q}\). We can assume the \(p\) and \(q\) have no common factors. (If not, we can factor out the common factor.) By our assumption,
    \begin{align*}
        r^2             & = 2 \\
        \frac{p^2}{q^2} & = 2
    \end{align*}
    It follows that, \begin{align*}
        p^2 = 2q^2
    \end{align*}
    Such that \(p^2\) is an even number because if \(p\) were odd, then \(p^2\) would be odd. There exists \(x \in \Z\) such that \(p = 2x\). Recall that \(p^2 = 2q^2\). Thus\begin{align*}
        (2x)^2 & = 2q^2 \\
        4x^2   & = 2q^2 \\
        2x^2   & = q
    \end{align*} Thus, \(q^2\) is even. Hence \(q\) is also even. So \(p\) and \(q\) are both divisible by 2. This contradicts that \(p\) and \(q\) have no common factors. Thus, our supposition is false. Therefore, there does not exist a rational number \(r\) such that \(r^2 = 2\)
}

So we are going to work with a larger set called the real numbers, \(\R\).
\begin{itemize}
    \item \(\N \subset \Z \subset \Q \subset \R\)
    \item You can:
          \begin{itemize}
              \item Add,
              \item Subtract,
              \item Multiply,
              \item Divide
          \end{itemize}
    \item In other words, all field axioms apply.
    \item Totally ordered set for any \(x,y \in \R\). Thus, one of these are true:
          \begin{enumerate}
              \item \(x < y\),
              \item \(x > y\),
              \item \(x = y\)
          \end{enumerate}
    \item  Think of it as a number line.
    \item \(\Q\) is dense: \\
          If \(a,b \in \Q\) with \(a \ne b\), there exists \(c \in \Q\) which is between \(a\) and \(b\) such that \(a < c < b\). One example is \(\frac{a + b}{2}\).
    \item \(\Q\) is not \textit{complete}, but \(\R\) is.
          \begin{itemize}
              \item \textit{Complete:} Think, ``no gaps.''
          \end{itemize}
\end{itemize}

\renewcommand{\theenumi}{\arabic{enumi}}
\renewcommand{\labelenumi}{\theenumi.}
\section{Preliminaries}

Things to remember from Intro and Discrete.

\begin{table}[htbp]
    \begin{tabular}{c||c}
        Set Notation & Complement                     \\
        \hline
        \(x \in A\)  & \(A^c\) (not \(\overline{A}\)) \\
        \(A \cup B\) & \(\R \setminus A\)             \\
        \(A \cap B\)
    \end{tabular}
\end{table}

\begin{itemize}
    \item \(\displaystyle \bigcup^\infty_{n=1} A_n = A_1 \cup A_2 \cup A_3 \cup \dots\). \\

    \item \(\displaystyle \bigcap^\infty_{n=1} A_n = A_1 \cap A_2 \cap \dots\). \\
\end{itemize}

\begin{definition}
    \textit{De Morgan's Laws} are defined as \((A \cup B)^c = A^c \cap B^c\) and \((A \cap B)^c = A^c \cup B^c\). 
\end{definition}


\subsection{Infinite Unions and Intersections} \hfill

For each \(n \in \N\), define \(A_n = \{n, n+1, n+2,\dots\} = \{k \in \N \ | \ k \geq n\}\). In other words, each subsequent element in the subset will start at \(n\). For example, \(A_1 = \{1,2,\dots\}\), whereas \(A_5 = \{5,6,\dots\}.\)   \\

\(\bigcup^\infty_{n=1} A_n = \N\). To show a number \(\in \N\) belongs in the set \(A_n\), we can start with that, \(k \in \N\). Then \(k \in A_k\). Thus, \(k \in A_k \subseteq \bigcup^\infty_{n=1}A_n\). Therefore, \(\N \subseteq \bigcup^\infty_{n=1}A_n\). \\

\(\bigcap^\infty_{n=1}A_n = \emptyset\). Obviously, we know that the empty set is a subset of \(A_n\), but to prove that \(\bigcap^\infty_{n=1}A_n\) is a subset of the empty set, we should suppose a \(k \in \N\) such that \(k \in \bigcap^\infty_{n=1}A_n\). Notice that \(k \notin \bigcap^\infty_{n=1}A_n\). So, \(\bigcap^\infty_{n=1}A_n = \emptyset\).

\subsection{Functions and Notation} \hfill

\(f \colon A \rightarrow B\) where \(f\) is a function, \(A\) is a domain, and \(B\) is the co-domain. Thus, \(f(x) = y\) such that \(x \in A\) and \(y \in B\). \\

\subsubsection{Some definitions to keep in mind}


\begin{definition}
    The \textit{Dirichlet Function} is defined as \[f(x) = \begin{cases} 1 & \text{if } x \in \Q \\ 0 & \text{if } x \notin \Q \end{cases}\]
\end{definition}

\begin{definition}
    Let \(f \ \colon \ \R \rightarrow \R\). If \(E \subseteq \R\), then \(f(E) = \{f(x) \ | \ x \in E\}\).
\end{definition}

Example: \(g \ \colon \ \R \rightarrow \R\), when we say \(y \in g(A)\) implies there exists an \(x\) such that \(g(x) = y\) \\

\hypertarget{Triangle Inequality}{}
\begin{definition}
    The \textit{Triangle Inequality} is defined as: For any \(a,b \in \R\), \(|a + b| \leq  |a| + |b|\).
\end{definition}
The most common application: For any \(a,b,c \in \R\), \(|a - b| \leq |a-c| + |c-b|\), with the intermediate step of \(a - b = (a-c) + (c-b)\).

\begin{definition}
    A function \(f\) is \textit{injective} (or \textit{one-to-one}) if \(a_1 \ne a_2\) in \(A\) implies that \(f(a_1) \ne f(a_2)\) in \(B\). Note the contrapositive of this definition: If \(f(a_1) = f(a_2)\), then \(a_1 = a_2\).
\end{definition}

\begin{definition}
    A function \(f\) is \textit{surjective} (or \textit{onto}) if for every \(b \in B\), there exists an \(a \in A\) such that \(f(a) = b\). Note the contrapositive of this definition: If there exists a \(b \in B\) such that there is no \(a \in A\) such that \(f(a) = b\), then the function is not surjective.
\end{definition}

\subsection{Common Strategies for Analysis Proofs}

\setcounter{BoxCounter}{5}
\begin{theorem}
    Let \(a,b \in \R\). Then, \[a = b \text{ if and only if for all } \epsilon > 0, |a-b| < \epsilon.\]
\end{theorem}

\iffpf
{Assume \(a = b\). Let \(\epsilon > 0\). Then \(|a-b| = 0 < \epsilon\)}
{Assume for all \(\epsilon > 0\), \(|a-b| < \epsilon\). Suppose \(a \ne b\). Then \(a - b \ne 0\). So, \(|a-b| \ne 0\). Now, Consider \(\epsilon_0 = |a-b|\). By our assumption we know that \(|a-b| < \epsilon_0\). It is not true that \(|a-b| < |a - b|\). Therefore, it must be the case that \(a = b\).}
{Therefore, by showing both sides of the implication accomplish the same thing as the other side, we know that \(a = b \text{ if and only if for all } \epsilon > 0, |a-b| < \epsilon.\)}
{xblue}

\subsection{Mathematical Induction} \hfill

\textit{Inductive Hypothesis:} Let \(x_1 = 1\). For all \(n \in \N\), let \(x_{n+1} = \frac{1}{2}x_n + 1\). \\

\textit{Inductive Step:} \(x_1 = 1, x_2 = 1.5, x_3 = 1.75, x_4 = 1.875\). \\

\begin{example}
    {Induction} The sequence \((x_n)\) is increasing. In other words, for all \(n\in \N\), \(x_n \leq x_{n+1}\).
\end{example}

\epf{
Suppose the sequence \((x_n)\) is increasing. We will prove this point by using induction. \\

\textbf{Base Case:} We see that \(x_1 = 1\) and \(x_2 = 1.5\). Thus, \(x_1 \leq x_2\). \\

\textbf{Inductive Hypothesis:} For \(n \in \N\), assume \(x_n \leq x_{n+1}\). \\

\indent \textit{Scratch work:} We want: \(x_{n+1} \leq x_{n+2}\). We know: \(x_{n+1} = \frac{1}{2}x_{n+1}+1\). \\

\textbf{Inductive Step:} Then \(\frac{1}{2}x_n \leq \frac{1}{2}x_{n+1}\). Hence, \(\frac{1}{2}x_{n} + 1 \leq \frac{1}{2}x_{n+1} + 1\). Therefore we have proven through induction that, \(x_{n+1} \leq x_{n+2}\).
}

\renewcommand{\theenumi}{\alph{enumi}}
\renewcommand{\labelenumi}{(\theenumi)}
\subsection{Exercises}

\begin{exercise}
    {1.2.3}Decide which of the following represent true statements about
    the nature of sets. For any that are false, provide a specific example where the
    statement in question does not hold.
    \begin{enumerate}
        \item If \(A_1 \supseteq A_2 \supseteq A_3 \supseteq A_4\dots\) are all sets containing an infinite number of
              elements, then the intersection \(\bigcap^\infty_{n=1} A_n\) is infinite as well.
        \item If \(A_1 \supseteq A_2 \supseteq A_3 \supseteq A_4\dots\) are all finite, nonempty sets of real numbers,
              then the intersection \(\bigcap^\infty_{n=1} A_n\) is finite and nonempty.
        \item \(A \cap (B \cup C) = (A \cap B) \cup C\)
        \item \(A \cap (B \cap C) = (A \cap B) \cap C\)
        \item \(A \cap (B \cup C) = (A \cap B) \cup (A \cap C)\)
    \end{enumerate}
\end{exercise}

\sol{\hfill
    \begin{enumerate}
        \item This is false. Consider the following as a counterexample: If we define \(A_1\) as \(A_n = \{n, n+1, n+2,\dots\} = \{k \in \N \ | \ k \geq n\}\), we can see why the intersection of these sets of infinite numbers are actually empty. Consider a number \(m\) that actually satisfies \(m \in A_n\) for every \(A_n\) in our collection of sets. Because \(m\) is not an element of \(A_{m + 1}\), no such \(m\) exists and the intersection is empty.
        \item This is true.
        \item False. Consider sets \(A = \{1, 2, 3\}\), \(B = \{3, 6, 7\}\) and \(C = \{5\}\). Note that
              \(A \cap (B \cup C) = \{3\}\) is not equal to \((A \cap B) \cup C = \{3, 5\}\).
        \item  This is true. A proof would start with \(x \in A \cap (B \cap C)\).
        \item This is true. A proof would start with \(x \in A \cap (B \cup C)\).
    \end{enumerate}
}

\newpage
\begin{exercise}
    {1.2.5}\hyperref[def:1.2.1]{De Morgan's Laws} Let \(A\) and \(B\) be subsets of \(\R\).
    \begin{enumerate}
        \item  If \(x \in (A \cap B)^c\), explain why \(x \in A^c \cup B^c\). This shows that \((A \cap B)^c \subseteq A^c \cup B^c\).
        \item Prove the reverse inclusion \((A \cap B)^c \supseteq A^c \cup B^c\), and conclude that \((A \cap B)^c = A^c \cup B^c\).
        \item Show \((A \cup B)^c = A^c \cap B^c\) by demonstrating inclusion both ways.
    \end{enumerate}
\end{exercise}

\sol{\hfill
    \begin{enumerate}
        \item If \(x \in (A \cap B)^c\), and we know that \(A^c = \{x \in \R \ \colon x \notin A\}\), then we know \(x\) must cannot exist in \(A^c\) and \(B^c\) because \((A \cap B)^c = \{x \in \R \ \colon x \notin (A \cap B)\}\). Thus, \(x\) is in either \(A^c\) or \(B^c\). Put another way \(x \in A^c \cup B^c\). Since we have shown that an element that started in \((A \cap B)^c\) ended up in \(A^c \cup B^c\), then we know \((A \cap B)^c \subseteq A^c \cup B^c\).
        \item Assume there exists a \(y \in A^c \cup B^c\). Thus, it must be the case that \(y \notin A\) or \(y \notin B\). Hence, \(y\) cannot be exist in both sets at the same time, so \(y \in (A \cap B)^c\). Because we have taken an element that started in \(A^c \cup B^c\) and have shown that it exists in \((A \cap B)^c\), we have proven \(A^c \cup B^c \subseteq (A \cap B)^c\).
        \item\hfill \sbseteqpf
              {Assume there exists \(x \in (A \cap B)^c\), and we know that \(A^c = \{x \in \R \ \colon x \notin A\}\), then we know \(x\) must cannot exist in \(A^c\) and \(B^c\) because \((A \cap B)^c = \{x \in \R \ \colon x \notin (A \cap B)\}\). Thus, \(x\) is in either \(A^c\) or \(B^c\). Put another way \(x \in A^c \cup B^c\). Since we have shown that an element that started in \((A \cap B)^c\) ended up in \(A^c \cup B^c\), then we know \((A \cap B)^c \subseteq A^c \cup B^c\).}
              {Now assume there exists a \(y \in A^c \cup B^c\). Thus, it must be the case that \(y \notin A\) or \(y \notin B\). Hence, \(y\) cannot be exist in both sets at the same time, so \(y \in (A \cap B)^c\). Because we have taken an element that started in \(A^c \cup B^c\) and have shown that it exists in \((A \cap B)^c\), we have proven \(A^c \cup B^c \subseteq (A \cap B)^c\).}
              {Therefore, we have shown through proving both sides of the implication, that these two statements are logically equivalent. In that, all elements of \(A^c \cup B^c\) are the same elements that are in \((A \cap B)^c\)}
              {xgray}
    \end{enumerate}
}



\begin{exercise}
    {1.2.7}Given a function \(f\) and a subset \(A\) of its domain, let \(f(A)\)
    represent the range of \(f\) over the set \(A\); that is, \(f(a) = \{f(x) \colon x \in A\}\).
    \begin{enumerate}
        \item Let \(f(x) = x^2\). If \(A = [0,2]\) (the closed interval \(\{x \in \R \colon 0 \leq x \leq 2\}\)) and \(B = [1,4]\), find \(f(A)\) and \(f(B)\). Does \(f(A\cap B) = f(A) \cap f(B)\) in this case? Does \(f(A \cup B) = f(A) \cup f(B)\)?
        \item Find two sets \(A\) and \(B\) for which \(f(A\cap B) \ne f(A) \cap f(B)\).
        \item Show that, for an arbitrary function \(g \ \colon \ \R \rightarrow \R\), it is always true that \(g(A \cap B) \subseteq g(A) \cap g(B)\) for all sets \(A,B\subseteq \R\).
        \item Form and prove a conjecture about the relationship between \(g(A \cup B)\) and \(g(A) \cup g(B)\) for an arbitrary function \(g\).
    \end{enumerate}
\end{exercise}

\sol{\hfill
    \begin{enumerate}
        \item Since \(f(x) = x^2\), the intervals of \(f(A)\) would be \([0,4]\) and \(f(B)\) would be \([1,16]\). The interval of the intersection of \(A \cap B\) is \([1,2]\). Take this through our function, we get \(f(A \cap B) = [1,4]\). On the other side of the equation, we already know the intervals of \(f(A)\) and \(f(B)\), and the intersection of theirs would be \([1,4]\). So they do equal each other. We know \(f(A \cup B)\) and \(f(A) \cup f(B)\) will be equivalent because \(f(A \cup B)\) has an interval of \([0,16]\), and \(f(A) \cup f(B)\) also has an interval of \([0,16]\) because taking the union of \([0,4] \cup [1,16]\) is \([0,16]\).
        \item Two sets could be \(A = [5,6]\) and \(B = [0,0]\). Because the sets have nothing in common even after taking their function, they do not equal each other.
        \item \hfill \innerpf{Let \(x \in g(A \cap B)\). Using the definition of function, we know there exists a \(y \in A \cap B\) to which that \(y\) is mapped to as \(g(y) = x\). From the definition of intersection, we know \(y \in A\) and \(y \in B\) such that \(x = g(y) \in g(A)\) and \(x = g(y) \in g(B)\) because \(y \in A \cap B\). Putting it together, we have \(x \in g(A) \cap g(B)\) thus proving \(g(A \cap B) \subseteq g(A) \cap g(B)\)}
        \item Conjecture: For any function \(g\) defined as \(g \ \colon \R \rightarrow \R\) and for any subsets \(A,B \subseteq \R\), the following holds: \[g(A \cup B) = g(A) \cup g(B)\]
              \sbseteqpf
              {Take any element \(x \in g(A \cap B)\). By definition of function, we know there exists some \(y \in A \cup B\) such that \(g(x) = y\). From the definition of union, we know \(y \in A\) or \(y \in B\) such that \(x = g(y) \in g(A)\) or \(x = g(y) \in g(B)\) or both. Putting it together, we have \(x \in g(A) \cup g(B)\) thus proving \(g(A \cup B) \subseteq g(A) \cup g(B)\).}
              {Take any element \(p \in g(A) \cap g(B)\). By definition of union, we know \(p\) is either in \(g(A)\) or \(g(B)\) or both. From the definition of function, we know that if \(p \in g(A)\) or \(p \in g(B)\) then there exists some \(q \in A\) or \(q \in B\) such that \(g(q) = p\). Putting it together, we have \(q \in A \cup B\). Moreover, this means \(p = g(x) \in g(A \cup B)\). And since \(p \in g(A) \cup g(B)\) implies \(p \in g(A \cup B)\), we know \(g(A) \cup g(B) \subseteq g(A \cup B)\).}
              {Therefore, since we have proven that both expressions are functions of each other, we have proved that they are equal.}
              {xgray}
    \end{enumerate}
}

\begin{exercise}
    {1.2.8}Given a function \(f \ \colon A \rightarrow B\) can be defined as either \hyperref[def:1.2.5]{injective} or \hyperref[def:1.2.6]{surjective}, give an example of each or state that the request is impossible:
    \begin{enumerate}
        \item \(f \ \colon \N \rightarrow \N\) that is 1-1 but not onto.
        \item \(f \ \colon \N \rightarrow \N\) that is onto but not 1-1.
        \item \(f \ \colon \N \rightarrow \Z\) that is 1-1 and onto.
    \end{enumerate}
\end{exercise}

\sol{\hfill
    \begin{enumerate}
        \item The function \(f(a) + 1\) is 1-1 because when \begin{align*}
                  f(a_1)  & = f(a_2)  \\
                  a_1 + 1 & = a_2 + 1 \\
                  a_1     & = a_2
              \end{align*}
              However, the function is not onto because the entire co-domain is not covered. That being 1.
        \item \textit{We need to find a function that will cover every entry in the co-domain, while also avoiding a scenario where \(a_1 = a_2\)...} Consider the function,
              \[f(a) =
                  \begin{cases}
                      a     & \text{ if \textit{a} is odd,} \\
                      a - 1 & \text{ if \textit{a} is even}
                  \end{cases}\]
              This function is onto because every natural number is covered, but it is not 1-1 because \(a_1 \ne a_2 - 1\).
        \item This request is not possible. There is no way to map every natural number to every integer because we are simply missing 0! (Not 0 factorial, we do have the number 1, I just mean the number 0 in a exclamatory sense.)
    \end{enumerate}
}


\begin{exercise}
    {1.2.13}For this exercise, assume \hyperref[ex:1.2.5]{Exercise 1.2.5} has been successfully completed.
    \begin{enumerate}[label=(\alph*)]
        \item Show how induction can be used to conclude that
              \[
                  (A_1 \cup A_2 \cup \cdots \cup A_n)^c = A_1^c \cap A_2^c \cap \cdots \cap A_n^c
              \]
              for any finite \(n \in \mathbb{N}\).

        \item It is tempting to appeal to induction to conclude
              \[
                  \left( \bigcup_{i=1}^{\infty} A_i \right)^c = \bigcap_{i=1}^{\infty} A_i^c,
              \]
              but induction does not apply here. Induction is used to prove that a particular statement holds for every value of \(n \in \mathbb{N}\), but this does not imply the validity of the infinite case. To illustrate this point, find an example of a collection of sets \(B_1, B_2, B_3, \dots\) where
              \[
                  \bigcap_{i=1}^{n} B_i \neq \emptyset \quad \text{is true for every } n \in \mathbb{N},
              \]
              but
              \[
                  \bigcap_{i=1}^{\infty} B_i = \emptyset
              \]
              fails.

        \item Nevertheless, the infinite version of De Morgan’s Law stated in (b) is a valid statement. Provide a proof that does not use induction.
    \end{enumerate}
\end{exercise}
\newpage
\sbseteqpf%
%
[In this proof, we plan to prove (c). Thus, we need to show that:
    \[
        \left( \bigcup_{i=1}^{\infty} A_i \right)^c \subseteq \bigcap_{i=1}^{\infty} A_i^c
    \]
    and
    \[
        \left( \bigcup_{i=1}^{\infty} A_i \right)^c \supseteq \bigcap_{i=1}^{\infty} A_i^c.
    \]]%
%
{Let \(x \in \left( \bigcup_{i=1}^{\infty} A_i \right)^c\). This means \(x\) is in the union set of \(A_i\) for all \(i \in \N\). Then, because we are taking the complement of \(\left( \bigcup_{i=1}^{\infty} A_i \right)\), that means \(x \notin A_i\) for all \(i \in \mathbb{N}\). Hence, \(x\) is in the complement of each \(A_i\). Thus, we can use the definition of intersection to assert \(x \in \bigcap_{i=1}^{\infty} A_i^c\).

    Therefore, we have shown:
    \[
        \left( \bigcup_{i=1}^{\infty} A_i \right)^c \subseteq \ \bigcap_{i=1}^{\infty} A_i^c.
    \]}%
%
%
%
{Similar to before, let \(x \in \bigcap_{i=1}^{\infty} A_i^c\). Because \(x \in A^c_i\) for all \(i \in \N\) we know \(x \notin A_i\). Hence, \(x \notin \left(\bigcup_{i=1}^{\infty} A_i\right)\), which means \(x \in \left( \bigcup_{i=1}^{\infty} A_i \right)^c\).

    Therefore, we have shown:
    \[
        \left( \bigcup_{i=1}^{\infty} A_i \right)^c \supseteq \ \bigcap_{i=1}^{\infty} A_i^c.
    \]}%
%
%
{By showing both inclusions, we see that:
    \[
        \left( \bigcup_{i=1}^{\infty} A_i \right)^c = \bigcap_{i=1}^{\infty} A_i^c.
    \] \qedhere
}%
%
%
{xgray}




\setcounter{section}{2}
\renewcommand{\theenumi}{\arabic{enumi}}
\renewcommand{\labelenumi}{\theenumi.}
\section{Axiom of Completeness}

\begin{axiom}
    {Completeness}Every nonempty set of real numbers that is \hyperref[def:1.3.1]{bounded} has a \hyperref[def:1.3.2]{least upper bound}.
\end{axiom}

Think about \(\Q\) and \(\R\).
\begin{itemize}
    \item Both are fields.
          \begin{itemize}
              \item Both have \(+,-,\times,\div\) operations.
          \end{itemize}
    \item Both are totally ordered
          \begin{itemize}
              \item \(a < b\),
              \item \(a > b\),
              \item or \(a = b\)
          \end{itemize}
    \item \(\R\) is complete. \(\Q\) is not.
\end{itemize}

\subsection{Least Upper Bounds and Greatest Lower Bounds}

\begin{definition}
    A set \(A \subseteq \R\) is \textit{bounded above} if there exists a number \(b \in \R\) such that \(a \leq b\) for all \(a \in A\). The number \(b\) is called an \textit{upper bound} of \(A\).

    Similarly, a set \(A \subseteq \R\) is \textit{bounded below} if there exists a \textit{lower bound} \(l \in \R\) satisfying \(l \leq a\) for every \(a \in A\).
\end{definition}

\textbf{Note that upper bounds are not unique!} For example, consider the line, \(A\), from \(0\) to \(1\). There are infinitely many upper bounds past \(1\) because \(A\) is bounded.

\begin{definition}
    A number \(s\) is a \textit{least upper bound} for a set \(A \subseteq \R\) if it meets the following two criteria:
    \begin{enumerate}[label=(\roman*)]
        \item \(s\) is an upper bound for \(A\);
        \item if \(b\) is any upper bound for \(A\), then \(s \leq b\).
    \end{enumerate}
\end{definition}

We often call the least upper bound the \textit{supremum} of a set. 

\begin{example}
    {Supremum} Imagine a number line from \((1,8)\). Note that parenthesis mean \(<\) and not \(\leq\). Hence, the supremum is \(8\). Wrote simply as \(\sup A\).
\end{example}




\begin{example}
    {Supremum and Infimum 1}Consider a set, \(B = [-5,-2] \cup (3,6) \cup \{13\}\). What is the supremum and the infimum?
\end{example}

\lesol{
    \(\sup B = 13\); \(\inf B = -5\) because \(-5\) is the greatest lower bound.
}

\begin{example}
    {Supremum and Infimum 2}Consider the set, \(\C = \{\frac{1}{n} \ \colon \ n \in \N\}\). What is the supremum and the infimum?
\end{example}

\lesol{
    \(\sup \C = 1\), \(\inf \C = 0\).
}

\begin{example}
    Let \(A \subseteq \R\) be nonempty and bounded above, and let \(c \in \R\). Define the set \(c + A\) by \[c + A = \{c+a \ \colon \ a \in A\}\] Then \(\sup(c + A) = c + \sup A\).
\end{example}

\lesol{
    To properly verify this we focus separately on each part of \hyperref[def:1.3.2]{Definition 1.3.2}. Setting \(s = \sup A\), we see that \(a \leq s\) for all \(a \in A\), which implies \(c + a \leq c + s\) for all \(a \in A\). Thus, \(c + s\) is an upper bound for \(c + A\) and condition (i) is verified. For (ii), let \(b\) be an arbitrary upper bound for \(c + A\); i.e., \(c + a \leq b\) for all \(a \in A\). This is equivalent to \(a \leq b - c\) for all \(a \in A\), from which we conclude that \(b - c\) is an upper bound for \(A\). Because \(s\) is the least upper bound of \(A\), \(s \leq b - c\), which can be rewritten as \(c + s \leq b\). This verifies part (ii) of \hyperref[def:1.3.2]{Definition 1.3.2}, and we conclude \(\sup (c + A) = c + \sup A\).
}

\setcounter{BoxCounter}{3}

\begin{definition}
    A real number \(a_0\) is a \textit{maximum} of the set \(A\) if \(a_0\) is an element of \(A\) and \(a_a \geq a\) for all \(a \in A\). Similarly, a number \(a_1\) is a \textit{minimum} of \(A\) if \(a_1\) is an element of \(A\) and \(a_1 \leq a\) for all \(a \in A\).
\end{definition}

Note that some sets have a maximum and some sets do not. You cannot refer to a maximum without first knowing it exists. This is the same with minimums.

\begin{lemma}
    Assume \(s\) is an \hyperref[def:1.3.1]{upper bound} for a set \(A \subseteq \R\). Then, \(s\) is the supremum of \(A\) if and only if for every \(\epsilon > 0\), there exists \(x\in A\) such that \(s - \epsilon < x\).
\end{lemma}

This lemma allows us to take any positive number and take a ``step back.'' In essence, you can verify something as an upper bound if you continuously back up over and over until you cannot back up any longer.

\iffpf
{
    Assume \(s = \sup A\). Let \(\epsilon > 0\). Suppose there are no elements \(x\) of \(A\) such that \(s - \epsilon < x\). Then \(s - \epsilon\) would be an upper bound. This contradicts that \(s\) is the least upper bound. Therefore, there must exist an element \(x \in A\) such that \(s - \epsilon < x\).
}{
    Assume for every \(\epsilon > 0\), there exists \(x \in A\) such that \(s - \epsilon < x\). Let \(t\) be an upper bound of \(A\). Suppose \(t < s\). Consider \(\epsilon_0 = s - t > 0\). By our assumption, there exists \(x \in A\) such that \(s - \epsilon_0 < x\). So, \(t < x\). This contradicts that \(t\) is an upper bound of \(A\). So, \(t \geq s\). Thus, \(s\) is the least upper bound
}{
    Therefore, by proving both the right and left implication, we have shown the statement to be true.
}{xpurple}
Analogous statement about infimums: Assume \(z\) is a lower bound of a set \(A \subseteq \R\). Then \(z = \inf A \iff \) for all \( \epsilon >0\), there exists \(y \in A\) such that \(y < z + \epsilon\).

\renewcommand{\theenumi}{\alph{enumi}}
\renewcommand{\labelenumi}{(\theenumi)}
\subsection{Exercises}

\begin{exercise}
    {1.3.4}Let \(A_1,A_2,A_3\dots\) be a collection of nonempty sets each of which is bounded above.
    \begin{enumerate}
        \item Find a formula for \(\sup(A_1 \cup A_2)\). Extend this to \(\sup (\bigcup^n_{k=1}A_k)\).
        \item Consider \(\sup(\bigcup^\infty_{k=1} A_k)\). Does the formula in (a) extend to the infinite case?
    \end{enumerate}
\end{exercise}

\sol{\hfill
    \begin{enumerate}
        \item Let \(A_1\) and \(A_2\) be nonempty sets, each bounded above. To find the largest of the two suprema, we can use the following: \(\sup(A_1 \cap A_2) = \max\{\sup A_1, \sup A_2\}\). If we extend this notion to \(\sup (\bigcup^n_{k=1}A_k)\), we can use the same idea from before and write it as \(\sup (\bigcup^n_{k=1}A_k) = \max\{\sup A_1, \sup A_2, \dots, \sup A_n\}\).
        \item The formula does not extend to the infinite case. Consider the counterexample \(\bigcup^\infty_{k=1} A_k\) where \(A_k := [k, k+1]\). Even though these sets are bounded above, when we take the union of them, we approach infinity, which is not bounded: \(\bigcup^\infty_{k=1} A_k = [1,2] \cup [2,3] \cup \dots = [1, \infty)\).
    \end{enumerate}
}

\begin{exercise}
    {1.3.5}As in Example 1.3.7, let \(A \subseteq \R\) be nonempty and bounded above, and let \(c \in \R\). This time define the set \(cA = \{ca \ \colon a \in A\}\).
    \begin{enumerate}
        \item If \(c \geq 0\), show that \(\sup(cA) = c\sup A\).
        \item Postulate a similar type of statement for \(\sup(cA)\) for the case \(c < 0\).
    \end{enumerate}
\end{exercise}

\sol{
    \begin{enumerate}
        \item Let \(A \subseteq \R\) be nonempty and bounded above. Define the set \(cA := \{ca \ \colon a \in A\}\). From the axiom of completeness, because \(A\) is bounded above, we know there is a least upper bound, \(s = \sup A\). Following from Example 1.3.7, we see that \(a \leq s\) for all \(a \in A\) which implies \(ca \leq cs\) for all \(a \in A\). Thus, \(cs\) is an upper bound for \(cA\), and the first condition of \hyperref[def:1.3.2]{Definition 1.3.2} is satisfied. For the second condition, we need to look at both \(c = 0\) and \(c > 0\) to avoid dividing by zero. So, we have two cases:
              \begin{itemize}
                  \item \textbf{\(c = 0\):} If \(c = 0\), then \(cA = \{0 \colon \ a \in A\} = \{0\}\). Since the only element in \(cA\) is \(0\), \(\sup(cA) = 0\). Similarly, because \(c = 0\), \(c \sup A = 0 \cdot \sup A = 0\). Therefore, \(\sup(cA) = c \sup(A)\).
                  \item \textbf{\(c > 0\):} Let \(b\) be an arbitrary upper bound for \(cA\) and \(c > 0\). In other words, \(ca \leq b\) for all \(a \in A\). This is equivalent to \(a \leq b / c\) where \(c \ne 0\), from which we can see that \(b / c\) is an upper bound for \(A\). Because \(s\) is the least upper bound of \(A\), \(s \leq b / c\), which can be rewritten as \(cs \leq b\). This verifies the second part of \hyperref[def:1.3.2]{Definition 1.3.2}, and we conclude \(\sup(cA) = c \sup A\).
              \end{itemize}
        \item Postulate: If \(c < 0\), then \(\sup(cA) = c \inf (A)\). \\
              % \expf{Assume \(c < 0\), Multiplying \(a \leq s\) by \(c\) reverses the inequality, so \(ca \geq cs\) for all \(a \in A\). Therefore, \(cs\) is a lower bound of for \(cA\). To verify this as the lowest upper bound, let \(m  = \inf A\). Using the same logic as before, for all \(a \in A\), \(m \leq a\), so \(cm \geq ca\). This shows that \(cm\) is the least upper bound of \(cA\), and thus \(\sup(cA) = cm = c \inf A\).}
    \end{enumerate}
}

\begin{exercise}
    {1.3.8}Compute, without proofs, the suprema and infima (if they exist) of the following sets:
    \begin{enumerate}
        \item \(\left\{\frac{m}{n} : m, n \in \mathbb{N} \text{ with } m < n\right\}\).
        \item \(\left\{\frac{{(-1)}^{m}}{n} : m, n \in \mathbb{N}\right\}\).
        \item \(\left\{\frac{n}{3n+ 1} : n \in \mathbb{N}\right\}\).
        \item \(\left\{\frac{m}{m+ n} : m, n \in \mathbb{N}\right\}\).
    \end{enumerate}
\end{exercise}

\sol{To avoid writing out every set definition, I am going to denote each set as \(A_n\) where \(n\) corresponds to the numerical value of the list from (a) - (d).
    \begin{enumerate}
        \item \(\sup A_1 = 1\), \(\inf A_1 = 0\)
        \item \(\sup A_2 = 1\), \(\inf A_2 = -1\)
        \item \(\sup A_3 = \frac{1}{3}\), \(\inf A_3 = \frac{1}{4}\)
        \item \(\sup A_4 = 1\), \(\inf A_3 = 0\)
    \end{enumerate}
}

\renewcommand{\theenumi}{\arabic{enumi}}
\renewcommand{\labelenumi}{\theenumi.}
\section{Consequences of Completeness}

\begin{ntheorem}
    {Nested Interval Property}For each \(n \in \N\), assume we are given a closed interval \(I_n = [a_n, b_n]\). Assume \(I_n\) contains \(I_{n+1}\). This results in a nested sequence of intervals. \[I_1 \supseteq I_2 \supseteq I_3 \supseteq I_4\dots\] Then, \(\bigcap^\infty_{n=1} \ne \emptyset\).
\end{ntheorem}

\textbf{tl;dr} there has to be something that is common to all of the sets.

\tpf{
    Notice that the sequence, \(a_1,a_2,a_3,\dots\) is increasing. In other words, for each \(n \in \N\), since \(I_n \supset I_{n + 1}\) we have \(a_n \leq a_{n+1}\). If we consider the set \(A = \{a_n \ \colon \ n\in \N\}\). The element \(b_1\) is an upper bound of \(A\). (Note that \(b_1\) and \(a_1\) corresponds to the endpoints of the first set, \(I_1\). Think of this as a tornado looking structure where the larger the \(I_n\), the smaller the number line.) For each \(n \in \N\), \(a_n \leq b_n \leq b_1\). \\

    Since \(A\) has an upper bound, it must have a least upper bound. Hence, let \(\alpha = \sup A\). We claim that \(\alpha \in \bigcap^\infty_{n=1} I_n\). We said \(b_1\) was an upper bound. In fact, every \(b_n\) is an upper bound of \(A\). Choose any \(n,m \in \N\). We want to show that \(a_n \leq b_m\). Consider the following cases: \\

    \textbf{Case 1:} If \(n < m\), then \(a_n \leq a_m \leq b_m\). (Think: two number lines stacked on top of each other. The top number line is larger, call it \(I_n\) and it has \(a_n\) and \(b_n\) as endpoints. Consider a contained line (\(I_n \supseteq I_m\)) that is smaller, and has endpoints \(a_m\) and \(b_m\).) \\

    \textbf{Case 2:} If \(n > m\), then \(a_n \leq b_n \leq b_m\). So every \(b_n\) is an upper bound of \(A\).\\

    Hence,
    \begin{itemize}
        \item Because \(\alpha = \sup A\), we have \(\alpha \geq a_n\).
        \item Since \(b_n\) is an upper bound of \(A\), we have \(\alpha \leq b_n\).
    \end{itemize}
    so, \(\alpha \in [a_n, b_n] = I_n\). Thus, \(\alpha \in \bigcap^\infty_{n=1} I_n.\)
}

Nested, closed, Bounded Intervals \(\Rightarrow\) non-empty intersection.

\begin{ntheorem}
    {Archimedean Principle}
    \begin{enumerate}
        \item Given any number \(x \in R\), there exists an \(n \in N\) satisfying \(n > x\).
        \item Given any real number \(y > 0\), there exists an \(n \in N\) satisfying \(1/n < y\).
    \end{enumerate}
\end{ntheorem}

\tpf{
    \begin{enumerate}
        \item If \(\N\) was bounded, then we can let \(s \in \N = \sup \N\). However, we know that there is always a higher number (e.g., \(n+1\)) for any \(n \in \N\) that is given. Thus, by contradiction, there must exist \(n \geq x\).
        \item For any \(x > 0\), there exists \(n \in \N\) such that \(\frac{1}{n} < x\).
    \end{enumerate}
}
\begin{ntheorem}
    {Density of the Rationals in the Reals}For any \(a,b \in \R\) with \(a < b\), there exists \(q \in \Q\) such that \(a < q < b\).
\end{ntheorem}

\tpf{
    Since \( b - a > 0\), there exists \(n \in \N\) such that \(\frac{1}{n} < b - a\). From the \namrefthm{Archimedean Principle}, since \(a\times n \in \R\), there exists \(m \in \N\) such that \(a\times n < m\). Let \(m\) be there smallest such natural numbers (by the well ordered principle). Since \(m\) is the smallest such natural number, it follows that \(m - 1 \leq a\times n < m\). We then see that \(a < \frac{m}{n}\). Now, we need to find some \(\frac{m}{n} < b\). \begin{align*}
        m - 1       & \leq a\times n       \\
        m           & \leq a \times n + 1  \\
        \frac{m}{n} & \leq a + \frac{1}{n} \\
        \frac{m}{n} & < a + (b - a)        \\
        \frac{m}{n} & < b
    \end{align*} We now have that \(a < \frac{m}{n} < b\) so \( \frac{m}{n}\) is a rational number in \((a,b)\)
}

\begin{exercise}
    {1.4.1}Recall that \(\mathbb{I}\) stands for the set of irrational numbers.
    \begin{enumerate}
        \item Show that if \(a,b \in \Q\), then \(ab\) and \(a + b\) are elements of \(\Q\) as well.
        \item Show that if \(a \in \Q\) and \(t \in \mathbb{I}\), then \(a + t \in \mathbb{I}\) and \(at \in \mathbb{I}\) as long as \(a \ne 0\).
        \item Part (a) can be summarized by saying that \(\Q\) is closed under addition and multiplication. Is \(\mathbb{I}\) closed under addition and multiplication? Given two irrational numbers \(s\) and \(t\), what can we say about \(s + t\) and \(st\)? In other words, are there two irrational numbers that can be added and multiplied such that you get a number \(x\) such that \(x \notin \mathbb{I}\).
    \end{enumerate}
\end{exercise}

\sol{
    \begin{enumerate}
        \item Let \(a,b \in \Q\). This means there exists some \(p,q,a,b \in \Z\) such that \[a = \frac{p}{q}\] and \[b = \frac{a}{b}\] where \(q,b \ne 0\). The product of these numbers is \[ab = \frac{p}{q} \cdot \frac{a}{b} = \frac{pa}{qb}.\] Since \(pa,qb \in \Z\), \(ab \in \Q\). The sum of these numbers is \[a + b = \frac{p}{q} + \frac{a}{b} = \frac{pb + aq}{qb}.\] Since \(pb + aq, qb \in \Z\), \(a + b \in \Q\).
        \item Let \(a \in \Q\) and \(t \in \I\). Assume, for contradiction, that \(a + t \in \Q\). This would imply \(t = (a + t) - a\) (because we can subtract \(t + a\) from the original equation and rearrange terms). Since \(a + t, a \in \Q\) their sum would be rational because the rational numbers are closed under addition. However, that would contradict the assumption that \(t \in \I\). Hence, \(a + t \in \I\).
        \item For \(\I\), it is not closed under addition and multiplication. Consider the following counterexample: \(\sqrt{2} + (-\sqrt{2}) = 0\) which is not in the irrationals. For multiplication, consider \(\sqrt{2} \cdot \sqrt{2} = 2\), which is also not in the irrationals.
    \end{enumerate}
}

\renewcommand{\theenumi}{\arabic{enumi}}
\renewcommand{\labelenumi}{\theenumi.}
\section{Cardinality}

Two sets have the same \textit{cardinality} if there exists a bijection between them. Thus, the natural numbers, the integers, and the rational numbers have the same cardinality. A set is \textit{countably infinite} if it has the same cardinality as \(\N\). (If it can be put into one-to-one correspondence with \(\N\).) A set is \textit{countable} if it is countably infinite or finite. \\

\setcounter{BoxCounter}{5}
\begin{theorem}
    \(\R\) is not countable.
\end{theorem}

\tpf{%
    \textit{1} (most common) \\
    Suppose \(\R\) is countable. Then we can list them all, or we can enumerate them. \(\R = \{x_1,x_2,x_3,x_4,\dots\}\). We can write the decimal expansion of each of these. Consider the following table:
    \begin{center}
        \begin{tabular}{lllllll}
            \(x_1 =\) & \(\boxed{a_{10}}\) & \(a_{11}\)         & \(a_{12}\)         & \(a_{13}\)         & \(a_{14} \)        & \dots \\ \midrule
            \(x_2 =\) & \(a_{20} \)        & \(\boxed{a_{21}}\) & \(a_{22} \)        & \(a_{23} \)        & \(a_{24} \)        & \dots \\ \midrule
            \(x_3 =\) & \(a_{30} \)        & \(a_{31} \)        & \(\boxed{a_{32}}\) & \(a_{33} \)        & \(a_{34} \)        & \dots \\ \midrule
            \(x_4 =\) & \(a_{40} \)        & \(a_{41} \)        & \(a_{42} \)        & \(\boxed{a_{43}}\) & \(a_{44} \)        & \dots \\ \midrule
            \(x_5 =\) & \(a_{50} \)        & \(a_{51} \)        & \(a_{52} \)        & \(a_{53} \)        & \(\boxed{a_{54}}\) & \dots \\ \midrule
            \(x_6 =\) & \(a_{60} \)        & \(a_{61} \)        & \(a_{62} \)        & \(a_{63} \)        & \(a_{64} \)        & \dots
        \end{tabular}
    \end{center}

    We will now construct a number that is not in this list. Focus on diagonal entries. For each \(n \in \N\), let \(b_n\) be a digit that is different fron \(a_{nn}\). Now consider the number \(y = 0.b_1b_2b_3b_4b_5\dots\). This number \(y\) is not in our list. So our list did not include all of \(\R\). Avoid repeating \(9\)s.
}

\tpf{
    \textit{2} (uses nested interval theorem) \\
    Suppose \(\R\) is countable. Then we can enumerate \(\R\) \(\R = \{x_1,x_2,x_3,\dots\}\). Let \(I_1\) be any closed interval that does not contain \(x_1\). Next, we will find another closed interval \(I_2\) that:\begin{itemize}
        \item \(I_2 \subseteq I_1\)
        \item \(x_2 \notin I_2\)
    \end{itemize} Continue in this fashion creating a sequence of nested closed intervals: \(I_1\supseteq I_2 \supseteq I_3 \supseteq \dots\) such that for all \(k \in \N\), \(x_k \notin I_k\). Now consider: \[\bigcap^\infty_{n=1} I_n\]
    \begin{itemize}
        \item For each \(k \in \N\), since \(x_k \notin I_k\), we see \(x_k \notin \bigcap^\infty_{n=1} I_n\).
        \item By the nested interval theorem, there exists \(x \in \R\) such that \(x \in \bigcap^\infty_{n=1} I_n\). So \(x\) is a real number that is not included in our list.
    \end{itemize}
}

\begin{theorem}
    A countable collection of finite sets is \textit{countable}.
\end{theorem}

\begin{theorem}
    \begin{enumerate}[label=(\roman*)]
        \item The union of two countable sets is \textit{countable}.
        \item A countable union of countable sets is \textit{countable}.
    \end{enumerate}
\end{theorem}

From Theorem \ref{thm:1.5.6}, we know that \(\R\) is uncountable, but what about \((0,1)\)? It does have the same cardinality of \(\R\) because we can make a one-to-one and onto function between both the sets. Similarly, \((a,b)\) also has the same cardinality. What about \([a,b]\)?

\textbf{Recap:} \(\N\) is countable, and \(\R\) is uncountable and has a different cardinality than \(\N\). Thus, the question is, do all uncountable sets have the same cardinality as \(\R\)? The answer is \textbf{no}.

\begin{ntheorem}
    {Canter's Theorem} For any set \(A\), there does not exist an onto map from \(A\) into \(\mathcal{P}\).
\end{ntheorem}

\tpf{
    Suppose there exists an onto function, \(f \ \colon \ A \rightarrow \mathcal{P}(A)\). So each \(a \in A\) is mapped to an element \(f(a) \in \mathcal{P}(A)\). Then, \(f(a) \subseteq A\). We are going to construct an element of \(\mathcal{P}(A)\) which is not mapped to by \(f\).

    Consider \(B = \{a \in A \ \colon \ a \notin f(a)\}\). Since \(f\) is onto there exists \(a' \in A\) such that \(B = f(a')\). Thus, there are two cases to consider:

    \begin{itemize}
        \item \textbf{Case 1:} If \(a' \in B = f(a')\), then \(a' \notin B\).
        \item \textbf{Case 2:} If \(a' \notin B = f(a')\), then \(a' \in B\).
    \end{itemize}

    As evidenced, both cases lead to contradictions, so \(B\) is not the image of any \(a \in A\). Therefore \(f\) is not onto.
}

\begin{example}
    {Set and Power Set Matching}\(A = \{a,b,c\}\).
\end{example}

\sol{
    \(\mathcal{P}(A) = \emptyset, \{a\}, \{b\}, \{c\}, \{a,b\}, \{a,c\}, \{b,c\}, \{a,b,c\}\). Note that you can map \(\{a\}\), \(\{b\}\), \(\{c\}\), to elements such as \(\emptyset\), \(\{a,b\}\), \(\{a,b,c\}\), but there are still more elements that are left unmapped. We can extrapolate from our proof a set \(B\) such that \(B = \{a,c\}\) because those elements are not mapped to.
}

All of this is to show \(\mathcal{P}(\R)\) has a larger cardinality than \(\R\). Then \(\mathcal{P}(\mathcal{P}(\R))\) has a larger cardinality than \(\mathcal{P}(\R)\).

\setcounter{chapter}{1}
\chapter{Sequences and Series}\label{sequences and series}
\vspace*{-0.25in}
\renewcommand{\theenumi}{\arabic{enumi}}
\renewcommand{\labelenumi}{\theenumi.}
\section{Discussion: Rearrangement of Infinite Series}








\noindent \textbf{Questions:}

\textbf{What is a \textit{sequence}?} \\

A countable, ordered list of elements. An example could be \(1,2,3,4,5,\dots\). Note that this is \textit{ordered}, therefore distinguishing it from a sequence like \(3,1,2,4,5,6,\dots\). Hence, order matters.

% This needs to be a part of the definition environment
A \textit{sequence} is a function whose domain is \(\N\). \textbf{Note:} The domain \(\N\) refers to each element's position in the list. For example, \((a_n) = a_1,a_2,a_3,\dots\).

We will focus on the \textit{limit} of a sequence. We use sequences to approximate other things.

\begin{example}
    {Sequence} \(3,3.1,3.14,3.141,3.1415,\dots \approx \pi\). \\
    \(x, x - \frac{x^3}{6}, x - \frac{x^3}{6} + \frac{x^5}{120}, x - \frac{x^3}{6} + \frac{x^5}{120} - \frac{x^7}{7!}, \dots \approx \sin(x)\)
\end{example}

\noindent\textbf{What is a \textit{series}?} \\

An infinite sum. We look at the sequence of partial sums. We ask, do the partial sums approach a limit?

\begin{example}
    {Alternating Harmonic Series} \(1 - 1/2 + 1/3 - 1/4 + 1/5 - 1/6 + \dots\)
\end{example}

We can rearrange these terms such that we can `force' the series to converge to a specific number. Therefore, we will need to be careful with our definitions.



\renewcommand{\theenumi}{\arabic{enumi}}
\renewcommand{\labelenumi}{\theenumi.}
\section{The Limit of a Sequence}

\begin{definition}
    A \textit{sequence} is a function whose domain is \(\N\). We write \((a_n) = a_1,a_2,a_3,\dots\).
\end{definition}
\setcounter{BoxCounter}{2}
\begin{definition}
    The sequence \((a_n)\) \textit{converges} to \(L\) if for all \(\epsilon > 0\), there exists \(N \in \N\) such that for all \(n \geq \N\), \(|a_n - L| < \epsilon\). In other words, there exists \(N \in \N\) such that
    \begin{itemize}
        \item \textbf{(In the interval)} \(a_N \in (L - \epsilon, L + \epsilon)\).
        \item \textbf{(Stays in the interval)} \(\forall n \geq N\), \(a_n \in (L - \epsilon, L + \epsilon)\).
    \end{itemize}
\end{definition}

\begin{example}
    {Limit Proof 1}Let \(a_n = \frac{1}{n}\). \((a_n) = (1,\frac{1}{2},\frac{1}{3},\dots\))
\end{example}

\epf{
    Our claim is \(\displaystyle \lim_{n \rightarrow \infty} \frac{1}{n} = 0\). Thus, let \(\epsilon = .01\). Does the sequence eventually get inside \((-.01,.01)\)? We will set \(N = 101\). So, for any \(n \geq |0|\), \[\left|\frac{1}{n} - 0\right| = \frac{1}{n} \leq \frac{1}{101} < .01.\] From \(A_n\) and on, the sequence stayed within \(\epsilon\) of \(0\). But what about \(\epsilon = .001\), \(\epsilon = .00001\) and so on? \\

    Actual proof let \(\epsilon > 0\). By the Archimedean Principle, there exists \(N \in \N\) such that \(N > \frac{1}{\epsilon}\). Now, for any \(n \geq N\), \[\left|\frac{1}{n}\right| = \frac{1}{n} \leq \frac{1}{N} < \frac{1}{1/\epsilon}.\] (Where \(\frac{1}{1/\epsilon} = \epsilon\), but is in that form for demonstration purposes.) Therefore \(\displaystyle \lim_{n \rightarrow \infty} \frac{1}{n} = 0\)
}

``To get close'' means is that we are finding a bigger and bigger \(N\) as \(\epsilon\) gets smaller. Note that the choice of \(N\) certainly depends on \(\epsilon\). This idea of ``getting close'' can be seen in the following definition:

\begin{tcolorbox}%
    [enhanced, sharp corners=downhill, colframe=xred, colback=xred!15, coltitle=white, label=def:2.2.3B, title={\large \textbf{Definition 2.2.3B}},attach boxed title to top right={xshift=-0.5cm}, boxed title style={colback=xred}]
    A sequence \( (a_n) \) converges to \( a \) if, given any \( \epsilon \)-neighborhood \( V_\epsilon(a) \) of \( a \), there exists a point in the sequence after which all of the terms are in \( V_\epsilon(a) \). In other words, every \( \epsilon \)-neighborhood contains all but a finite number of the terms of \( (a_n) \).
\end{tcolorbox}

\subsection{Basic Structure of a Limit Proof}

Claim: \(\displaystyle \lim_{n \rightarrow \infty} a_n = L\). \\
Proof: Let \(\epsilon > 0\). There exists \(N \in \N\) such that \{something involving \(\epsilon\)\}. Assume \(n \geq N\). Then, \[|a-n - L| \ \boxed{\dots} < \epsilon\] (Where \(\boxed{\dots}\) is going to be where the majority of the work is going to lie.)

\begin{example}
    {Limit Proof 2}Claim: \(\displaystyle \lim_{n \rightarrow \infty} \frac{2n-3}{2n} = 1\)
\end{example}

\epf{
    Let \(\epsilon > 0\). \textit{Scratch paper:} Solve for: \[\left|\frac{2n - 3}{2n} - 1\right| = \left|\frac{-3}{2n}\right| = \frac{3}{2n} < \epsilon \Rightarrow \frac{3}{2\epsilon} < n.\] By the \namrefthm{Archimedean Principle}, there exists \(N \in \N\) such that \(N > \frac{3}{2\epsilon}\). Assume \(n \geq N\), (want to know what happens past this point) \[\left|\frac{2n-3}{2n}-1\right| \leq \frac{3}{2N} < \frac{3}{2 \cdot 3/2\epsilon} = \epsilon.\] Therefore, \(\displaystyle \lim_{n \rightarrow \infty} \frac{2n-3}{2n} = 1\)
}

\begin{example}
    {Limit Proof 3}Claim: \(\displaystyle \limn \frac{2n^2 + 1}{n^2} = 2\)
\end{example}

\epf{
Let \(\epsilon > 0\). By the \namrefthm{Archimedean Principle}, there exists \(N \in \N\) such that [leave off] \textit{Scratch paper:} Solve for \[\disabs{\frac{2n^2 + 1}{n^2} - 2} =  \frac{2n^2}{n^2} < \epsilon \Rightarrow \frac{3}{2\epsilon} < n\] [pick up] there exists \(N \in \N\) such that \[N > \frac{1}{\sqrt{\epsilon}}.\] Assume \(n \geq N\), then
\begin{align*}
    \disabs{\frac{2n^2 + 1}{n^2} - 2} & = \frac{1}{n^2}                                \\
                                      & \leq \frac{1}{N^2}                             \\
                                      & < \frac{1}{\left(1/(\sqrt{\epsilon})\right)^2} \\
                                      & = \frac{1}{1/\epsilon}                         \\
                                      & = \epsilon
\end{align*}
Therefore, \(\limn \frac{2n^2 + 1}{n^2} = 2\)
}

\begin{example}
    {Limit Proof 4}Claim: \(\displaystyle \limn \frac{7n + 8}{3n + 6} = \frac{7}{3} \)
\end{example}

\epf{
    \begin{align*}
        \disabs{\frac{7n+8}{3n+6} - \frac{7}{3}} & = \disabs{\frac{21n + 24}{3(3n + 6)} - \frac{21n + 42}{3(3n + 6)}} \\
                                                 & = \disabs{\frac{-18}{9n + 18}}                                     \\
                                                 & = \frac{18}{9n + 18} < \epsilon**                                  \\
                                                 & = \frac{18}{3} < 9n + 18                                           \\
                                                 & = \frac{18}{3} - 18 < 9n                                           \\
                                                 & = \frac{18/\epsilon - 18}{9} < n
    \end{align*}
    \(** \frac{18}{9n + 8} < \frac{18}{9n} < \epsilon \Rightarrow \frac{2}{\epsilon} < N\). \(\exists N \in \N\) such that \(N > \frac{2}{\epsilon}\). Assume \(n \geq N\),
    \begin{align*}
        \disabs{\frac{7n + 8}{3n + 6} - \frac{7}{3}} & = \frac{18}{9n + 18}   \\
                                                     & = \frac{2}{n+2}        \\
                                                     & < \frac{2}{n}          \\
                                                     & \leq \frac{2}{N}       \\
                                                     & < \frac{2}{\epsilon/2} \\
                                                     & = \epsilon. \qedhere
    \end{align*}
}

\textbf{Does every sequence have a limit?}

\begin{ntheorem}
    {Uniqueness of Limits}The limit when it exists, is unique.
\end{ntheorem}


\tpf{
    Let \((x_n)\) be a convergent sequence. Suppose \(L\) and \(M\) are limits of this sequence. Without the loss of generality, we are going to assume \(M > L\) Let \[\epsilon = \frac{M - L}{3}.\] Since \(n_x\) converges to \(L\), there exists \(N_1 \in \N\) such that for all \(n \geq N\), \(\disabs{(x_n) - L} < \epsilon\). Since \((x_n)\) converges to \(M\), there exists an \(N_2 \in \N\) such that for all \(n \geq N_2\), \(\disabs{(x_n) - M} < \epsilon\). Consider \(n = \max\{N_1, N_2\}\). Since \(n \geq N_1\), \(\disabs{(x_n) - L} < \epsilon\). Since \(n \geq N_2\), \(\disabs{(x_n) - M} < \epsilon\). Then \(L - \epsilon < x_n < L + \epsilon\) and \(M - \epsilon < x_n < M + \epsilon\). By our choice of \(\epsilon\), we now have \[(x_n) < L + \epsilon < M - \epsilon < (x_n).\] This is a contradiction. Thus, \((x_n)\) cannot have two different limits.
}

\begin{example}
    {Limit Proof 5}Let \((x_n) = \displaystyle \frac{\cos(n)}{3n}\). Claim: \(\displaystyle \limn (x_n) = 0\)
\end{example}

\epf{
    Let \(\epsilon > 0\). By the \namrefthm{Archimedean Principle}, there exists \(N \in \N\) such that \(N > \frac{1}{3\epsilon}\) for all \(n \geq N\),
    \begin{align*}
        \disabs{\frac{\cos(n)}{3n} - 0} & = \disabs{\frac{\cos(n)}{3n}} \\
                                        & \leq \frac{1}{3n}             \\
                                        & \leq \frac{1}{3N}             \\
                                        & < \frac{1}{3(1/ 3\epsilon)}   \\
                                        & = \epsilon
    \end{align*}
}

\begin{example}
    {Limit Proof 6}Let \(\displaystyle (y_n) = \frac{4n-1}{n^2}\). Claim: \(\displaystyle \limn y_n = 0\).
\end{example}

\epf{
    Let \(\epsilon > 0\). By the \namrefthm{Archimedean Principle}, there exists \(N \in \N\) such that \(N > \frac{1}{\epsilon}\). For all \(n \geq N\),
    \begin{align*}
        \disabs{\frac{4n-1}{n^2} - 0} & = \disabs{\frac{4n-1}{n^2}} \\
                                      & = \frac{4n-1}{n}            \\
                                      & < \frac{4n}{n^2}            \\
                                      & = \frac{4}{n}               \\
                                      & \leq \frac{4}{N}            \\
                                      & < \frac{4}{4 / \epsilon}    \\
                                      & = \epsilon
    \end{align*}
}

% \renewcommand{\theenumi}{\alph{enumi}}
% \renewcommand{\labelenumi}{(\theenumi)}
% \subsection{Exercises}

% \begin{exercise}
%     {2.1.1}What happens if we reverse the order of the quantifiers in \defref{2.2.3}?

%     \textit{Definition:} A sequence \(x_{n}\) \textit{verconges} to \(x\) if there exists an \(\epsilon >0\) such that for all \(N \in \N\) it is true that \(n \geq N\) implies \(\disabs{x-n - x} < \epsilon \).
%     \begin{enumerate}
%         \item Give an example of a vercongent sequence.
%         \item Is there an example of a vercongent sequence that is divergent?
%         \item Can a sequence verconge to two different values?
%         \item What exactly is being described in this strange definition?
%     \end{enumerate}
% \end{exercise}

% \sol{
%     \begin{enumerate}
%         \item Pick \(\epsilon = 2\), \(x_{n} = (-1)^{n}\) and \(x = 0\). This sequence will stay within the bounds of \((-2,2)\) for all \(N \in \N\) and \(n \geq N\).
%         \item There cannot be a divergent vercongent sequence because vercongence wants us to be bounded, and divergence wants it to grow outside the bounds. These two ideas are mutually exclusive.
%         \item Yes. For example, \(x_{n} = 0\) and \(x_{n} = 1\).
%         \item This definition is describing a sequence that is bounded. It is a sequence that is not growing outside of a certain range.
%     \end{enumerate}
% }

% \begin{exercise}
%     {2.2.2}Verify, using \defref{2.2.3}, that
%     the following sequences converge to the proposed limit.
%     \begin{enumerate}
%         \item \(\limn \frac{2n + 1}{5n + 4} = \frac{2}{5}\).
%         \item \(\limn \frac{2n^2}{n^3+3} = 0\)
%     \end{enumerate}
% \end{exercise}

% \expf{\hfill
%     \begin{enumerate}
%         \item Let \(\epsilon > 0\). Choose \(N \in \N\) such that \(N > \dfrac{3}{25\epsilon}\). Then for all \(n \geq N\),
%               \begin{align*}
%                   \left|\frac{2n + 1}{5n + 4} - \frac{2}{5}\right| & = \left|\frac{-3}{5(5n + 4)}\right| \\
%                                                                    & = \frac{3}{25n + 20}                \\
%                                                                    & \leq \frac{3}{25n}                  \\
%                                                                    & \leq \frac{3}{25N}                  \\
%                                                                    & < \epsilon
%               \end{align*}
%               Therefore, \(\displaystyle \lim_{n \to \infty} \frac{2n + 1}{5n + 4} = \frac{2}{5}\).
%         \item Let \(\epsilon > 0\). By the \namrefthm{Archimedean Principle}, there exists an \(N \in \N\) such that \(N > \frac{2}{\epsilon}\). Then, for \(n \geq N\),
%               \begin{align*}
%                   \disabs{\frac{2n^2}{n^3 + 3} - 0} & = \disabs{\frac{2n^2}{n^3 + 3}} \\
%                                                     & = \frac{2n^2}{n^3 + 3}          \\
%                                                     & < \frac{2n^2}{n^3}              \\
%                                                     & = \frac{2}{n}                   \\
%                                                     & \leq \frac{2}{N}                \\
%                                                     & = \frac{2}{2/\epsilon}          \\
%                                                     & = \epsilon.
%               \end{align*}
%               Therefore, \(\limn \frac{2n^2}{n^3 + 3} = 0\). \qedhere
%     \end{enumerate}
% }

% \begin{exercise}
%     {2.2.3} Describe what we would have to demonstrate in order to disprove each of the following statements.
%     \begin{enumerate}
%         \item At every college in the United States, there is a student who is at least
%               seven feet tall.
%         \item For all colleges in the United States, there exists a professor who gives every student a grade of either A or B.
%         \item There exists a college in the United States where every student is at least six feet tall.
%     \end{enumerate}
% \end{exercise}

% \sol{
%     \begin{enumerate}
%         \item There is at least one college in the United States where all students are less than seven feet tall.
%         \item There is at least one college in the United States where all professors give at least one student a grade of C or lower.
%         \item For all colleges in the United States, there exists a student who is less than six feet tall.
%     \end{enumerate}
% }

% \begin{exercise}
%     {2.2.4} Give an example of each or state that the request is impossible.
%     For any that are impossible, give a compelling argument for why that is the case.
%     \begin{enumerate}
%         \item A sequence with an infinite number of ones that does not converge to one.
%         \item A sequence with an infinite number of ones that converges to a limit not equal to one.
%         \item A divergent sequence such that for every \(n \in \N\) it is possible to find \(n\) consecutive ones somewhere in the sequence.
%     \end{enumerate}
% \end{exercise}

% \sol{
%     \begin{enumerate}
%         \item Possible. Consider the sequence \(a_n = \begin{cases}
%                   1 & \text{if } n \text{ is even} \\
%                   0 & \text{if } n \text{ is odd}
%               \end{cases}\). This sequence has infinitely many ones but does not converge to one.

%         \item Impossible. Suppose \((a_n)\) is a sequence that converges to a limit \(L \neq 1\) and has infinitely many ones. Since \((a_n)\) converges to \(L\), for any \(\epsilon > 0\), there exists \(N \in \N\) such that for all \(n \geq N\), \(|a_n - L| < \epsilon\). Choose \(\epsilon = \frac{|1 - L|}{2} > 0\). Then, for \(n \geq N\), \(|a_n - L| < \epsilon\), which implies \(a_n \neq 1\) beyond this \(N\). This contradicts the existence of infinitely many ones. Therefore, such a sequence is impossible.

%         \item Possible. Define a sequence by concatenating increasing blocks of ones separated by zeros: \( (0, 1, 0, 1, 1, 0, 1, 1, 1, 0, \dots) \). Specifically, the sequence consists of \(n\) ones followed by a zero for \(n = 1, 2, 3, \dots\). For every \(n \in \N\), there is a block of \(n\) consecutive ones somewhere in the sequence. The sequence does not converge, so it is divergent.
%     \end{enumerate}
% }

% \begin{exercise}
%     {2.2.5} Let \([[x]]\) be the greatest integer less than or equal to \(x\). For example, \([[\pi]] = 3\) and \([[3]] = 3\). For each sequence, find \(\lim_{n \to \infty} a_{n}\) and verify it with the definition of convergence.
%     \begin{enumerate}
%         \item \(a_{n} = [[5/n]]\)
%         \item \(a_{n} = [[(12 + 4n)/3n]]\)
%     \end{enumerate}
%     Reflecting on these examples, comment on the statement following \defref{2.2.3B} that ``the smaller the \(\epsilon\)-neighborhood, the larger \(N\) may have to be.''
% \end{exercise}

% \sol{
%     \begin{enumerate}
%         \item We will show that \(\displaystyle \lim_{n \to \infty} a_{n} = 0\).

%               \innerpf{
%               For \(n \geq 6\), we have \(\displaystyle \frac{5}{n} \leq \frac{5}{6} < 1\), so \(a_{n} = [[5/n]] = 0\).

%               Let \(\epsilon > 0\). Choose \(N = 6\). Then for all \(n \geq N\),
%               \[
%                   |a_{n} - 0| = |0 - 0| = 0 < \epsilon.
%               \]
%               Therefore, by the definition of convergence, \(\displaystyle \lim_{n \to \infty} a_{n} = 0\).
%               }

%         \item We will show that \(\displaystyle \lim_{n \to \infty} a_{n} = 1\).

%               \innerpf{
%               Observe that:
%               \[
%                   a_{n} = \left[ \dfrac{12 + 4n}{3n} \right] = \left[ \dfrac{4n + 12}{3n} \right] = \left[ \dfrac{4}{3} + \dfrac{4}{n} \right].
%               \]
%               As \(n \to \infty\), \(\dfrac{4}{n} \to 0\), so \(\dfrac{4}{3} + \dfrac{4}{n} \to \dfrac{4}{3} \approx 1.333\).

%               For \(n \geq 7\), we have:
%               \[
%                   \dfrac{4}{n} \leq \dfrac{4}{7} \approx 0.571,\quad \dfrac{4}{3} + \dfrac{4}{n} \leq 1.333 + 0.571 = 1.904.
%               \]
%               Since \(1 < \dfrac{4}{3} + \dfrac{4}{n} < 2\) for \(n \geq 7\), we have:
%               \[
%                   a_{n} = \left[ \dfrac{4}{3} + \dfrac{4}{n} \right] = 1.
%               \]
%               Let \(\epsilon > 0\). Choose \(N = 7\). Then for all \(n \geq N\),
%               \[
%                   |a_{n} - 1| = |1 - 1| = 0 < \epsilon.
%               \]
%               Therefore, by the definition of convergence, \(\displaystyle \lim_{n \to \infty} a_{n} = 1\).
%               }
%     \end{enumerate}

%     \textbf{Reflection:} In these examples, we see that once the sequence reaches a certain point (i.e., \(n \geq N\)), the terms remain constant. This means that for any \(\epsilon > 0\), we can find a fixed \(N\) to satisfy the definition of convergence, regardless of how small \(\epsilon\) is. However, in general, smaller \(\epsilon\)-neighborhoods may require larger \(N\) because the sequence may not settle into its limit as neatly as it does in these cases.
% }

% \tcbset{highlight math style={enhanced,%<-- needed for the `remember' options
% colframe=red,colback=red!10!white,boxsep=0pt}}
% \begin{align*}
% \tcbhighmath[remember as=fx]{f(x)}
% &= \int\limits_{1}^{x} \frac{1}{t^2}~dt
% = \left[ -\frac{1}{t} \right]_{1}^{x}\\
% &= -\frac{1}{x} + \frac{1}{1}\\
% &=
% \tcbhighmath[remember,overlay={%
% \draw[blue,very thick,->] (fx.south) to[bend right] ([yshift=2mm]frame.west);}]
% {1-\frac{1}{x}.}
% \end{align*}

% \begin{exercise}
%     {2.2.6} Prove the \namrefthm{Uniqueness of Limits} theorem. To get started, assume \((a_{n}) \rightarrow a \) and \((a_{n}) \rightarrow b\). Now argue \(a = b\).
% \end{exercise}

% \expf{
% Since \((a_{n}) \rightarrow a\), this means for all \(\epsilon > 0\), there exists an \(N_{1} \in \N\) such that for all \(n \geq N\), \(\disabs{a_{n} - a} < \epsilon/2\). Similarly, since \((a_{n}) \rightarrow b\), this means for all \(\epsilon > 0\), there exists an \(N_{2} \in \N\) such that for all \(n \geq N_{2}\), \(\disabs{a_{n} - b} < \epsilon/2\).

% Now, let \(N = \max\{N_{1},N_{2}\}\) so that
% \begin{align*}
%     \disabs{a - b} & = \disabs{a - a_{n} + a_{n} - b}        \\
%                    & \leq \disabs{(a_{n} - a) + (a_{n} - b)} \\
%                    & < \epsilon/2 + \epsilon/2               \\
%                    & < \epsilon
% \end{align*}
% Then, by \numrefthm{1.2.6}, \(a = b\).
% }

% \begin{exercise}
%     {2.2.7} {Here are two useful definitions:}
%     \begin{enumerate}[label=(\roman*)]
%         \item A sequence \((a_n)\) is \emph{eventually} in a set \(A \subseteq \mathbb{R}\) if there exists an \(N \in \mathbb{N}\) such that \(a_n \in A\) for all \(n \geq N\).
%         \item A sequence \((a_n)\) is \emph{frequently} in a set \(A \subseteq \mathbb{R}\) if, for every \(N \in \mathbb{N}\), there exists an \(n \geq N\) such that \(a_n \in A\).
%               \begin{enumerate}[label=(\alph*)]
%                   \item Is the sequence \((-1)^n\) eventually or frequently in the set \(\{1\}\)?
%                   \item Which definition is stronger? Does frequently imply eventually, or does eventually imply frequently?
%                   \item Give an alternate rephrasing of \defref{2.2.3B} using either frequently or eventually. Which is the term we want?
%                   \item Suppose an infinite number of terms of a sequence \((x_n)\) are equal to \(2\). Is \((x_n)\) necessarily eventually in the interval \((1.9, 2.1)\)? Is it frequently in \((1.9, 2.1)\)?
%               \end{enumerate}
%     \end{enumerate}

% \end{exercise}

% \sol{
%     \begin{enumerate}[label=(\alph*)]
%         \item The sequence \((-1)^n\) is \textit{frequently} in the set \(\{1\}\) because for every \(N \in \mathbb{N}\), we can find an \(n \geq N\) such that \((-1)^n = 1\).
%         \item The definition of \textit{eventually} is stronger because \textit{eventually} implies \textit{frequently}, but \textit{frequently} does not imply \textit{eventually}.
%         \item An alternate rephrasing of Definition 2.2.3B using \textit{eventually} is: A sequence \((a_n)\) converges to \(a\) if, given any \(\epsilon\)-neighborhood---\(V_\epsilon(a)\) of \(a\)---\((a_n)\) is \textit{eventually} in \(V_\epsilon(a)\). The term we want is eventually.
%         \item If an infinite number of terms of a sequence \((x_n)\) are equal to \(2\),  \((x_n)\) is not \textit{eventually} in \((1.9, 2.1)\) because we can have a sequence \((a_n)\) that will not settle in \((1.9, 2.1)\). For example, \((a_n) = (0,2,0,2,\cdots)\) does not settle in \((1.9, 2.1)\). Whereas, \((x_n)\) is \textit{frequently} in the interval \((1.9, 2.1)\) because for every \(N \in \mathbb{N}\) there exists an \(n \geq N\) such that \(x_n \in (1.9, 2.1)\) for all \(n \geq N\). We can see an instance of this being true by examining the previous example.
%     \end{enumerate}
% }



\renewcommand{\theenumi}{\arabic{enumi}}
\renewcommand{\labelenumi}{\theenumi.}
\section{The Algebraic and Order Limit Theorems}

\begin{definition}
    A sequence \((x_n)\) is \textit{bounded} if there exists some \(M > 0\) such that every term in the sequence belongs to \([-M,M]\).
\end{definition}

\begin{theorem}
    Every convergent sequence is bounded.
\end{theorem}

\tpf{
Let \((x_n)\) be a convergent sequence with limit \(L\). There exists \(N \in \N\) such that for all \(n \geq N\), \(\disabs{(x_n) - L} < 1\). Equivalently, \((x_n) \in (L - 1, L + 1)\). Let \[M = \max\{|x_1|,|x_2|,\dots,|x_{N-1}|,|L + 1|, |L - 1|\}.\] We claim that for all \(n \in \N\), \(|x_n| \leq M\).
\begin{enumerate}
    \item This is true for \(n < N\).
    \item For \(n \geq N\), we know \(L - 1 < x_n < L + 1\), so \((x_n) \leq \max\{| L - 1 |, |L + 1|\} \)
\end{enumerate}
Thus, every term is in \([-M,M]\).
}

\begin{ntheorem}
    {Algebraic Limit Theorem}Let \(\limn a_n = a\) and \(\limn b_n = b\). Then,
    \begin{enumerate}[label=(\roman*)]
        \item \(\limn ca_n = ca\) for all \(c \in \R\);
        \item \(\limn (a_n + b_n) = a + b\);
        \item \(\limn (a_nb_n) = ab\);
        \item \(\limn \frac{a_n}{b_n} = \frac{a}{b}\) provided \(b \neq 0\).
    \end{enumerate}
\end{ntheorem}

\textit{Scratch Paper:}
\begin{align*}
    \disabs{ca_n - ca} = \disabs{c}\disabs{a_n - a} & < \epsilon                    \\
    \disabs{a_n - a}                                & < \frac{\epsilon}{\disabs{c}}
\end{align*}
\text{Leave off and go back to proof\(^1\)}
\tpf{(i) \\
Let \(\epsilon > 0\).\(^1\) Since \((a_n)\) converges to \(a\), there exists \(N \in \N\) such that for all \(n \geq N\), \(\disabs{a_n - a} < \frac{\epsilon}{\disabs{c}}\). Now, for any \(n \geq N\) we have two case because we want to avoid dividing by 0:
\begin{itemize}
    \item If \(c = 0\): \\
          then each \(ca_n = 0\). So \((ca_n)\) converges to 0, which can equal \(ca\).
    \item If \(c > 0\): \\
          \(\disabs{ca_n - ca} = \disabs{c}\disabs{a_n - a} < \disabs{c}\frac{\epsilon}{\disabs{c}} = \epsilon\).
\end{itemize}

(ii) \\
\textit{Scratch paper:}
\begin{align}
    \disabs{(a_n + b_n)} & = \disabs{(a_n - a) + (b_n - b)}          \\
                         & \leq \disabs{a_n - a} + \disabs{b_n - b}  \\
                         & < \frac{\epsilon}{2} + \frac{\epsilon}{2}
\end{align}
Note that (2.2) is from the triangle inequality. Now, we will pick up to back at \(\epsilon > 0\). \\

Let \(\epsilon > 0\). Since \((a_n)\) converges to \(a\), there exists \(N_1 \in \N\) such that for all \(n \geq N_1\), \(\disabs{a_n - a} < \frac{\epsilon}{2}\). Since \((b_n)\) converges to \(b\), there exists \(N_2 \in \N\) such that for all \(n \geq N_2\), \(\disabs{b_n - b} < \frac{\epsilon}{2}\). Now, let \(N = \max\{N_1, N_2\}\). Thus, for any \(n \geq N\), (refer back to scratch paper). \\

(iii) \\
\textit{Scratch paper:}
\begin{align}
    \disabs{a_nb_n - ab} & = \disabs{a_nb_n - ab_n + ab_n - ab}                                    \\
                         & = \disabs{a_n(b_n - b) + b(b_n - b)}                                    \\
                         & \leq \disabs{a_n}\disabs{b_n - b} + \disabs{b}\disabs{b_n - b}          \\
                         & \leq M\disabs{b_n - b} + M\disabs{a_n - a}.                             \\
                         & < M\left(\frac{\epsilon}{2M}\right) + M\left(\frac{\epsilon}{2M}\right) \\
                         & = \epsilon
\end{align}
Note that: (2.4) is where we added 0, (2.5) is from the triangle inequality, and (2.6) is just factored. Additionally, we choose \(N\) to get the fractions in (2.8) Now, we will pick up to back at \(\epsilon > 0\). \\

Let \(\epsilon > 0\). Since convergent sequences are bounded, then there exists \(M > 0\) such that for all \(n \in \N\), \(|a_n| \leq M\). We can choose \(M\) so that \(|b_n| \leq M\) as well. Since \((a_n)\) converges to \(a\), there exists \(N_1 \in \N\) such that for all \(n \geq N_1\), \(\disabs{a_n - a} < \frac{\epsilon}{2M}\). Since \((b_n)\) converges to \(b\), there exists \(N_2 \in \N\) such that for all \(n \geq N_2\), \(\disabs{b_n - b} < \frac{\epsilon}{2M}\). Now, let \(N = \max\{N_1, N_2\}\). Thus, for any \(n \geq N\), (refer back to scratch paper, and change (2.4)'s sign from an `=' to `\(\leq\)'). \\

(iv) \\
\textit{Scratch paper:}
\begin{align*}
    \disabs{\frac{a_n}{b_n} - \frac{a}{b}} & = \disabs{\frac{a_nb - ab_n}{b_nb}}                                                     \\
                                           & = \disabs{\frac{a_nb - ab_n + ab_n - ab}{b_nb}}                                         \\
                                           & = \disabs{\frac{a_n(b - b_n) + b(b_n - b)}{b_nb}}                                       \\
                                           & = \disabs{\frac{a_n(b - b_n)}{b_nb} + \frac{b(b_n - b)}{b_nb}}                          \\
                                           & \leq \disabs{\frac{a_n}{b_n}}\disabs{b - b_n} + \disabs{b}\disabs{\frac{b_n - b}{b_nb}} \\
                                           & < \epsilon
\end{align*}
Let \(\epsilon > 0\). Since \((b_n)\) converges to \(b\), there exists \(N_1 \in \N\) such that for all \(n \geq N_1\), \(\disabs{b_n } > \disabs{\frac{b}{2}}\). There also exists \(N_2 \in \N\) such that for all \(n \geq N_2\), \(\disabs{b_n - b} < \frac{\epsilon|b|^{2}}{2}\). Now, let \(N = \max\{N_1, N_2\}\). Let \(n \geq N\), (refer back to scratch paper). \qedhere
}

\begin{lemma}
    Let \((a_n)\) and \(c < a\). There exists \(N \in \N\) such that for all \(n \geq N\), \(a_n > c\). Similarly, if \(a < d\), there exists \(N \in \N\) such that for all \(n \geq N\), \(a_n < d\).
\end{lemma}

\subsection{Limits and Order}

\begin{ntheorem}
    {Order Limit Theorem}Let \((a_n)\) and \((b_n)\) be sequences. If \(\limn a_n = a\) and \(\limn b_n = b\), then
    \begin{enumerate}[label=(\roman*)]
        \item If \(a_n \geq c\) for all \(n \in \N\), then \(a \geq c\).
        \item If \(a_n \leq b_n\) for all \(n \in \N\), then \(a \leq b\).
        \item If there exists \( c \in \mathbb{R} \) for which \( c \leq b_n \) for all \( n \in \mathbb{N} \), then \( c \leq b \). Similarly, if \( a_n \leq c \) for all \( n \in \mathbb{N} \), then \( a \leq c \).

    \end{enumerate}
\end{ntheorem}



% \renewcommand{\theenumi}{\alph{enumi}}
% \renewcommand{\labelenumi}{(\theenumi)}
% \subsection{Exercises}

% \begin{exercise}
%     {2.3.1}
%     \begin{enumerate}
%         \item If \(\lim_{n \to \infty} x_n = 0\), show that \(\lim_{n \to \infty} \sqrt{x_n} = 0\).
%         \item If \(\lim_{n \to \infty} x_n = x\), show that \(\lim_{n \to \infty} \sqrt{x_n} = \sqrt{x}\).
%     \end{enumerate}
% \end{exercise}

% \expf{ \hfill
%     \begin{enumerate}
%         \item \textit{Solution.} Let \(\epsilon > 0\). Since \(\lim_{n \to \infty} x_n = 0\), there exists \(N \in \mathbb{N}\) such that for all \(n \geq N\),
%               \[
%                   |x_n| < \epsilon^2.
%               \]
%               Then, for all \(n \geq N\),
%               \[
%                   \left|\sqrt{x_n} - 0\right| = \sqrt{x_n} < \sqrt{\epsilon^2} = \epsilon.
%               \]
%               Therefore, \(\lim_{n \to \infty} \sqrt{x_n} = 0\).

%         \item \textit{Solution.} Let \(\epsilon > 0\). Since \(\lim_{n \to \infty} x_n = x\), for any \(\delta > 0\), there exists \(N \in \mathbb{N}\) such that for all \(n \geq N\),
%               \[
%                   |x_n - x| < \delta.
%               \]
%               We consider two cases:

%               \textbf{Case 1:} \(x > 0\).

%               Since \(x > 0\), choose \(\delta = \min\left\{\epsilon\left(2\sqrt{x}\right),\, \dfrac{x}{2}\right\}\). Then for all \(n \geq N\), we have \(x_n > x - \dfrac{x}{2} = \dfrac{x}{2} > 0\). Thus,
%               \[
%                   \sqrt{x_n} + \sqrt{x} \geq \sqrt{\dfrac{x}{2}} + \sqrt{x} > 0.
%               \]
%               Now,
%               \[
%                   \left|\sqrt{x_n} - \sqrt{x}\right| = \dfrac{|x_n - x|}{\sqrt{x_n} + \sqrt{x}} \leq \dfrac{\delta}{\sqrt{\dfrac{x}{2}}} \leq \epsilon.
%               \]

%               \textbf{Case 2:} \(x = 0\).

%               From part (1), we have \(\lim_{n \to \infty} \sqrt{x_n} = 0 = \sqrt{0}\).

%               Therefore, \(\lim_{n \to \infty} \sqrt{x_n} = \sqrt{x}\).
%     \end{enumerate}
% }

% \begin{exercise}
%     {2.3.2}Using only \defref{2.2.3}, prove that if \((x_n) \rightarrow 2\), then
%     \begin{enumerate}
%         \item \(\displaystyle \left(\frac{2x_n - 1}{3}\right) \rightarrow 1\);
%         \item \(\displaystyle \left(1 / x_n\right) \rightarrow 1 / 2\).
%     \end{enumerate}
%     (For this exercise the Algebraic Limit Theorem is off-limits, so to speak.)
% \end{exercise}

% \sol{ \hfill
%     \begin{enumerate}
%         \item \textit{Proof.} Let \(\epsilon > 0\). Since \((x_n)\) converges to 2, there exists \(N \in \N\) such that for all \(n \geq N\), \(\disabs{x_n - 2} < \epsilon\). Now, for any \(n \geq N\),
%               \begin{align*}
%                   \disabs{\frac{2x_n - 1}{3} - 1} & = \disabs{\frac{2x_n - 1 - 3}{3}} \\
%                                                   & = \disabs{\frac{2x_n - 4}{3}}     \\
%                                                   & = \frac{2}{3}\disabs{x_n - 2}     \\
%                                                   & < \disabs{x_n - 2}                \\
%                                                   & < \epsilon
%               \end{align*}
%               Therefore, \(\frac{2x_n - 1}{3} \rightarrow 1\) \qed
%         \item \textit{Proof.} Let \(\epsilon > 0\). Since \((x_n)\) converges to 2, there exists \(N_1 \in \N\) such that for all \(n \geq N_1\), \(x_n \geq 1\). Then, we will choose \(N_2\) so that \(|x_n - 2 | < \epsilon\) for all \(n \geq N_2\). Afterwards, we take \(N = \max\{N_1, N_2\}\). And note that for \(n \geq N\),

%               \begin{align*}
%                   \disabs{\frac{1}{x_n} - \frac{1}{2}} & = \disabs{\frac{2 - x_n}{2x_n}} \\
%                                                        & < \frac{\disabs{2 - x_n}}{2}    \\
%                                                        & < \frac{\epsilon}{2}            \\
%                                                        & < \epsilon
%               \end{align*} \qed
%     \end{enumerate}
% }

% \renewcommand{\theenumi}{\arabic{enumi}}
% \renewcommand{\labelenumi}{\theenumi.}
\section{The Monotone Convergence Theorem and a First Look at Infinite Series}

\begin{definition}
    A sequence \(a_n\) is \textit{increasing} if \(a_n \leq a_{n+1}\) for all \(n \in \N\) and \textit{decreasing} if \(a_n \geq a_{n+1}\) for all \(n \in \N\). A sequence is \textit{monotone} if it is either increasing or decreasing.
\end{definition}

\begin{ntheorem}
    {Monotone Convergence Theorem}If a sequence is monotone and bounded, then it converges.
\end{ntheorem}

\tpf{
    Let \((a_n)\) be an increasing and bounded sequence. Since \((a_n)\) is bounded, the set \(A = \{a_n \ |\ n \in \N\}\) is clearly also bounded. Since \(A\) is bounded, \(\sup{A}\) exists. We claim that \(\limn a_n = \sup{A}\). Thus, for all \(\epsilon > 0\) and by our definition of supremum, there exists \(N \in \N\) such that \(\sup{A} - \epsilon < a_N \leq \sup{A}\). Since \((a_n)\) is increasing, for all \(n \geq N\), \(\sup{A} - \epsilon < a_N \leq a_n \leq \sup{A}\). It follows that \(\disabs{a_n - \sup{A}} < \epsilon\). Therefore, \(\limn a_n = \sup{A}\).
}

\begin{example}
    {MCT}Consider the recursively defined sequence \(x_n\) where \(x_1 = 3\) and for all \(n \in \N\), \(x_{n + 1} = \frac{1}{4 - x_n}\). Show that \(x_n\) converges.
\end{example}

\epf{
    We will show that \(x_n\) is monotone and bounded.
    \begin{itemize}
        \item \textbf{Part 1: Monotone Decreasing}
              \begin{itemize}
                  \item \underline{Base case}: \(x_1 = 3\), \(x_2 = 1\).
                  \item \underline{Induction step}: Assume for some \(n \in \N\), \(x_n \geq x_{n+1}\). It follows that
                        \begin{align*}
                            x_n               & \geq x_{n + 1}               \\
                            4 - x_n           & \leq 4 - x_{n + 1}           \\
                            \frac{1}{4 - x_n} & \geq \frac{1}{4 - x_{n + 1}} \\
                            x_{n + 1}         & \geq x_{n + 2}
                        \end{align*}
              \end{itemize}
        \item \textbf{Part 2: Bounded Below}
              Claim: Sequence is bounded below by 0.
              \begin{itemize}
                  \item \underline{Base case}: \(x_1 = 3 > 0\).
                  \item \underline{Induction step}: Assume for some \(n \in \N\), \(x_n \geq 0\). It follows that \(4 - x_n \leq 4\), and when we take the reciprocal, we get \begin{align*}
                            \frac{1}{4 - x_n} & \leq \frac{1}{4} \\
                            x_{n+1}           & \geq 1/4         \\
                                              & > 0
                        \end{align*}
                        By math induction, \(x_n\) is bounded below by 0.
              \end{itemize}
    \end{itemize}
    By the Monotone Convergence Theorem, \(x_n\) converges.

    So, what is the limit? We know \((x_n)\) converges so let \(L = \limn x_n\). Then, \(\limn x_{n+1} = L\). We also know \(x_{n+1} = \frac{1}{4 - x_n}\). So \(L = \limn x_{n+1} = \limn \frac{1}{4- x_n} = \frac{1}{4 - L}\). It must be true that \(L = \frac{1}{4 - L}\). Solving for \(L\), we get \begin{align*}
        L(4 - L)     & = 1 \\
        4L - L^2     & = 1 \\
        L^2 - 4L + 1 & = 0
    \end{align*} Hence, \(L = 2 - \sqrt{3}\) or \(L = 2 + \sqrt{3}\). Notice that it cannot be the latter because it is bigger than 3.
}

\subsection{Recap and Summary}

We use limits to define multiple things in calculus. This is why we are focusing so heavily upon it. For example,

\begin{enumerate}
    \item Derivatives: \(\lim_{h \to 0} \frac{f(x + h) - f(x)}{h}\)
    \item Integrals: \(\lim_{n \to \infty} \sum_{i = 1}^{n} f(x_i) \Delta x\)
    \item Infinite Series: \(\lim_{n \to \infty} \sum_{i = 1}^{n} a_i\)
          Consider geometric series, \(C_a\) such that each term is multiplied by a ratio \(r\). This is represented as \(\sum_{n=0}^\infty ar^n = 1 + r + r^2 + r^3\dots\). When we look at partial sums, we get \(S_n = 1 + r + r^2 + r^3 + \dots + r^n\). We can then multiply by \(r\) to get \(rS_n = r + r^2 + r^3 + r^4 + \dots + r^{n+1}\). Subtracting the two, we get \((1 - r)S_n = 1 - r^{n+1}\). Thus, \[S_n = \frac{1 - r^{n+1}}{1 - r}.\] If \(|r| < 1\), then \(\lim_{n \to \infty} r^n = 0\). Thus, \(\lim_{n \to \infty} S_n = \frac{1}{1 - r}\).
\end{enumerate}

Looking to the future, we are going to use functions and summations together. For example, when we have \(f(x) = \sum_{n=0}^\infty (a_n)x^n\) such that \(f'(x) = \sum_{n=0}^\infty (a_n)x^{n-1}\).



\begin{definition}
    Let \((x_n)\) be a bounded sequence. Then the \textit{limit inferior} is \(\liminf_{n \to \infty} x_n = \lim_{n \to \infty} \inf\{x_k \ | \ k \geq n\}\). This is the largest a limit can get. The \textit{limit superior} is \(\limsup_{n \to \infty} x_n = \lim_{n \to \infty} \sup\{x_k \ | \ k \geq n\}\). This is the smallest a limit can get.
\end{definition}

\textbf{See \hyperref[ex:2.4.7]{Exercise 2.4.7} in the book for more information.}

\begin{example}
    {Monotone Decreasing Sequence}The following sequence is an example of a monotone decreasing sequence.
    \begin{align*}
        x_1, x_2, x_3, x_4, x_5, x_6, & \dots \sup\{x_k \mid k \geq 1\} = S. \\
        x_2, x_3, x_4, x_5, x_6,      & \dots \sup\{x_k \mid k \geq 2\} = S. \\
        x_3, x_4, x_5, x_6,           & \dots \sup\{x_k \mid k \geq 3\} = S. \\
        x_4, x_5, x_6,                & \dots \sup\{x_k \mid k \geq 4\} = S.
    \end{align*}
\end{example}

\(\limsup_{n \to \infty} x_n\) is guaranteed to exist by the \namrefthm{Monotone Convergence Theorem}.

\begin{example}
    {liminf} Let \(x_n = (-1)^n(1 + \frac{1}{n})\). Thus, \(x_{1,2,3} = -2,1\frac{1}{2},-1\frac{1}{3}\dots\)
\end{example}

\begin{example}
    {Convergence Towards 0}Let \(x_n = (-1)^n\frac{1}{n}\). Thus, \(x_{1,2,3} = -1, \frac{1}{2}, -\frac{1}{3}\dots\).
\end{example}

\begin{theorem}
    A sequence \(x_n\) is convergent if, and only if, \(\liminf_{n \to \infty} x_n = \limsup_{n \to \infty} x_n\).
\end{theorem}

\textbf{See Theorem 2.4.6 in the book for another view.}


% \renewcommand{\theenumi}{\alph{enumi}}
% \renewcommand{\labelenumi}{(\theenumi)}
% \subsection{Exercises}

% \begin{exercise}
%     {2.4.7 (Limit Superior)}Let \((a_{n})\) be a bounded sequence.
%     \begin{enumerate}
%         \item Prove that the sequence defined by \(y_{n} = \sup\{a_{k} \colon k \geq n\}\) converges.
%         \item The \textit{limit superior} of \((a_{n})\) or \(\limsupn{a_{n}}\), is defined by \[
%                   \limsupn{a_{n}} = \lim_{n \rightarrow \infty} y_{n},
%               \]
%               where \(y_{n}\) is the sequence from part (a) of this exercise. Provide a reasonable definition for \(\liminfn{a_{n}}\) and briefly explain why it always exists for any bounded sequence.
%         \item Prove that \(\liminfn{a_{n}} \leq \limsupn{a_{n}}\) for every bounded sequence, and give an example of a sequence for which the inequality is strict.
%         \item Show that \(\liminfn{a_{n}} = \limsupn{a_{n}}\) if and only if \(\lim a_{n}\) exists. In this case, all three share the same value.
%     \end{enumerate}
% \end{exercise}

% \begin{align*}
%     \left|\frac{3n - 1}{2n + 1} - \frac{3}{2}\right| & = \left|\frac{2(3n - 1) - 3(2n + 1)}{2(2n + 1)}\right|   \\
%                                                      & = \left|\frac{6n - 2 - 6n - 3}{4n + 2}\right|            \\
%                                                      & = \left|\frac{-5}{4n + 2}\right|                         \\
%                                                      & = \frac{5}{4n + 2}                                       \\
%                                                      & = \frac{5}{4N} \quad \text{Let } N = \frac{5}{4\epsilon} \\
%                                                      & = \epsilon
% \end{align*}

% The sequence \((x_n)\) is given by \(x_n = \frac{3n - 1}{2n + 1}\). Determine its limit and give a formal proof.

% I have:
% Let \(\epsilon > 0\). We want
% \[
%     \left|\frac{3n - 1}{2n + 1} - \frac{3}{2}\right| = \frac{5}{4n + 2} < \epsilon.
% \]
% For large \(n\): \(\frac{1}{4n + 2} < \frac{1}{4n}\) so
% \[
%     \frac{5}{4n + 2} < \frac{5}{4n}.
% \]
% To ensure \(\frac{5}{4n} < \epsilon\), choose \(N\) such that \(N > \frac{5}{4\epsilon}\). The Archimedean Principle ensures this number.

% For \(n \geq N\), we have \(n > \frac{5}{4\epsilon}\), so \(\frac{5}{4n} \leq \frac{5}{4N}\). Since \(N > \frac{5}{4\epsilon}\), we know \(\frac{5}{4N} < \epsilon\). Therefore,
% \[
%     \left|\frac{3n - 1}{2n + 1} - \frac{3}{2}\right| = \frac{5}{4n + 2} < \epsilon.
% \]

% \sol{
%     \begin{enumerate}
%         \item We will show that \((y_{n})\) converges.

%               \innerpf{
%                   Since \((a_n)\) is bounded, there exists \(M > 0\) such that \(|a_n| \leq M\) for all \(n\).

%                   For each \(n\), define \(y_n = \sup\{a_k : k \geq n\}\). As \(n\) increases, the set \(\{a_k : k \geq n\}\) becomes smaller, so the supremum cannot increase. Therefore, the sequence \((y_n)\) is non-increasing:
%                   \[
%                       y_{n+1} \leq y_n \quad \text{for all } n.
%                   \]
%                   Additionally, since \((a_n)\) is bounded below, so is \((y_n)\). Therefore, \((y_n)\) is a bounded, non-increasing sequence.

%                   By the Monotone Convergence Theorem, every bounded, monotonic sequence converges. Thus, \((y_n)\) converges.
%               }

%         \item A reasonable definition for \(\liminfn{a_{n}}\) is to define \(z_n = \inf\{a_k \colon k \geq n\}\) for each \(n\). Then, the \textit{limit inferior} of \((a_n)\) is defined by:
%               \[
%                   \liminfn{a_{n}} = \lim_{n\to\infty} z_n.
%               \]
%               Since \((a_n)\) is bounded, each \(z_n\) exists and the sequence \((z_n)\) is non-decreasing. As \(n\) increases, the set \(\{a_k : k \geq n\}\) becomes smaller, so the infimum cannot decrease. Therefore, \((z_n)\) is a bounded, non-decreasing sequence, which converges by the \namrefthm{Monotone Convergence Theorem}. Hence, \(\liminfn{a_{n}}\) always exists for any bounded sequence.

%         \item We will show that \(\liminfn{a_{n}} \leq \limsupn{a_{n}}\) for every bounded sequence.
%               \innerpf{
%                   For each \(n\), we have \(z_n = \inf\{a_k : k \geq n\} \leq a_n \leq \sup\{a_k : k \geq n\} = y_n\). This implies:
%                   \[
%                       z_n \leq y_n \quad \text{for all } n.
%                   \]
%                   Taking limits as \(n \to \infty\), we get:
%                   \[
%                       \lim_{n\to\infty} z_n \leq \lim_{n\to\infty} y_n,
%                   \]
%                   which means:
%                   \[
%                       \liminfn{a_{n}} \leq \limsupn{a_{n}}.
%                   \]
%                   For an example where the inequality is strict, consider the sequence \(a_n = (-1)^n\). Then:
%                   \[
%                       y_n = \sup\{(-1)^k : k \geq n\} = 1, \quad z_n = \inf\{(-1)^k : k \geq n\} = -1.
%                   \]
%                   Therefore:
%                   \[
%                       \limsupn{a_{n}} = 1, \quad \liminfn{a_{n}} = -1, \quad \liminfn{a_{n}} < \limsupn{a_{n}}.
%                   \]
%               }

%         \item We will show that \(\liminfn{a_{n}} = \limsupn{a_{n}}\) if and only if \(\lim a_{n}\) exists. In this case, all three share the same value.

%               \iffpf{
%                   Suppose \(\liminfn{a_{n}} = \limsupn{a_{n}} = L\). We will show that \(\lim a_n\) exists and equals \(L\).

%                   Let \(\epsilon > 0\). Since \(\limsupn{a_{n}} = L\), there exists \(N_1\) such that for all \(n \geq N_1\):
%                   \[
%                       y_n = \sup\{a_k : k \geq n\} < L + \epsilon.
%                   \]
%                   Similarly, since \(\liminfn{a_{n}} = L\), there exists \(N_2\) such that for all \(n \geq N_2\):
%                   \[
%                       z_n = \inf\{a_k : k \geq n\} > L - \epsilon.
%                   \]
%                   Let \(N = \max\{N_1, N_2\}\). Then, for all \(n \geq N\):
%                   \[
%                       L - \epsilon < z_n \leq a_n \leq y_n < L + \epsilon,
%                   \]
%                   which implies:
%                   \[
%                       |a_n - L| < \epsilon.
%                   \]
%                   Therefore, \(\lim a_n = L\).
%               }{Conversely, suppose \(\lim a_n = L\). Then, for every \(\epsilon > 0\), there exists \(N\) such that for all \(n \geq N\):
%                   \[
%                       |a_n - L| < \epsilon.
%                   \]
%                   This implies that for all \(n \geq N\), the set \(\{a_k : k \geq n\}\) is contained in \((L - \epsilon, L + \epsilon)\). Therefore:
%                   \[
%                       y_n = \sup\{a_k : k \geq n\} \leq L + \epsilon, \quad z_n = \inf\{a_k : k \geq n\} \geq L - \epsilon.
%                   \]
%                   Taking limits, we get:
%                   \[
%                       \limsupn{a_{n}} \leq L + \epsilon, \quad \liminfn{a_{n}} \geq L - \epsilon.
%                   \]
%                   Since \(\epsilon > 0\) is arbitrary, it follows that \(\limsupn{a_{n}} = \liminfn{a_{n}} = L\). \qedhere}{}{xgray}
%     \end{enumerate}
% }

% \begin{exercise}
%     {2.4.10 (Infinite Products)}

%     A close relative of infinite series is the infinite product
%     \[
%         \prod_{n=1}^{\infty} b_n = b_1 b_2 b_3 \cdots,
%     \]
%     which is understood in terms of its sequence of partial products
%     \[
%         p_m = \prod_{n=1}^{m} b_n = b_1 b_2 b_3 \cdots b_m.
%     \]

%     Consider the special class of infinite products of the form
%     \[
%         \prod_{n=1}^{\infty} (1 + a_n) = (1 + a_1)(1 + a_2)(1 + a_3) \cdots, \quad \text{where} \quad a_n \geq 0.
%     \]

%     \begin{enumerate}
%         \item Find an explicit formula for the sequence of partial products in the case where \( a_n = \frac{1}{n} \) and decide whether the sequence converges. Write out the first few terms in the sequence of partial products in the case where \( a_n = \frac{1}{n^2} \) and make a conjecture about the convergence of this sequence.
%         \item Show, in general, the sequence of partial products converges if and only if \(\sum_{n=1}^{\infty} a_n \text{ converges.}\)
%               (The inequality \(1 + x \leq 3^{x} \text{ for positive } x\) will be useful in one direction.)
%     \end{enumerate}
% \end{exercise}

% \sol{
%     \begin{enumerate}
%         \item For \( a_n = \dfrac{1}{n} \):

%               The sequence of partial products is:
%               \[
%                   p_m = \prod_{n=1}^m \left(1 + \dfrac{1}{n}\right) = \prod_{n=1}^m \dfrac{n + 1}{n}
%               \]
%               This telescopes:
%               \[
%                   p_m = \dfrac{2}{1} \times \dfrac{3}{2} \times \dfrac{4}{3} \times \cdots \times \dfrac{m + 1}{m} = \dfrac{m + 1}{1} = m + 1
%               \]
%               Therefore, the sequence diverges as \( m \to \infty \).

%               For \( a_n = \dfrac{1}{n^2} \):

%               Compute the first few terms:
%               \begin{align*}
%                   p_1 & = 1 + \dfrac{1}{1^2} = 2                                                                                \\
%                   p_2 & = \left(1 + \dfrac{1}{1^2}\right)\left(1 + \dfrac{1}{2^2}\right) = 2 \times \dfrac{5}{4} = \dfrac{5}{2} \\
%                   p_3 & = p_2 \times \left(1 + \dfrac{1}{3^2}\right) = \dfrac{5}{2} \times \dfrac{10}{9} = \dfrac{25}{9}        \\
%                   p_4 & = p_3 \times \left(1 + \dfrac{1}{4^2}\right) = \dfrac{25}{9} \times \dfrac{17}{16} = \dfrac{425}{144}
%               \end{align*}
%               The sequence increases slowly, suggesting that the infinite product is monotone increasing, and thus it converges.

%         \item We will provide an if and only if proof below.

%               \iffpf[We will show that the infinite product \(\prod_{n=1}^\infty (1 + a_n)\) converges if and only if the series \(\sum_{n=1}^\infty a_n\) converges.]{If \(\sum_{n=1}^\infty a_n\) converges, then \(a_n \to 0\). For \(a_n \geq 0\), we have \(\ln(1 + a_n) \leq a_n\). Thus,
%                   \[
%                       \sum_{n=1}^\infty \ln(1 + a_n) \leq \sum_{n=1}^\infty a_n < \infty
%                   \]
%                   So the series \(\sum_{n=1}^\infty \ln(1 + a_n)\) converges, which implies that the product \(\prod_{n=1}^\infty (1 + a_n)\) converges.}{Conversely, if \(\prod_{n=1}^\infty (1 + a_n)\) converges, then the partial products are bounded. For \(a_n \geq 0\) and \(1 + x \geq e^{x/2}\) for small \(x\), we have
%                   \[
%                       \ln(1 + a_n) \geq \dfrac{a_n}{2}
%                   \]
%                   For sufficiently large \(n\), this gives
%                   \[
%                       \sum_{n=1}^\infty a_n \leq 2 \sum_{n=1}^\infty \ln(1 + a_n)
%                   \]
%                   Since \(\sum_{n=1}^\infty \ln(1 + a_n)\) converges, so does \(\sum_{n=1}^\infty a_n\).\qedhere}{}{xgray}
%     \end{enumerate}
% }


% \renewcommand{\theenumi}{\arabic{enumi}}
% \renewcommand{\labelenumi}{\theenumi.}
\section{Subsequences and the Bolzano-Weierstrass Theorem}

\begin{definition}
    Let \(a_n\) be a sequence of real numbers, and let \(n_1 < n_2 < n_3 < \dots\) be an increasing sequence of natural numbers. Then, the sequence \(a_{n_1}, a_{n_2}, a_{n_3}, \dots\) is called a \textit{subsequence} of \(a_n\) and is denoted by \(a_{n_k}\), where \(k \in \N\) indexes the subsequence.
\end{definition}

\begin{theorem}
    Subsequences of a convergent sequence converge to the same limit as the original sequence.
\end{theorem}

\tpf{
    Let \(x_{n_k}\) be a subsequence of \(x_n\), and let \(L = \limn x_n\). We want to show that \(\limn x_{n_k} = L\). Let \(\epsilon > 0\). Since \(x_n\) converges to \(L\), there exists \(N \in \N\) such that for all \(n \geq N\), \(\disabs{x_n - L} < \epsilon\). Since \(n_k\) is increasing, there exists \(M \in \N\) such that \(n_k \geq N\) for all \(k \geq M\). Thus, for all \(k \geq M\), \(\disabs{x_{n_k} - L} < \epsilon\). Therefore, \(\limn x_{n_k} = L\).


    Let \(x_{n_k}\) be a subsequence of \(x_n\). Let \(\epsilon > 0\). Since \((x_n) \rightarrow L\), there exists an \(N \in \N\) such that for all \(n \geq N\), \(\disabs{x_n - L} < \epsilon\).

    Now, looking at \(x_{n_k}\), notice that \(n_k \geq k\) for all \(k\). Consider \(k = N\). For any \(n \geq N\), \(n \geq N \geq k\). Thus, \(\disabs{x_{n_k} - L} < \epsilon\). Therefore, \(\limn x_{n_k} = L\).
}

\begin{ntheorem}
    {Divergence Criterion}If \(x_n\) has two subsequences that converge to different limits, then \(x_n\) diverges.
\end{ntheorem}

Building upon this idea of Divergence, we can list some other ways a sequence can diverge:
\begin{enumerate}
    \item Find one subsequence that diverges.
    \item Find tow subsequences that converge to separate limits.
    \item Negate the \defword{definition of convergence}{2.2.3}.
          \begin{itemize}
              \item For example, a sequence converges to \(L\) if there exists \(\epsilon > 0\) such that for all \(N \in \N\) there exists \(n \geq N\) such that \(\disabs{a_n - a} \geq \epsilon\). There exists a subsequence \((a_{n_k})\) such that for all \(k \in \N\), \(\disabs{a_{n_k} - L} \geq \epsilon\).
                    % implement a tikz visialization of this.
          \end{itemize}
\end{enumerate}

\begin{ntheorem}
    {Bolzano-Weierstrass Theorem}Every bounded sequence in \(\R\) has a convergent subsequence.
\end{ntheorem}

\tpf{
Let \(x_n\) be a bounded sequence. There exists an \(M > 0\) such that every term \(x_n\) belongs to \([-M,M]\). To prove this theorem, we will be utilizing a recursive argument style. Thus, let \(I_0 = [-M,M]\). \(I_0\) has length \(2M\). Cut \(I_0\) in half with \(I_1\) and \(I_2\) both being half as long as \(I_0\). Since \(x_n\) is bounded, there exists an \(I_L\) or \(I_R\) that contains infinitely many terms of \(x_n\). We will pick one, call it \(I_1\) that is contained in \(I_L\). \(I_1\) has length \(M\). Pick one of those terms inside \(I_1\) and call it \(x_{n_1}\). Now, cut \(I_1\) in half with equal length in intervals. One of them contains infinitely many terms. Call that interval \(I_2\). \(I_2\) has length \(\frac{M}{2}\). Pick one of those terms inside \(I_2\) and call it \(x_{n_2}\). Continue this process indefinitely for all \(n \geq \N\) with \(n_1 > n_2\). Continue this process, and we get
\begin{itemize}
    \item a sequence of closed intervals \(I_n\).
          \begin{itemize}
              \item \(I_n\) has length \(\frac{2M}{2^n}\).
              \item They are nested, \(I_n \subseteq I_{n - 1}\).
          \end{itemize}
    \item a subsequence \(x_{n_k}\)
          \begin{itemize}
              \item for all \(k_1\), \(x_{n_k} \in I_k\).
          \end{itemize}
\end{itemize}



The \namrefthm{Nested Interval Property} states that \(\bigcup_{n = 1}^\infty I_n\) is non empty. Let \(L\) be a point in \(\bigcup_{n = 1}^\infty I_n\). We claim \(\limn x_{n_k} = L\). Let \(\epsilon > 0\). There exists an \(N \in \N\) such that \(\frac{2M}{2^n} < \epsilon\). (Since \(\limn \frac{2M{2^n}} = 0\). See \numrefthm{2.5.5}) For any \(k \geq N\), recall that \(x_{n_k}\), \(L \in I_k\). Since \(I_k\) has length \(\frac{2M}{2^n}\). Thus, \(\disabs{x_{n_k} - L} < \epsilon\). Therefore, \(\limn x_{n_k} = L\) and \((x_n)\) has a convergence subsequence.
}

\begin{tikzpicture}

    % Define the interval length
    \def\M{6} % Length of I_0

    % Draw the first interval I_0
    \draw[thick] (-\M, 0) -- (\M, 0) node[anchor=west] {\(I_0 = [-M,M]\)};
    \draw[thick] (-\M, 0.1) -- (-\M, -0.1) node[below] {\(-M\)};
    \draw[thick] (\M, 0.1) -- (\M, -0.1) node[below] {\(M\)};

    \draw[thick] (5, 0.2) node[above] {\(Right\)};
    \draw[thick] (-5, 0.2) node[above] {\(Left\)};

    % Draw the midpoint of I_0
    \draw[thick] (0, 0.1) -- (0, -0.1) node[below] {\(0\)};

    % Draw subintervals I_1, I_L, I_R
    \def\halfM{\M / 2}
    \draw[thick] (0, -1) -- (\halfM, -1) node[anchor=west] {\(I_1\)};
    \draw[thick] (0, -0.9) -- (0, -1.1) node[below] {};
    \draw[thick] (\halfM, -0.9) -- (\halfM, -1.1) node[below] {};

    % Draw I_2 within I_1
    \def\quarterM{\M / 4}
    \draw[thick] (\quarterM, -2) -- (\halfM, -2) node[anchor=west] {\(I_2\)};
    \draw[thick] (\quarterM, -1.9) -- (\quarterM, -2.1) node[below] {};
    \draw[thick] (\halfM, -1.9) -- (\halfM, -2.1) node[below] {};

    % Further subintervals
    \def\eighthM{\M / 8}
    \draw[thick] (\quarterM, -3) -- (\quarterM + \eighthM, -3) node[anchor=west] {\(I_3\)};
    \draw[thick] (\quarterM + \eighthM, -2.9) -- (\quarterM + \eighthM, -3.1) node[below] {};
    \draw[thick] (\quarterM, -2.9) -- (\quarterM, -3.1) node[below] {};

    % Points representing x_n's
    \filldraw[red] (2.5, 0) circle (2pt) node[anchor=north] {\(x_{n_1}\)};
    \filldraw[red] (2.25, -1) circle (2pt) node[anchor=north] {\(x_{n_2}\)};
    \filldraw[red] (2.125, -2) circle (2pt) node[anchor=north] {\(x_{n_3}\)};
    \filldraw[red] (2.0625, -3) circle (2pt) node[anchor=north] {\(x_{n_4}\)};

\end{tikzpicture}

\begin{theorem}
    Let \(b \in (0,1)\). Then \(\limn b^n = 0\).
\end{theorem}

\tpf{
    The sequence \((b^n)\) is monotone decreasing. This is because \(b^{n + 1} = b^nb < b^n\). This sequence is also bounded by 0. Hence, by the \namrefthm{Monotone Convergence Theorem}, \((b^n)\) converges. Now, let \(L = \limn b^n\). Consider the subsequence \(b^{2n}\). This sequence also converges to \(L\). Thus, \begin{align*}
        L & = \limn b^{2n}       \\
          & = \limn b^nb^n       \\
          & = \limn b^n\limn b^n \\
          & = L^2.
    \end{align*} Thus, \(L = 0\) or \(L = 1\). The limit cannot be 1 because \(b^n\) is decreasing away from 1. Therefore, \(L = 0\).
}

% \renewcommand{\theenumi}{\alph{enumi}}
% \renewcommand{\labelenumi}{(\theenumi)}
% \subsection{Exercises}

% \begin{exercise}
%     {2.5.1}Give an example of each of the following, or argue that such a request is impossible.
%     \begin{enumerate}
%         \item A sequence that has a subsequence that is bounded but contains no subsequence that converges.
%         \item A sequence that does not contain 0 or 1 as a term but contains subsequences converging to each of these values.
%     \end{enumerate}
% \end{exercise}

% \sol{
%     \begin{enumerate}
%         \item \textbf{Impossible.} This violates the \namrefthm{Bolzano-Weierstrass Theorem}. It assures us that every bounded sequence has a convergent subsequence. If a subsequence is bounded, then it must have a convergent subsequence.
%         \item Consider the sequence \((\frac{1}{2},\frac{1}{2},\frac{1}{3},\frac{2}{3},\frac{1}{4},\frac{3}{4}, \cdots \frac{1}{n}, \frac{(n-1)}{n})\). From this, you can have a subsequence \((\frac{1}{2}, \frac{1}{3}, \dots, \frac{1}{n})\) which converges to 0, and also a subsequence \((\frac{1}{2}, \frac{2}{3}, \dots \frac{n - 1}{n})\), which converges to 1.
%     \end{enumerate}
% }

% \begin{exercise}
%     {2.5.2} Decide whether the following propositions are true or false, providing a short justification for each conclusion.
%     \begin{enumerate}
%         \item If every proper subsequence of \((x_n)\) converges, then \((x_n)\) converges as well.
%         \item If \((x_{n})\) contains a divergent subsequence, then \((x_n)\) diverges.
%         \item If \((x_n)\) is bounded and diverges, then there exist two subsequences of \((x_n)\) that converge to different limits.
%     \end{enumerate}
% \end{exercise}

% \sol{
%     \begin{enumerate}
%         \item \textbf{True.} If every proper subsequence of \((x_n)\) converges, then \((x_n)\) must converge to the same limit. If \((x_n)\) did not converge, there would exist at least one divergent subsequence or two subsequences converging to different limits, contradicting the assumption.

%         \item \textbf{True.} If \((x_n)\) contained a divergent subsequence, then \((x_n)\) cannot converge. A convergent sequence has all its subsequences converging to the same limit, so the existence of a divergent subsequence implies that \((x_n)\) diverges (contrapositive).

%         \item \textbf{True.} Since \((x_n)\) is bounded and diverges, the \namrefthm{Bolzano-Weierstrass Theorem} guarantees the existence of at least one convergent subsequence. Let this subsequence converge to \(L_1\). Because \((x_n)\) does not converge to \(L_1\), there is an \(\epsilon > 0\) and infinitely many terms of \((x_n)\) such that \(|x_n - L_1| \geq \epsilon\). Extracting a subsequence from these terms, the Bolzano-Weierstrass Theorem ensures a further subsequence converging to a limit \(L_2 \neq L_1\). Thus, \((x_n)\) has two subsequences converging to different limits.
%     \end{enumerate}
% }


% \begin{exercise}
%     {2.5.5} Assume \((a_n)\) is a bounded sequence with the property that every convergent subsequence of \((a_n)\) converges to the same limit \(a \in \mathbb{R}\). Show that \((a_n)\) must converge to \(a\).
% \end{exercise}

% \expf{
% Suppose that \((a_n)\) does not converge to \(a \in \R\). By the definition of convergence, this means there is a positive real number \(\epsilon_0\) such that no matter how large we choose \(N \in \N\), there will always exist some \(n > N\) where \(|a_n - a| \geq \epsilon_0\). In a formal way, this shows that \((a_n)\) does not converge to \(a\) within the \(\epsilon_0\)-neighborhood.

% We aim to demonstrate that this leads to a contradiction by constructing a subsequence of \((a_n)\) that stays outside this neighborhood. Begin by selecting \(n_1\) such that \(|a_{n_1} - a| \geq \epsilon_0\). Next, since the condition holds for all \(N \in \N\), we can find another index \(n_2 > n_1\) such that \(|a_{n_2} - a| \geq \epsilon_0\). Continuing this process, we generate an increasing sequence of indices \(n_1 < n_2 < n_3 < \dots\) such that for each \(i \in \N\), \(|a_{n_i} - a| \geq \epsilon_0\).

% Now consider the subsequence \((a_{n_i})\) we have built. Since \((a_n)\) is bounded by assumption, its subsequence \((a_{n_i})\) is also bounded. By the \namrefthm{Bolzano-Weierstrass Theorem}, every bounded sequence has a convergent subsequence. Let \((a_{n_{i_k}})\) denote a convergent subsequence of \((a_{n_i})\). According to our assumption, any convergent subsequence of \((a_n)\) must converge to \(a\).

% However, each term of \((a_{n_{i_k}})\) remains outside the \(\epsilon_0\)-neighborhood of \(a\). Thus, it is impossible for \((a_{n_{i_k}})\) to converge to \(a\). This contradiction implies that our initial assumption—that \((a_n)\) does not converge to \(a\)—is false. Therefore, the sequence \((a_n)\) must converge to \(a\).
% }



% \begin{exercise}
%     {2.5.6} Use a similar strategy to the one in \numrefthm{2.5.5} to show
%     \[\lim b^{1/n} \text{ exists for all } b \geq 0\] and find the value of the limit. (The results in \hyperref[exerc:2.3.1]{Exercise 2.3.1} may be assumed.)
% \end{exercise}

% \expf{
%     We will show that \(\lim_{n \to \infty} b^{1/n}\) exists for all \(b \geq 0\) and find its value.
%     \begin{itemize}
%         \item \textbf{Case 1:} \(b = 0\).

%               When \(b = 0\), the sequence becomes \(a_n = 0^{1/n} = 0\) for all \(n\). Thus, \(\lim_{n \to \infty} b^{1/n} = 0\).

%         \item \textbf{Case 2:} \(b > 0\).

%               Suppose, for contradiction, that \(\lim_{n \to \infty} b^{1/n} \ne 1\). Then there exists \(\epsilon > 0\) and infinitely many \(n\) such that \(|b^{1/n} - 1| \geq \epsilon\). Extract a subsequence \((b^{1/n_k})\) where this inequality holds for all \(k\).

%               Since \(b^{1/n} > 0\) and bounded, by the \namrefthm{Bolzano-Weierstrass Theorem}, the subsequence \((b^{1/n_k})\) has a further subsequence that converges to a limit \(L\). According to \hyperref[ex:2.3.1]{Exercise 2.3.1}, any convergent subsequence of \((b^{1/n})\) must have its limit equal to \(\lim_{n \to \infty} b^{1/n}\).

%               Consider \(\ln b^{1/n} = \frac{\ln b}{n}\). As \(n \to \infty\), \(\frac{\ln b}{n} \to 0\), so \(\ln b^{1/n} \to 0\), which implies \(b^{1/n} \to e^{0} = 1\).

%               This contradicts the assumption that \(|b^{1/n_k} - 1| \geq \epsilon\), so \(\lim_{n \to \infty} b^{1/n} = 1\).

%     \end{itemize}

%     \textbf{Conclusion:}

%     \[
%         \lim_{n \to \infty} b^{1/n} =
%         \begin{cases}
%             0, & \text{if } b = 0,                \\
%             1, & \text{if } b > 0. \hfill\qedhere
%         \end{cases}
%     \]
% }


% \renewcommand{\theenumi}{\arabic{enumi}}
% \renewcommand{\labelenumi}{\theenumi.}
\section{The Cauchy Criterion}

\begin{center}
    \textbf{Recall}
\end{center}
How do we prove \(x_n\) converges?
\begin{enumerate}
    \item We know and prove the limit \(\rightarrow\) claim \(L\), show terms get close to \(L\).
    \item \namrefthm{Monotone Convergence Theorem}.
\end{enumerate}

\begin{definition}
    A sequence \((x_n)\) is a \textit{Cauchy sequence} if for all \(\epsilon > 0\), there exists \(N \in \N\) such that for all \(m,n \geq N\), \(\disabs{x_m - x_n} < \epsilon\).
\end{definition}

This says that as terms get close to each other and stay close together, there's some value they're all getting close to.

\begin{center}
    \textbf{Geometric Series Review}
\end{center}

Remember that geometric series consist of terms that are multiplied by a ratio \(r\). For example, that could look like \(1 + r + r^{2} + r^{3} + \cdots\).

We are most interested in \textbf{partial sums}. That is,
\[
    1 + r + r^{2} + \cdots r^{n - 1} + r^{n} = S_{n}.
\]
From here, we we would multiply both sides by \(r\). This gives
\[
    r + r^{2} + \cdots r^{n} + r^{n + 1} = rS_{n}.
\]
When we subtract these two from each other, we get
\[
    1 - r^{n + 1} = S_{n} - rS_{n}.
\]
This yields the identity
\[
    1 + r + r^{2} + \cdots r^{n} = \frac{1-r^{n + 1}}{1 - r}.
\]

\begin{example}
    {Cauchy Sequence}Consider the sequence \(a_1 = 1, a_2 = 2,\) where \[a_n = \frac{a_{n-1} + a_{n-2}}{2} \text{ for all } n \geq 2.\] Show this sequence is Cauchy.
\end{example}

\lepf{
Look at the differences of consequtive terms, \(\disabs{a_1 - a_2} = 1\), \(\disabs{a_2 - a_3} = 1/2\), we can see a formula \(a_n - a_{n+1} = {1/2}^{n-1}\). Assume \(\disabs{a_n - a_m}= \disabs{a_n - a_{n+1} - a_{n+2}} - \cdots - a_{m - 1} - a_m\) with \(n <m\). From the \hyperref[def:1.2.4]{Triangle Inequality},
\begin{align}
    \disabs{a_n - a_m} & \leq \disabs{a_n - a_{n+1}} + \disabs{a_{n+1} - a_{n+2}} + \cdots + \disabs{a_{m-1} - a_m} \\
                       & = \frac{1}{2^{n-1}} + \frac{1}{2}^{n} + \cdots + \frac{1}{2}^{m-2}                         \\
                       & = \frac{1}{2^{n-1}} \left(1 + \frac{1}{2} + \cdots + \frac{1}{2}^{m-n-1}\right)            \\
                       & = \frac{1}{2^{n-1}} \left(\frac{1 - \frac{1}{2}^{m-n}}{1 - \frac{1}{2}}\right)             \\
                       & = \frac{1}{2^n} \left(1 - \frac{1}{2^{m-n}}\right)                                         \\
                       & < \frac{1}{2^{n}}.
\end{align}
Notice that we were able to pull out the \(1/2\) and use the geometric series formula at step 2.12. From here we know that \(\disabs{a_n - a_m} < \frac{1}{2^n}\).

Now, conclude the proof by letting \(\epsilon > 0\). We know \((1/2^n) \rightarrow 0\). Thus, there exists \(N \in \N\) such that \(\frac{1}{2^N} < \epsilon\). For all \(n,m \geq N\), (without loss of generality \(n < m\)) \(\disabs{a_n - a_m} < \frac{1}{2^{n}} \leq \frac{1}{2^{N}}< \epsilon\). Therefore, \(a_n\) is \hyperref[def:2.6.1]{Cauchy} and it converges.
}

\begin{note}
    \textbf{Note:} To find the limit of this series, a proof strategy is finding subsequences that are odd and even, and show the converge to the same limit.
\end{note}

\begin{ntheorem}
    {Cauchy Criterion}A sequence \(x_n\) converges if, and only if, it is a Cauchy sequence.
\end{ntheorem}

\iffpf{Assume \((x_n)\) is a convergent seqeuence in \(\R\). Given \(\epsilon > 0\). Let \(L = \limn x_n\). Since \((x_n) \rightarrow L\), there exists an \(N \in \N\) such that for all \(n \geq N\), \(\disabs{x_n - L} < \frac{\epsilon}{2}\). For all \(n,m \geq N\), \begin{align*}
        \disabs{x_m - x_n} & = \disabs{x_m - L + L - x_n}              \\
                           & \leq \disabs{x_m - L} + \disabs{L - x_n}  \\
                           & < \frac{\epsilon}{2} + \frac{\epsilon}{2} \\
                           & = \epsilon.
    \end{align*}
    Therefore, \(x_n\) is a Cauchy sequence.}{Assume \(x_n\) is a Cauchy sequence.
    \begin{itemize}
        \item \textbf{Step 1:} Show that \(x_n\) is bounded.

              Since \(x_n\) is Cauchy, there exists \(N \in \N\) such that for all \(n,m \geq N\), \(\disabs{x_n - x_m} < l\). It follows that for all \(n \geq N\), we need to account for \(x_1\dots,x_{n-1}\). Thus, let \(M = \max\{\disabs{x_1},\disabs{x_2},\dots,\disabs{x_{n-1}},\disabs{x_n}+1\}\). Then for all \(n \in \N\), \(\disabs{x_m} < M\).
        \item \textbf{Step 2:} Since \(x_n\) is bounded, there exits a convergent subsequence \(x_{n_k}\) by the \namrefthm{Bolzano-Weierstrass Theorem}. Let \(L\) be the limit of the subsequence.
        \item \textbf{Step 3:} Show that \(x_n\) converges to \(L\).

              If some get close to \(L\) and all get close to each other, they all get close to \(L\). Let \(\epsilon > 0\). Since \(x_{n_k}\) converges to \(L\), there exists \(N \in \N\) such that for all \(k \geq N\), \(\disabs{x_{n_k} - L} < \frac{\epsilon}{2}\). Since \(x_n\) is Cauchy, there exists \(M \in \N\) such that for all \(n,m \geq M\), \(\disabs{x_n - x_m} < \frac{\epsilon}{2}\). Let \(M_0 = \max \{N_1, n_k\}\). By the \namrefthm{Archimedean Principle}, there exists \(N_0\) such that \(n_{k_0} \geq M_0\). Then, from the \hyperlink{Triangle Inequality}{Triangle Inequality}, we say that for all \(n \geq N_0\), \begin{align*}
                  \disabs{x_n - L} & \leq \disabs{x_n - x_{n_{k_0}}} + \disabs{x_{n_{k_0}} - L} \\
                                   & < \frac{\epsilon}{2} + \frac{\epsilon}{2}                  \\
                                   & = \epsilon.
              \end{align*} Therefore, \((x_n) \rightarrow L\).
    \end{itemize}}{By proving both directions of the inequality, we found that a sequence \((x_n)\) converges if, and only if, it is a Cauchy sequence.}%
{xblue}%

\begin{definition}
    A sequence is called \textit{contracting} is there exists \(0 < C < 1\) such that for all \(n \in \N\), \(\disabs{x_{n+1} - x_n} \leq C\disabs{x_n - x_{n-1}}\).
\end{definition}

How this works: we take a sequence \(a_1, a_2, \dots\) and subtract \(a_1 - a_2\). Then, we have the inequality: \begin{align*}
    \disabs{a_2 - a_1} & \leq C\disabs{a_1 - a_0}                            \\
    \disabs{a_3 - a_2} & \leq C\disabs{a_2 - a_1} \leq C^2\disabs{a_1 - a_2} \\
    \disabs{a_4 - a_3} & \leq C\disabs{a_3 - a_2} \leq C^3\disabs{a_1 - a_2} \\
                       & \quad \vdots
\end{align*} From this, a theorem emerges:

\begin{theorem}
    If a sequence is contracting, then it is Cauchy, and thus converges.
\end{theorem}

\tpf{
Let \((a_n)\) be a contracting sequence; that is, there exists a constant \(0 < C < 1\) such that for all \(n \in \mathbb{N}\),
\[
    |a_{n+1} - a_n| \leq C |a_n - a_{n-1}|.
\]
We will show that \((a_n)\) is a Cauchy sequence.

First, we observe by induction that for all \(k \geq 1\),
\[
    |a_{n+k} - a_{n+k-1}| \leq C^{k} |a_n - a_{n-1}|.
\]
\textit{Proof by induction}:

\textbf{Base case \((k=1)\):}
\[
    |a_{n+1} - a_{n}| \leq C |a_n - a_{n-1}|.
\]

\textbf{Inductive step:} Assume that for some \(k \geq 1\),
\[
    |a_{n+k} - a_{n+k-1}| \leq C^{k} |a_n - a_{n-1}|.
\]
Then,
\begin{align*}
    |a_{n+k+1} - a_{n+k}| & \leq C |a_{n+k} - a_{n+k-1}|                \\
                          & \leq C \left( C^{k} |a_n - a_{n-1}| \right) \\
                          & = C^{k+1} |a_n - a_{n-1}|.
\end{align*}
Thus, the inequality holds for \(k+1\), completing the induction.

Next, for any integers \(m > n\), we have:
\[
    |a_m - a_n| = \left| \sum_{j=n}^{m-1} \left( a_{j+1} - a_j \right) \right| \leq \sum_{j=n}^{m-1} |a_{j+1} - a_j|.
\]
Applying the inequality obtained from the induction,
\[
    |a_{j+1} - a_j| \leq C^{j - n +1} |a_n - a_{n-1}|.
\]
Therefore,
% \begin{align*}
%     |a_m - a_n| & \leq |a_n - a_{n-1}| \sum_{j=n}^{m-1} C^{j - n +1}                                 \\
%                 & = |a_n - a_{n-1}| \sum_{k=1}^{m - n} C^{k} \quad (\text{Let } k = j - n +1 \text)} \\
%                 & = |a_n - a_{n-1}| \left( \frac{C (1 - C^{m - n})}{1 - C} \right)
% \end{align*}
Since \(C^{m - n} \geq 0\), we have:
\[
    |a_m - a_n| \leq |a_n - a_{n-1}| \left( \frac{C}{1 - C} \right).
\]
As \(n \to \infty\), the term \(|a_n - a_{n-1}|\) tends to zero because:
\[
    |a_n - a_{n-1}| \leq C^{n-1} |a_1 - a_0| \to 0 \quad \text{as } n \to \infty.
\]
Therefore, for any \(\epsilon > 0\), there exists \(N \in \mathbb{N}\) such that for all \(n \geq N\),
\[
    |a_n - a_{n-1}| < \epsilon \left( \frac{1 - C}{C} \right).
\]
Then, for all \(m, n \geq N\) (with \(m > n\)),
\[
    |a_m - a_n| \leq |a_n - a_{n-1}| \left( \frac{C}{1 - C} \right) < \epsilon.
\]
This shows that \((a_n)\) is a Cauchy sequence. Since every Cauchy sequence in \(\mathbb{R}\) converges, the sequence \((a_n)\) converges.
}




\chapter{Basic Topology of Real Numbers}\label{topology}
\vspace*{-0.25in}
\renewcommand{\theenumi}{\arabic{enumi}}
\renewcommand{\labelenumi}{\theenumi.}
\section{Discussion: The Cantor Set}

We will build this set through an iterative process. Start with a number line \(C_0\) that stretches from 0 to 1. Remove the middle third of the interval, leaving two intervals of length \(\frac{1}{3}\). We will call the set of points removed from \(C_0\) \(C_1\). Next, remove the middle third of each of the two intervals, leaving four intervals of length \(\frac{1}{9}\). We will call the set of points removed from \(C_1\) \(C_2\). Continue this process indefinitely. 

% Put a tikz picture here that shows the cantor set at each step

\begin{definition}
    The \textit{Cantor set}, \(C\), is defined as \(C = \bigcap_{n=0}^\infty C_n\). This set is 
    \begin{enumerate}
        \item non-empty. All end points stay within the interval.
        \item uncountable.
    \end{enumerate}
\end{definition}

The second part of that definition is a bit tricky to prove, but a visual will do for now. We can put all elements of the Cantor set in a one-to-one correspondence with the set of all 0s and 1s. This shows that not only is it uncountable, but it also has the same cardinality as \([0,1]\). % Put a tikz picture here from the notes.

The total length of removed elements, \(\frac{1}{3} + \frac{2}{9} + \frac{4}{27} + \frac{8}{81} + \cdots = \frac{1}{3}(1 + \frac{2}{3} + \frac{4}{9})\). Notice the resemblance to the geometric series? We can write this as \[\frac{1}{3} \sum_{n=0}^\infty \left(\frac{2}{3}\right)^n = \frac{1}{3} \left(\frac{1}{1 - \frac{2}{3}}\right) = 1.\] 

In summary, \begin{enumerate}
    \item Start with intercal \([0,1]\).
    \item Remove countably disjoint intervals.
    \item Uncountably many points between these intervals, all isolated from each other.
    \item The space taken up by the leftover points has ``length'' 0.
\end{enumerate}

When we review the properties of a fractal, we see that the Cantor set is a fractal. It is self-similar, and the dimension of the Cantor set is \(\log_3 2\).

\begin{note}
    For a cool look at the cantor set as a fractal, check out \url{https://en.wikipedia.org/wiki/File:Cantor_Set_Expansion.gif}.
\end{note}

\renewcommand{\theenumi}{\arabic{enumi}}
\renewcommand{\labelenumi}{\theenumi.}
\section{Open and closed Sets}

\begin{definition}
    For a point \(x \in \R\), and \(\epsilon > 0\), we define the \textit{epsilon-neighborhood} of \(x\) to be \(V_\epsilon(x) = \{y \in \R \mid |x - y| < \epsilon\}\).
\end{definition}

In other words, \(V_{\epsilon}(x)\) is the open interval \((x - \epsilon, x + \epsilon)\), centered at \(x\) with radius \(\epsilon\).

\subsection{Open Sets}

\begin{definition}
    A set \(A \subseteq \R\), is called an \textit{open set} if for every \(x \in A\), there exists an \(\epsilon > 0\) such that \(V_\epsilon(x) \subseteq A\).
\end{definition}

\subsubsection{Some Examples of Open Sets}

\begin{multicols}{2}
    \begin{itemize}
        \item All open intervals are also open sets.
        \item \(\R\) is open.
        \item \(\emptyset\) is open.
        \item \(\{1\}\) is not open.
        \item \([0,2]\) is not open.
        \item \(\Q\) is not open.
        \item \([4,6)\) is not open.
        \item \((0,1) \cup (1,3) \cup (5,10)\) is open.
        \item \((0,3] \cap [2,4)\) is open.
        \item Cantor set is not open.
    \end{itemize}
\end{multicols}

\begin{theorem}
    \begin{enumerate}[label=(\roman*)]
        \item The union of an arbitrary collection of open sets is open.
        \item The intersection of a finite collection of open sets is open. 
    \end{enumerate}
\end{theorem}

\tpf{
    (i) Let \(\{O_{\lambda} \colon \lambda \in A\}\) be a collection of open sets. Then, let \(O = \bigcup_{\lambda \in A} O_{\lambda}\). Let \(a\) be an element of \(O\). To show that \(O\) is open, we need to find an \(\epsilon\)-neighborhood that is completely contained within \(O\) to satisfy \defref{3.2.1}. But \(a \in O\) implies that \(a\) is an element of at least one particular \(O_{\lambda'}\). Because we are assuming \(O_{\lambda'}\) to be open, then we can use \defref{3.2.1} to assert that there exists \(V_{\epsilon}(a) \subseteq O_{\lambda'}\). The fact that \(O_{\lambda'} \subseteq O\) confirms that \(V_{\epsilon}(a) \subseteq O\). \\

    \noindent (ii) Let \(\{O_{1},O_{2},\dots,O_{N}\}\) be a finite collection of open sets. Then, let \(a \in \bigcap_{k = 1}^{N} O_{k}\). This means \(a\) is an element of every open set. \defref{3.2.1} tells us that for \(1 \leq k \leq K\), there exists an \(V_{\epsilon}(a) \subseteq O_{k}\). From this set, we are in search of one \(\epsilon\)-neighborhood of \(a\) that is contained in every \(O_{k}\), so the trick is to pick the smallest one. Letting \(\epsilon = \min\{\epsilon_{1},\epsilon_{2},\dots,\epsilon_{N}\}\), it follows that \(V_{\epsilon}(a) \subseteq V_{\epsilon_{k}} (a)\) for all \(k\), and hence \(V_{\epsilon}(a) \subseteq \bigcap_{k = 1}^{N} O_{k}\). 
}

Note that we cannot use this for cases with infinity. For example, consider \(A_n = (-\frac{1}{n}, \frac{1}{n})\). This is open, but \(\bigcap_{n = 1}^\infty\ A_n = \{0\}\), which is not open. 

\subsection{Closed Sets}

\begin{definition}
    Let \(A \subseteq \R\). We say \(x\) is a \textit{limit point} of \(A\) if for all \(\epsilon > 0\), there exists \(a \in A\) such that \(a \in V_\epsilon(x)\) that is not \(x\). Additionally, a point \(x \in \R\) is a \textit{limit point} if, and only if, there exists a sequence \((a_n)\) of points from \(A\) that are not \(x\). And \(\limn(a_n) = x\). 
\end{definition}

\begin{definition}
    A set \(B \subseteq \R\) is called a \textit{closed set} if \(B\) contains all its limit points. 
\end{definition}

\begin{note}
    \textbf{Important note:} Limit points could be outside a set. Consider \((0,1)\). Even though 0 and 1 do not belong to the set, they are considered limit points that are outside the set.
\end{note}

\subsubsection{Some Examples of Closed Sets}
\begin{multicols}{2}
    \begin{itemize}
        \item \([0,1]\) is closed.
        \item \((0,1)\) is not closed.
        \item \(\R\) is closed.
        \item \(\emptyset\) is closed.
        \item \(\Q\) is not closed.
        \item \([3, \infty)\) is closed.
        \item \(\frac{1}{n} \mid n \in \N\) not closed. (Because of 0)
        \item \([1,4] \cup \{8\}\) is closed. 
        \item \(\{1\}\) is closed. 
        \item \([1,2)\) is not closed. Note that this set is neither open or closed. 
    \end{itemize}
\end{multicols}

\begin{theorem}
    A set \(B \subseteq \R\) is closed if, and only if, its complement is open. Similarly, a set \(A \subseteq \R\) is open if, and only if, its complement is closed.
\end{theorem}

\iffpf{Assume \(B \subseteq \R\) is a closed set. We will show that \(B^c\) is open. Let \(x \in B^c\). So, \(x \notin B\). This means \(x\) is not a limit point. (From the negated definition of limit point:) There must exist \(\epsilon >0\) such that no elements of \(B\) belong to \(V_\epsilon(x)\). Then, \(V_\epsilon(x) \subseteq B^c\). Therefore, \(B^c\) is open.}{Assume \(B^c\) is open. We will show that \(B\) is closed. Let \(x\) be a limit point of \(B\). For all \(\epsilon > 0\), there exists a \(b \in B\) such that \(b \subseteq V_\epsilon(x)\). So, \(V_\epsilon(x)\) is not a subset of \(B^c\). This is true for every \(\epsilon\). Since \(B^c\) is open, it must be that \(x \notin B^c\). Thus, \(x \in B\). So, \(B\) contains all its limit points. Therefore, \(B\) is closed.\qedhere}{}{xblue}

\begin{definition}
    Let \(A \subseteq \R\) and let \(L\) be the set of limit points of \(A\). The \textit{closure} of \(A\) is defined as \(\bar{A} = A \cup L\).
\end{definition}

\begin{theorem}
    \begin{enumerate}[label=(\roman*)]
        \item The intersection of any collection of closed sets is closed. 
        \item The union of finitely many closed sets is closed.
    \end{enumerate}
\end{theorem}

\tpf{
    De Morgan's Laws state that for any collection of sets \(\{E_{\lambda} \colon \lambda \in A\}\) it is true that 
    \[
    \left(\bigcup_{\lambda \in A} E_{\lambda}\right)^{c} = \bigcap_{\lambda \in A} E_{\lambda}^{c} \qquad \text{and} \qquad \left(\bigcap_{\lambda \in A} E_{\lambda}\right)^{c} = \bigcup_{\lambda \in A} E_{\lambda}^{c}. 
    \]
    The result follows directly from these statements and \numrefthm{3.2.3}.
}

\begin{theorem}
    The closure of a set is a closed set.
\end{theorem}

\textbf{Note:} This theorem may seem trivial, but it answers the question of ``Are there limit points in \(L\) that are not accounted for?''

\tpf{
    We need to show that \(\bar{A}\) contains all the limit points of \(\bar{A}\). Let \(L\) be the limit points of \(A\). Thus, \(\bar{A} = A \cup L\). Let \(x\) be a limit point of \(\bar{A}\). There exists a sequence of points \((x_n)\) coming from \(\bar{A}\) such that \((x_n) \rightarrow x\). Then, for all \(n \in \N\), either \(x_n \in A\) or \(x_n \in L\).
    \begin{itemize}
        \item \textbf{Case 1:} \(x_n \in A\) \\
        There exists a subsequence \((x_{n_k})\) where each \(x_{n_k} \in A\). This subsequence also converges to \(x\), and we know the limit belongs to \(L\), so \(x \in L \subseteq \bar{A}\).

        \item \textbf{Case 2:} \(x_n \in L\) \\
        \(x_n\) belongs to \(A\) for only finitely many \(n \in \N\). Thus, a tail-end of the sequence is comprised entirely of points from \(L\). To simplify things, we will assume the entire sequence \((x_n)\) comes from \(L\). (We know that \((x_n)\) converges to \(x\), but we cannot assume those limit points converge as well.) Let \(n \in \N\). Since \(x_n \in L\), there exists \(a_n \in A\) such that \(\disabs{x_n - a_n} < \frac{1}{n}\). We now have \((x_n) \rightarrow x\) and \((x_n - a_n) \rightarrow 0\). Then, \((a_n) \rightarrow x\). Thus, \(x \in L \subseteq \bar{A}\).
    \end{itemize}
    Now that we have shown that either cases leads to the same conclusion, we know that \(x \in L \subseteq \bar{A}\), and therefore \(\bar{A}\) contains all its limit points.
}

\begin{theorem}
    The closure set \(\overline{A}\) is the \textit{smallest} closed set containing \(A\). (Where ``smallest'' refers to a subset of any other closed set containing \(A\).)
\end{theorem}

\tpf{
    If \(B\) is a closed set containing \(A\), then \(A \subseteq B\) and \(L \subseteq B\). Thus, \(\bar{A} = A \cup L \subseteq B\).
}

\begin{example}
    {Closed Sets 1}Generate countably many closed sets where the union is not closed.
\end{example}

\lesol{
    \(B_n = [\frac{1}{n}, 1 - \frac{1}{n}]\). Therefore, \(\bigcup_{n = 3}^\infty B_n = (0,1)\). For example, that would look like: \(\{\frac{1}{2}\} \cup \{\frac{1}{3}\} \cup \dots\)
}

\begin{example}
    {Closed Sets 2}What is the closure of the following sets? 
    \[
    (\text{a}) \ (0,1), \quad \text{(b)} \ \R, \quad (\text{c}) \ \left\{\frac{1}{n} \mid n \in \N\right\}, \quad (\text{d}) \ [0,1) \cup (1,3], \quad (\text{e}) \ \Q
    \]
\end{example}

\lesol{
    (a) \(\bar{A} = [0,1]\), \quad (b) \(\bar{A} = \R\), \quad (c) \( \bar{A} = A \cup \{0\}\), \quad (d) \(\bar{A} = [0,3]\), \quad (e) \(\bar{A} = \R\).
}

\renewcommand{\theenumi}{\alph{enumi}}
\renewcommand{\labelenumi}{(\theenumi)}
\subsection{Exercises}

\begin{exercise}
    {3.2.4}Let \(A\) be a nonempty and bounded above set so that \(s = \sup(A)\) exists. (See \defref{1.3.2} and \defref{3.2.7})
    \begin{enumerate}
        \item Show that \(s \in \bar{A}\).
        \item Can an open set contain its supremum?
    \end{enumerate}
\end{exercise}

\sol{
    \begin{enumerate}
        \item We need to show that \(s = \sup(A) \in \bar{A}\), where \(\bar{A} = A \cup L\), and \(L\) is the set of limit points of \(A\).

        Since \(A\) is nonempty and bounded above, \(s = \sup(A)\) exists.

        If \(s \in A\), then \(s \in \bar{A}\) trivially.

        Suppose \(s \notin A\). We will show that \(s\) is a limit point of \(A\), so \(s \in L \subseteq \bar{A}\).

        By definition, \(x\) is a limit point of \(A\) if for all \(\epsilon > 0\), there exists \(a \in A\) such that \(a \in V_\epsilon(x)\) and \(a \ne x\).

        Fix any \(\epsilon > 0\). Since \(s = \sup(A)\), for this \(\epsilon\), \(s - \epsilon\) is not an upper bound of \(A\). Therefore, there exists \(a \in A\) such that
        \[
        s - \epsilon < a \leq s.
        \]
        Since \(a \leq s\) and \(a > s - \epsilon\), we have \(|a - s| < \epsilon\), so \(a \in V_\epsilon(s)\) and \(a \ne s\).

        Therefore, \(s\) is a limit point of \(A\), and hence \(s \in \bar{A}\).

        \item An open set cannot contain its supremum if the supremum is finite.

        Assume \(A\) is an open set containing its supremum \(s\).

        Since \(A\) is open and \(s \in A\), there exists \(\epsilon > 0\) such that
        \[
        V_\epsilon(s) = \{ x \in \mathbb{R} \mid |x - s| < \epsilon \} \subseteq A.
        \]
        This means \(s + \frac{\epsilon}{2} \in A\).

        However, \(s\) is an upper bound of \(A\), so no element of \(A\) can be greater than \(s\).

        This is a contradiction.

        Therefore, an open set cannot contain its supremum.
    \end{enumerate}
}

\begin{exercise}
    {3.2.6}Decide whether the following statements are true or false. Provide counterexamples for those that are false, and supply proofs for those that are true.
    \begin{enumerate}
        \item An \hyperref[def:3.2.2]{open set} that contains every rational number must necessarily be all of \(\R\).
        \item The \namrefthm{Nested Interval Property} remains true if the term ``closed interval'' is replaced by ``\hyperref[def:3.2.5]{closed set}.''
        \item Every nonempty open set contains a rational number.
        \item Every bounded infinite closed set contains a rational number.
        \item The \hyperref[def:3.1.1]{Cantor set} is closed.
    \end{enumerate}
\end{exercise}

\sol{
    \begin{enumerate}
        \item \textbf{False.}

        \textit{Counterexample:} Consider the set \(U = \bigcup_{n=1}^\infty \left( q_n - \frac{1}{n},\ q_n + \frac{1}{n} \right)\), where \((q_n)\) is an enumeration of all rational numbers.

        Each interval \(\left( q_n - \frac{1}{n},\ q_n + \frac{1}{n} \right)\) is open, and the union \(U\) is open. Since every rational number is included in some interval, \(U\) contains all rationals. However, \(U \ne \mathbb{R}\) because there are irrational numbers not covered by these intervals. Therefore, an open set can contain all rational numbers without being all of \(\mathbb{R}\).

        \item \textbf{True.}

        \innerpf{The Nested Interval Property holds for any nested sequence of nonempty closed and bounded sets in \(\mathbb{R}\). If \(\{F_n\}\) is such a sequence with \(F_{n+1} \subseteq F_n\) for all \(n\), then the intersection \(\bigcap_{n=1}^\infty F_n\) is nonempty. This follows from the completeness of \(\mathbb{R}\), as every decreasing sequence of nonempty closed and bounded sets has a nonempty intersection. Therefore, replacing ``closed interval'' with ``closed set'' does not invalidate the property.}

        \item \textbf{True.}

        \innerpf{Let \(U\) be a nonempty open set. Then there exists \(x \in U\) and \(\epsilon > 0\) such that \((x - \epsilon,\ x + \epsilon) \subseteq U\). Since the rationals are dense in \(\mathbb{R}\), there exists a rational number \(q \in (x - \epsilon,\ x + \epsilon)\). Therefore, \(U\) contains a rational number.}

        \item \textbf{True.}

        \innerpf{Let \(F\) be a bounded infinite closed set. Since \(F\) is infinite and bounded, it must contain limit points. As \(F\) is closed, it contains its limit points. The real numbers are densely ordered with rationals between any two real numbers. Therefore, \(F\) must contain a rational number.}

        \item \textbf{True.}

        \innerpf{The \hyperref[def:3.1.1]{Cantor set} \(C\) is constructed as the intersection of a decreasing sequence of closed sets (finite unions of closed intervals). Since each of these sets is closed and the intersection of closed sets is closed, \(C\) is closed.}
    \end{enumerate}
}

\begin{exercise}
    {3.2.8} Assume \(A\) is an open set and \(B\) is a closed set. Determine if the following sets are definitely \hyperref[def:3.2.2]{open}, definitely \hyperref[def:3.2.5]{closed}, both, or neither.
    \begin{enumerate}
        \item \(\overline{A \cup B}\)
        \item \(A \setminus B = \{x \in A \mid x \notin B\}\)
        \item \((A^c \cup B)^c\)
        \item \((A \cap B) \cup (A^c \cap B)\)
        \item \(\overline{A}^c \cap \overline{A^c}\)
    \end{enumerate}
\end{exercise}

\sol{
    \begin{enumerate}
        \item \(\overline{A \cup B}\)

        The closure of any set is closed by \numrefthm{3.2.9}. Therefore, \(\overline{A \cup B}\) is closed.

        \textbf{Conclusion: Closed.}

        \item \(A \setminus B = \{x \in A \mid x \notin B\}\)

        Since \(B\) is closed, its complement \(B^c\) is open. The set \(A\) is open by assumption. The intersection of two open sets is open. Note that \(A \setminus B = A \cap B^c\). Therefore, \(A \setminus B\) is open.

        \textbf{Conclusion: Open.}

        \item \((A^c \cup B)^c\)

        By De Morgan's Law, \((A^c \cup B)^c = A \cap B^c\). Both \(A\) and \(B^c\) are open sets. The intersection of open sets is open. Thus, \((A^c \cup B)^c\) is open.

        \textbf{Conclusion: Open.}

        \item \((A \cap B) \cup (A^c \cap B)\)

        We can factor out \(B\):
        \[
        (A \cap B) \cup (A^c \cap B) = [ (A \cup A^c) \cap B ] = \mathbb{R} \cap B = B.
        \]
        Therefore, the set equals \(B\), which is closed.

        \textbf{Conclusion: Closed.}

        \item \(\overline{A}^c \cap \overline{A^c}\)

        Since \(A\) is open, its closure \(\overline{A}\) is a closed set containing all limit points of \(A\). Therefore, the complement \(\overline{A}^c\) is open.

        Similarly, \(A^c\) is closed (being the complement of an open set), so its closure is \(\overline{A^c} = A^c\), which is closed.

        Now, consider the intersection:
        \[
        \overline{A}^c \cap \overline{A^c} = \overline{A}^c \cap A^c
        \]
        Since \(\overline{A} \supseteq A\), we have \(\overline{A}^c \subseteq A^c\). Thus, the intersection simplifies to \(\overline{A}^c\).

        However, any point not in \(\overline{A}\) cannot be a limit point of \(A\) or belong to \(A\). In \(\mathbb{R}\), this set is empty unless \(A\) is either \(\emptyset\) or \(\mathbb{R}\).

        Therefore,
        \[
        \overline{A}^c \cap \overline{A^c} = \emptyset
        \]
        The empty set is both open and closed.

        \textbf{Conclusion: Both open and closed (since the set is empty).}

    \end{enumerate}
}

\begin{exercise}
    {3.2.11} \begin{enumerate}
        \item Prove that \(\overline{A \cup B} = \overline{A} \cup \overline{B}\).
        \item Does this result about closures extend to infinite unions of sets?
    \end{enumerate}
\end{exercise}

\sol{
    \begin{enumerate}
        \item We will prove that \(\overline{A \cup B} = \overline{A} \cup \overline{B}\).

        \innerpf{

            Recall that the closure of a set \(A\) is defined as \(\overline{A} = A \cup L_A\), where \(L_A\) is the set of limit points of \(A\).

            \textbf{Proof of \(\overline{A \cup B} \subseteq \overline{A} \cup \overline{B}\):}

            Let \(x \in \overline{A \cup B}\). Then \(x \in A \cup B\) or \(x\) is a limit point of \(A \cup B\).

            - If \(x \in A \cup B\), then \(x \in A\) or \(x \in B\), so \(x \in \overline{A}\) or \(x \in \overline{B}\), thus \(x \in \overline{A} \cup \overline{B}\).

            - If \(x\) is a limit point of \(A \cup B\), then every neighborhood \(V_\epsilon(x)\) contains a point \(y \ne x\) such that \(y \in A \cup B\). Therefore, \(y \in A\) or \(y \in B\), so \(x\) is a limit point of \(A\) or \(B\). Hence, \(x \in \overline{A}\) or \(x \in \overline{B}\), so \(x \in \overline{A} \cup \overline{B}\).

            Therefore, \(\overline{A \cup B} \subseteq \overline{A} \cup \overline{B}\).

            \textbf{Proof of \(\overline{A} \cup \overline{B} \subseteq \overline{A \cup B}\):}

            Let \(x \in \overline{A} \cup \overline{B}\). Then \(x \in \overline{A}\) or \(x \in \overline{B}\).

            - If \(x \in \overline{A}\), then \(x \in A\) or \(x\) is a limit point of \(A\). Since \(A \subseteq A \cup B\), \(x \in A \cup B\) or \(x\) is a limit point of \(A \cup B\). Thus, \(x \in \overline{A \cup B}\).

            - Similarly, if \(x \in \overline{B}\), then \(x \in \overline{A \cup B}\).

            Therefore, \(\overline{A} \cup \overline{B} \subseteq \overline{A \cup B}\).

            Hence, \(\overline{A \cup B} = \overline{A} \cup \overline{B}\).

        }

        \item The result does not necessarily extend to infinite unions of sets.

        Consider the sets \(A_n = \left( \frac{1}{n}, 1 - \frac{1}{n} \right)\) for \(n \in \N\). Then \(\overline{A_n} = \left[ \frac{1}{n}, 1 - \frac{1}{n} \right]\).

        The infinite union is \(A = \bigcup_{n=1}^\infty A_n = (0,1)\), so \(\overline{A} = [0,1]\).

        The union of the closures is \(\bigcup_{n=1}^\infty \overline{A_n} = (0,1)\), since none of the closed intervals \(\left[ \frac{1}{n}, 1 - \frac{1}{n} \right]\) include the endpoints 0 or 1.

        Therefore, \(\overline{A} \ne \bigcup_{n=1}^\infty \overline{A_n}\).

        Hence, the equality does not hold for infinite unions.

    \end{enumerate}
}

\renewcommand{\theenumi}{\arabic{enumi}}
\renewcommand{\labelenumi}{\theenumi.}
\section{Compact Sets}

\begin{definition}
    A set \(K \subseteq \R\) is a \textit{compact set} if every sequence from \(K\) has a convergent subsequence where the limit is also \(K\).
\end{definition}

\begin{theorem}
    A set \(K\) is compact if, and only if, it is closed and bounded.
\end{theorem}

\iffpf{Assume a set \(A \subseteq \R\) is closed and bounded. Thus, there exists a convergent subsequence by \namrefthm{Bolzano-Weierstrass Theorem}. Because \(A\) is closed, the limit is in the set by \defref{3.2.5}.}{Assume a set \(A\) is compact. If it is not bounded, then there exists an \((a_n)\) that heads toward infinity. This contradicts \defref{3.3.1}, so it must be bounded. Then, by the same definition, the limit points belong in the set, so it is closed.\qedhere}{}{xblue}

\begin{definition}
    Let \(A \in \R\). An \textit{open cover} for \(A\) is a collection of open sets \(\{O_\lambda \mid \lambda \in A\}\) whose union contains the set \(A\); that is \(A \subseteq \bigcup_{\lambda \in A}O_\lambda\). Given an open cover for \(A\), a \textit{finite subcover} is a finite sub-collection of open sets from the original open cover whose union still manages to completely contain \(A\).
\end{definition}

\begin{ntheorem}
    {Heine-Borel Theorem}Let \(K\) be a  subset of  \(\R\). All the following statements are equivalent in the sense that any of them implies the two others: 
    \begin{enumerate}[label=(\roman*)]
        \item \(K\) is compact.
        \item \(K\) is closed and bounded.
        \item Every open cover of \(K\) has a finite subcover.
    \end{enumerate}
\end{ntheorem}

\tpf{ The first set of ``if and only if proofs'' will be to prove (i) and (ii) are equivalent. Then, we will prove (ii) and (iii) are equivalent. 
    \begin{customframedproof}[linecolor=xgray]
        \begin{proofpart}[xgray]{(\(\Rightarrow\))}
            Assume \(K\) is compact. We need to show that \(K\) is closed and bounded. To show \(K\) is bounded, consider the open cover \(\mathcal{U} = \{(-n, n) \mid n \in \N\}\). \(\mathcal{U}\) covers all of \(\R\), so it certainly covers \(K\). Thus, there must exist a finite subcover. Consider the longest interval in the subcover. Then, \(K\) is a subset of this interval, so \(K\) is bounded. To show \(K\) is closed, we need to show every limit point belongs to \(K\). Assume \(x\) is a limit point of \(K\). From \defref{3.2.4} for every \(a \in K\), let \(\epsilon_a = \frac{1}{2}\disabs{a - x}\). Consider the open cover \(\mathcal{U} = \{V_{\epsilon_a}(a) \mid a \in K\} \). This covers every point on \(K\).
        \end{proofpart}
    \vspace{2mm}
            \begin{proofpart}[xgray]{(\(\Leftarrow\))}
                Because it is closed and bounded, by \numrefthm{3.3.2}, \(K\) is compact.
            \end{proofpart}
    \vspace{2mm}
    \end{customframedproof}
    \vspace{0.5cm}
    Now for the second part of the proof:

    \begin{customframedproof}[linecolor=xgray]
        \begin{proofpart}[xgray]{(\(\Rightarrow\))}
            Let \(x\) be a limit point of \(K\). This means there must exist a sequence \((x_n)\) in \(K\) with \(\limn{x_n} = x\). Suppose \(x \notin K\). For every \(y \in K\). Let \(\epsilon_y = \frac{1}{2}\disabs{y-x}\). Consider the open neighborhood \(V_{\epsilon_y}(y)\). Notice \(x \in V_{\epsilon_y} (y)\). Now, we will work with the collection of all such neighborhoods \(\mathcal{U} = \{V_{\epsilon_y}(y) \mid y \in K\}\). This \(\mathcal{U}\) is an open cover of \(K\). By our hypothesis there exists a finite subcover. There are some \(y_1,y_2,\dots,y_m\) such that \(K \subset \bigcup_{i = 1}^m V_{\epsilon_y}(y_i)\). Look at the distance from \(x\) to each \(y_i\): \((x - \epsilon_{y_i}\text{, } x + \epsilon_{y_i}) \cap V_{\epsilon_{y_i}} (y_i) = \emptyset\). Similar statements are for every \(y_i\). Let \(\epsilon = \min\{\epsilon_{y_1}, \epsilon_{y_2},\dots \epsilon_{y_M}\}\). Since there are infinitely many \(\epsilon > 0\), we see that \(V_\epsilon(x)\cap V_{\epsilon_y}(y_i) = \emptyset\) for every \(i \leq M\). So \(V_\epsilon (x) \cap K = \emptyset\). This gives us an \(\epsilon\)-neighborhood around \(x\) that does not intersect \(K\). Since \((x_n)\) approaches \(x\), there must be elements from the sequence that are inside of \(V_\epsilon(x)\). This creates a contradiction because we said \(x\) was a limit point. Therefore \(x \in u\) and \(K\) must be closed.
        \end{proofpart}
    \vspace{2mm}
            \begin{proofpart}[xgray]{(\(\Leftarrow\))}
                Let \(\mathcal{U}\) be an open cover of \(K\). Suppose there is no finite subcover. Since \(K\) is bounded there exists a closed interval \(I_0\) that contains \(K\). Bisect \(I_0\) and look at the two sub intervals \(A\) and \(B\). My claim is at least one of \(A \cap K\) and \(B \cap K\) does not have a finite subcover from \(\mathcal{U}\). If not, then we would have a finite subcover of all of K. Whichever half does not have a finites of cover will be called \(I_1\). Repeat this process. We get a sequence of nested closed intervals \(I_0 \supseteq I_1 \supseteq I_2 \supseteq \cdots\) such that for all \(j \in \N\), \(I_j \cap K\) does not have a finite subcover from \(\mathcal{U}\). Also, as the length of \(I_j\) approach 0, by the \namrefthm{Nested Interval Property}, there exists \(x \in \bigcap_{j = 1}^{\infty} I_j\). Since each \(I_j\) contains an element of \(K\) and the interval approaches 0, \(x\) must be a limit point of \(K\). Thus, since \(K\) is closed, \(x \in K\). There must be an open set \(U \in \mathcal{U}\) such that \(x \in U\). Since \(U\) is open and \(x \in U\), there exists an \(\epsilon > 0\) such that \(V_\epsilon(x) \subseteq U\). There is an \(I_j\) whose length is smaller than \(\epsilon\). This means \(I_j \subseteq V_\epsilon \subset U\). So \(\{U\}\) is a finite subcover of \(I_j \cap K\). This contradicts how we defined \(I_j\), therefore there must be a finite subcover from \(\mathcal{U}\). \qedhere
            \end{proofpart}
    \vspace{2mm}
    \end{customframedproof}
}


\begin{note}
    This allows us to take an infinite amount of \(\epsilon\)-neighborhoods and turn them into finite subcovers.
\end{note}
\begin{example}
    {Compactness}Let \(A = (0,1)\). Construct a set \(\mathcal{U}\) that is an open cover of \((0,1)\), but does not have a finite subcover.
\end{example}

\lesol{
    Consider \(\mathcal{U} = \{(0,t) \mid 0 < t < 1\}\). Thus, \(\mathcal{U}\) is an open cover, but does not contain a finite amount of subcovers because there will always be a point not covered. 
}

\renewcommand{\theenumi}{\alph{enumi}}
\renewcommand{\labelenumi}{(\theenumi)}
\subsection{Exercises}

\begin{exercise}
    {3.3.4}Assume \(K\) is \hyperref[def:3.3.1]{compact} and \(F\) is \hyperref[def:3.2.5]{closed}. Decide if the following sets are definitely compact, definitely closed, both, or neither. 
    \begin{enumerate}
        \item \(K \cap F\)
        \item \(\overline{F^c \cup K^c}\)
        \item \(K \setminus F = \{x \in K \mid x \notin F\}\)
        \item \(\overline{K \cap F^c}\)
    \end{enumerate}
\end{exercise}

\sol{
\begin{enumerate}
    \item Since \(K\) and \(F\) are closed, their intersection \(K \cap F\) is closed (\numrefthm{3.2.8}).

    To show that \(K \cap F\) is compact, let \(\mathcal{U}\) be any open cover of \(K \cap F\). Our goal is to extract a finite subcover from \(\mathcal{U}\). We can then use the \namrefthm{Bolzano-Weierstrass Theorem} (iii) to show that \(K \cap F\) is compact.

    Since \(F\) is closed, its complement \(F^c\) is open (\numrefthm{3.2.6}). Then \(K \setminus F = K \cap F^c\) is open as the intersection of an open set and \(K\).

    Consider the open cover \(\mathcal{U}' = \mathcal{U} \cup \{K \setminus F\}\) of \(K\). Every point in \(K\) is either in \(K \cap F\) (covered by \(\mathcal{U}\)) or in \(K \setminus F\) (covered by \(K \setminus F\)).

    Since \(K\) is compact, there exists a finite subcover \(\mathcal{U}'' \subseteq \mathcal{U}'\) that covers \(K\).

    If \(K \setminus F\) is in \(\mathcal{U}''\), remove it to obtain a finite subcollection of \(\mathcal{U}\) that still covers \(K \cap F\). If \(K \setminus F\) is not in \(\mathcal{U}''\), then \(\mathcal{U}'' \subseteq \mathcal{U}\) already covers \(K \cap F\).

    Therefore, \(K \cap F\) is compact.

    \textbf{Conclusion:} Both compact and closed.

    \item Since \(F\) and \(K\) are closed, \(F^c\) and \(K^c\) are open. The union \(F^c \cup K^c\) is open (\numrefthm{3.2.3}), so its closure \(\overline{F^c \cup K^c}\) is closed by \numrefthm{3.2.9}.

    This set may not be bounded, so it's not necessarily compact.

    \textbf{Conclusion:} Definitely closed.

    \item The set \(K \setminus F = K \cap F^c\) is the intersection of a compact set \(K\) and an open set \(F^c\). This set is open in \(K\) but not necessarily open or closed in \(\mathbb{R}\).

    Since \(K \setminus F\) is not necessarily closed, it may not be compact.

    \textbf{Conclusion:} Neither compact nor closed.

    \item The set \(K \cap F^c\) is open in \(K\), so its closure \(\overline{K \cap F^c}\) is closed by \numrefthm{3.2.9}.

    To show that \(\overline{K \cap F^c}\) is compact, let \(\mathcal{U}\) be any open cover of \(\overline{K \cap F^c}\).

    Since \(\overline{K \cap F^c} \subseteq K\) and \(K\) is compact, we can consider \(\mathcal{U}\) as an open cover of a subset of \(K\).

    By the definition of \hyperref[def:3.3.3]{open cover}, there exists a finite subcover of \(\mathcal{U}\) that covers \(\overline{K \cap F^c}\).

    Therefore, \(\overline{K \cap F^c}\) is compact.

    \textbf{Conclusion:} Both compact and closed.

\end{enumerate}
}


\begin{exercise}
    {3.2.8} Assume \(A\) is an open set and \(B\) is a closed set. Determine if the following sets are definitely \hyperref[def:3.2.2]{open}, definitely \hyperref[def:3.2.5]{closed}, both, or neither.
    \begin{enumerate}
        \item \(\overline{A \cup B}\)
        \item \(A \setminus B = \{x \in A \mid x \notin B\}\)
        \item \((A^c \cup B)^c\)
        \item \((A \cap B) \cup (A^c \cap B)\)
        \item \(\overline{A}^c \cap \overline{A^c}\)
    \end{enumerate}
\end{exercise}

\sol{
    \begin{enumerate}
        \item \(\overline{A \cup B}\)

        The closure of any set is closed by definition. Therefore, \(\overline{A \cup B}\) is definitely closed.

        \textbf{Conclusion: Closed.}

        \item \(A \setminus B = \{x \in A \mid x \notin B\}\)

        Since \(B\) is closed, its complement \(B^c\) is open. Since \(A\) is open, the intersection \(A \cap B^c = A \setminus B\) is the intersection of two open sets, which is open.

        \textbf{Conclusion: Open.}

        \item \((A^c \cup B)^c\)

        Applying De Morgan's Law:
        \[
        (A^c \cup B)^c = A \cap B^c
        \]
        Since \(A\) is open and \(B^c\) is open (because \(B\) is closed), their intersection \(A \cap B^c\) is open.

        \textbf{Conclusion: Open.}

        \item \((A \cap B) \cup (A^c \cap B)\)

        Simplify the expression:
        \[
        (A \cap B) \cup (A^c \cap B) = [A \cup A^c] \cap B = \R \cap B = B
        \]
        Thus, the set equals \(B\), which is closed.

        \textbf{Conclusion: Closed.}

        \item \(\overline{A}^c \cap \overline{A^c}\)

        Since \(A\) is open, its closure \(\overline{A}\) is closed, so \(\overline{A}^c\) is open.

        Since \(A^c\) is closed (being the complement of an open set), \(\overline{A^c} = A^c\) is closed.

        Therefore, \(\overline{A}^c \cap \overline{A^c}\) is the intersection of an open set and a closed set, which is generally open but not necessarily closed.

        For example, let \(A = (0,1)\). Then:
        \[
        \overline{A} = [0,1], \quad \overline{A}^c = (-\infty, 0) \cup (1,\infty)
        \]
        and
        \[
        \overline{A^c} = A^c = (-\infty, 0] \cup [1,\infty)
        \]
        Then:
        \[
        \overline{A}^c \cap \overline{A^c} = [(-\infty, 0) \cup (1,\infty)] \cap [(-\infty, 0] \cup [1,\infty)] = (-\infty, 0) \cup (1,\infty)
        \]
        Which is an open set.

        \textbf{Conclusion: Open.}

    \end{enumerate}
}

\chapter{Functional Limits and Continuity}\label{functional limits}
\vspace*{-0.25in}
\label{sec:Thomae}
\section{Discussion: Examples of Dirichlet and Thomae}

\begin{definition}
    The \textit{Dirichlet function} \(\limx c{g(x)}\) does not exist for any \(c \in \R\).
    \[
        g(x) = \begin{cases}
            1 & \text{if } x \in \Q,              \\
            0 & \text{if } x \in \R \setminus \Q.
        \end{cases}
    \]
\end{definition}
\begin{center}

    \begin{tikzpicture}

        \begin{axis}[
                axis lines = middle,
                xlabel = \(x\),
                ylabel = \(g(x)\),
                xmin = -1, xmax = 1,
                ymin = -0.5, ymax = 1.5,
                xtick = {-1, -0.5, 0, 0.5, 1},
                ytick = {0, 1},
            ]
            \addplot [
                domain=-1:1,
                samples=25,
                color=blue,
            ]
            {1};
            \addplot [
                domain=-1:1,
                samples=50,
                color=red,
            ]
            {0};

        \end{axis}

    \end{tikzpicture}
\end{center}

\begin{definition}
    The \textit{Thomae function} is defined as
    \[
        f(x) = \begin{cases}
            \frac{1}{q} & \text{if } x = \frac{p}{q} \text{ is in lowest terms}, \\
            0           & \text{if } x \in \R \setminus \Q.
        \end{cases}
    \]
\end{definition}

\textit{Thomae's function}, \(t(x)\) is continuous at all \(x \notin \Q\). It is not continuous at any \(x \in \Q\).

\newpage

\section{Functional Limits}

Recall from calculus I, that a function \(f(x)\) is continuous at \(x = c\) if \(\lim_{x\rightarrow c}f(x) = f(c)\).

\begin{definition}
    Let \(f \colon A \rightarrow \R\) be a function and let \(c\) be a limit point of \(A\). We say \(\lim_{x \rightarrow c}f(x) = L \), if for all \(\epsilon > 0\), there exists \(\delta > 0\) such that if \(0 < \disabs{x - c} < \delta\), then \(\disabs{f(x) - L} < \epsilon\).
\end{definition}


\begin{example}
    {Functional Limit (From book) 1}Let \(f(x) = 3x + 1\). Claim: \(\limx 2{f(x)} = 7\).
\end{example}

\pfscratch{
    Let \(\epsilon > 0\). After we have done our scratch work, we can choose \(\delta = \epsilon/ 3\), then \(0 < \disabs{x - 2} < \delta\) implies \(\disabs{f(x) - 7} < 3(\epsilon/3) = 3\).
}%
{
    \defref{4.2.1} requires that we produce a \(\delta > 0\) so that \(0 < \disabs{ x - 2} < \delta\) leads to the conclusion that \(\disabs{f(x) - 7} < \epsilon\). Notice that
    \[
        \disabs{f(x) - 7} = \disabs{3x + 1 - 7} = \disabs{3x - 6} = 3\disabs{x - 2}.
    \]
}{teal}


\begin{example}
    {Functional Limit (From book) 2}Let \(g(x) = x^{2}\). Claim: \(\limx 2{g(x)} = 4\).
\end{example}

\pfscratch{
    Let \(\epsilon > 0\). Choose \(\delta = \min\{1,\epsilon/5\}\). If \(0 < \disabs{x - 2} < \delta\), then
    \begin{align*}
        \disabs{g(x) - 4} & = \disabs{x^{2} - 4}           \\
                          & = \disabs{x - 2}\disabs{x + 2} \\
                          & < 5\delta                      \\
                          & = (5)\frac{\epsilon}{5}        \\
                          & = \epsilon.
    \end{align*}
}{
    Our goal this time is to make \(\disabs{g(x) - 4} < \epsilon\) by restricting \(\disabs{x - 2}\) to be smaller than some carefully chosen \(\delta\). As in the previous example, a little algebra reveals
    \[
        \disabs{g(x) - 4} = \disabs{x^{2} - 4} = \disabs{x - 2}\disabs{x + 2}.
    \]
    We can make \(\disabs{x + 2}\) as small as we like, but we need an upper bound on \(\disabs{x + 2}\) in order to know how small to choose \(\delta\). The presence of the variable \(x\) causes some initial confusion, but keep in mind that we are discussing the limt as \(x\) approaches 2. If we agree that our \(\delta-\)neighborhood around \(c = 2\) must have radius no bigger than \(\delta = 1\), then we get the upper bound \(\disabs{x + 2} < \disabs{3 + 2} = 5\) for all \(x \in V_{\delta}(c)\).
}{teal}


\begin{example}
    {Functional Limit 1}Let \(f(x) = 3x + 1\). Show that \(\lim_{x \rightarrow 2}f(x) = 7\).
\end{example}

\lepf{
    Let \(\epsilon > 0\). Set \(\delta = \frac{\epsilon}{3}\). Assume \(0 < \disabs{x - 2} < \delta\). Since \(\delta > 0\), \(2 - \delta < x < 2 + \delta\). Then,
    \begin{align*}
        \disabs{x - 2}    & < \delta,             \\
        \disabs{f(x) - 7} & = \disabs{3x + 1 - 7} \\
                          & = \disabs{3x - 6}     \\
                          & = 3\disabs{x - 2}     \\
                          & < 3\delta             \\
                          & = \epsilon.
    \end{align*}
    Therefore, \(\lim_{x \rightarrow 2} f(x) = 7\).
}



\begin{example}
    {Functional Limit 3}Let \(f(x) = x^{2}\). Claim: \(\limx 7{f(x)} = 49\)
\end{example}

\pfscratch{
    Let \(\epsilon > 0\). Set \(\delta = \min\{\frac{\epsilon}{8}, 1\}\). If \(0 < \disabs{x - 7} < \delta\), then \begin{align*}
        \disabs{f(x) - 49} & = \disabs{x^{2} - 49}              \\
                           & = \disabs{x - 7}\disabs{x + 7}     \\
                           & < 8\delta                          \\
                           & < 8\left(\frac{\epsilon}{8}\right) \\
                           & = \epsilon.
    \end{align*}
}%
{
    Always start with the goal statement: \(\disabs{f(x) - 49} = \disabs{x^{2} - 49}\). This factors into \(\disabs{x - 7} \disabs{x + 7}\). Then, if \(\delta < 1\), \(\disabs{x - 7} < \delta\) and \(\disabs{x + 7} < 8\). All together, we have \(8\delta < \epsilon < \frac{\epsilon}{8}\).
}%
{teal}

\begin{example}
    {Functional Limit 4}Claim: \(\limx 3{\frac{1}{x + 1}} = \frac{1}{4}\).
\end{example}

\pfscratch{
    Let \(\epsilon > 0\). Set \(\delta = \min\{12\epsilon,1\}\). If \(0 < \disabs{x - 3} < \delta\), then \begin{align*}
        \disabs{\frac{1}{x + 1} - \frac{1}{4}} & = \disabs{\frac{4 - (x + 1)}{4(x + 1)}} \\
                                               & = \disabs{\frac{3 - x}{4(x + 1)}}       \\
                                               & < \frac{\delta}{4(3)}                   \\
                                               & = \frac{12\epsilon}{12}                 \\
                                               & = \epsilon.
    \end{align*}
    Therefore, \(\limx 3{\frac{1}{x + 1}} = \frac{1}{4}\)
}{
    Goal: \(\disabs{\frac{1}{x + 1} - \frac{1}{4}}\). Hence,
    \begin{align*}
        \disabs{\frac{1}{x + 1} - \frac{1}{4}} & = \disabs{\frac{4 - (x + 1)}{4(x + 1)}} \\
                                               & = \disabs{\frac{3 - x}{4(x + 1)}}       \\
                                               & < \frac{\delta}{4\disabs{x + 1}}        \\
                                               & < \frac{\delta}{4(3)}                   \\
                                               & = \frac{\delta}{12}                     \\
                                               & < \epsilon.
    \end{align*}
    Thus, we need a \(\delta < 1\), and we can choose \(\delta = \min\{12\epsilon,1\}\). Note: When we are determining the value for \(\disabs{x + 2}\), we solve for \(\delta = 3 \pm 1 \Rightarrow x \in (2,4)\). Then, we find \(x + 1 = (3,5)\). We choose 3 rather than 5 because of division. We want to be as close as possible.
}{teal}

\begin{example}
    {Functional Limit 5}Claim: \(\limx 3{(x^{2} + 7x)} = 30\).
\end{example}

\lepf{
    Let \(\epsilon > 0\) and set \(\delta = \min\{\frac{\epsilon}{14}, 1\}\). If \(0 < \disabs{x - 3} < \delta\), then
    \begin{align*}
        \disabs{x^{2} + 7x - 30} & = \disabs{x - 3}\disabs{x + 10}      \\
                                 & < 14\delta                           \\
                                 & = 14\left(\frac{\epsilon}{14}\right) \\
                                 & = \epsilon.
    \end{align*}
}

\begin{example}
    {Functional Limit 6}Claim: \(\limx 3{\frac{2x + 3}{4x - 9}} = 3\).
\end{example}

\lepfscratch{
    Let \(\epsilon > 0\). Set \(\delta = \min\{\frac{\epsilon}{10}, \frac{1}{2}\}\). (Note: We are choosing \(\frac{1}{2}\) because we want to avoid having 0 anywhere in the interval.) Assume \(0 < \disabs{x - 3} < \delta\). Since \(\delta < \frac{1}{2}\), \(\frac{5}{2} < x < \frac{7}{2}\), then \(1 < \disabs{4x - 9} < 5\). (Thus, 0 can not possibly be in the denominator.)
}{
    \begin{align*}
        \disabs{\frac{2x + 3}{4x + 9} - 3} & = \disabs{\frac{2x + 3 - 3(4x + 9)}{4x + 9}} \\
                                           & = \disabs{\frac{2x + 3 - 12x - 27}{4x + 9}}  \\
                                           & = 10\disabs{\frac{x-3}{4x - 4}}              \\
                                           & < 10 \frac{\epsilon / 10}{1}                 \\
                                           & = \epsilon.
    \end{align*}}

\begin{example}
    {Functional Limit 7}Claim: \(\limx 4{\sqrt{x}} = 2\).
\end{example}

\lepfscratch{
    Let \(\epsilon > 0\). Set \(\delta = \min\{1,3\epsilon\}\). Assume \(0 < \disabs{x - 4} < \delta\). Then (refer to scratch work).
}{
    \begin{align*}
        \disabs{\sqrt{x} - 2} & = \disabs{\sqrt{x} - 2}                                             \\
                              & = \disabs{\frac{(\sqrt{x} - 2) \cdot (\sqrt{x} + 2)}{\sqrt{x} + 2}} \\
                              & = \disabs{\frac{x - 4}{\sqrt{x} + 2}}                               \\
                              & < \frac{\delta}{3}                                                  \\
                              & < \frac{3\epsilon}{3}                                               \\
                              & = \epsilon.
    \end{align*}
    Notice that we picked \(\delta < 1\) such that \(3 < x < 4\) so \(1 < \sqrt{x} < 2\) and \(3 < \sqrt{x} + 2 < 4\).
}


\begin{ntheorem}
    {Sequential Criterion for Functional Limits}
    The following statements are equivalent:

    \begin{enumerate}
        \item \(\lim_{x \rightarrow c} f(x) = L\).
        \item For all sequences \((x_n)\) where \(x_n \ne c\) and \((x_n) \rightarrow c\), we have \(\limn f(x_n) = L\).
    \end{enumerate}
\end{ntheorem}

\tpf{
    \((1) \rightarrow (2)\) \\

    Assume \(\limc{f(x)} = L\).

    Let \((x_n) \rightarrow c\) with \(x_n \ne c\)

    Let \(\epsilon > 0\).
    \begin{itemize}
        \item Since \(\limc{f(x)} = L\), there exists \(\delta > 0\) such that if \(0 < \disabs{x - c} < \delta\), then \(\disabs{f(x) - L} < \epsilon\).

        \item Since \(x_n \rightarrow c\), there exists \(N \in \N\) such that for all \(n \geq N\), \(\disabs{x_n - c} < \delta\).

        \item Now, for all \(n \geq N\), it follows that \(x_n - c < \delta\) and thus \(\disabs{f(x) - L} < \epsilon\).
    \end{itemize}

    Thus, \(\limn f(x_n) = L\). \\

    \((2) \rightarrow (1)\) \\

    Proof by contrapositive.

    Assume (1) is not true. Thus,
    \[
        \limc{f(x)} \ne L.
    \]

    There exists \(\epsilon_0 > 0\) such that for all \(\delta > 0\), there exists an \(x\) with \(0 < \disabs{x - c} < \delta\) and \(\disabs{f(x) - L} \geq \epsilon_0\).

    For each \(n \in \N\), consider \(\delta = \frac{1}{n}\). There exists \(x_n \in (c - \frac{1}{n}, c + \frac{1}{n})\) with \(x_n \ne c\) such that \(\disabs{f(x) - L} \geq \epsilon_0\).

    \begin{itemize}
        \item Since \(\disabs{x_n - c} < \frac{1}{n}\), we see that \((x_n) \rightarrow c\).
        \item Since for all \(\ninn\), \(\disabs{f(x) - L} \geq \epsilon_0\). Then, \(\limn f(x) \ne L\).
    \end{itemize}
    Thus, \(\neg(1) \rightarrow \neg(2)\). So \((2) \rightarrow (1)\) and \((1) \rightarrow (2)\).
}

If functional limits and sequential limits are the same thing, then everything we know about sequential limits is also true about functional limits.

Recall \namrefthm{Algebraic Limit Theorem}. From this, we can write the functional equivalent:

Assume \(\limc{f(x)} = L\) and \(\limc{g(x)} = M\). Then,
\begin{itemize}
    \item \(\limc{(f(x) + g(x))} = L + M\)
    \item \(\limc{(f(x) - g(x))} = L - M\)
    \item \(\limc{(f(x)g(x))} = LM\)
    \item \(\limc{(\frac{f(x)}{g(x)})} = \frac{L}{M}\) unless \(M = 0\).
\end{itemize}

\begin{ntheorem}
    {Divergence Criterion} Let \(f \colon A \rightarrow \R\) with \(c\) as a limit point of \(A\). If there exists two sequences \((x_n)\) and \((y_n)\) in \(A \setminus \{c\}\) (that both converge to \(c\)) such that \(\limn f(x_n) \ne \limn f(y_n)\), then \(\limc{f(x)}\) does not exist.
\end{ntheorem}

\begin{example}
    {Divergence Criterion 1} \(f(x) = \frac{x}{|x|} = \begin{cases}
        1, \text{ if } x > 0 \\
        -1, \text{ if } x < 0
    \end{cases}\)
    Our goal is to show that \(\lim_{x \rightarrow 0} f(x)\) does not exist.
\end{example}

\lepf{
    Let \((x_n) = (\frac{1}{n})\) and let \((y_n) = (\frac{-1}{n})\). We will see that as \(n \rightarrow \infty\), \(\limn f(x_n) = 1\) and \(\limn f(y_n) = -1\). Thus, \(\lim_{x \rightarrow 0} f(x)\) does not exist.
}

\begin{example}
    {Divergence Criterion 2} \(g(x) = \sin(\frac{1}{x})\). Show that \(\lim_{x \rightarrow 0} g(x)\) does not exist.
\end{example}

\lepf{
    Let \((x_n) = (\frac{1}{2\pi n})\) and let \((y_n) = (\frac{1}{2\pi n + \frac{\pi}{2}})\). We will see that as \(n \rightarrow \infty\), \(\limn g(x_n) = 1\) and \(\limn g(y_n) = -1\). Thus, \(\lim_{x \rightarrow 0} g(x)\) does not exist.
}



We say \(\limn x_n = \infty\) if for all \(M > 0\), there exists \(N \in \N\) such that for all \(n \geq N\), \(x_n > M\).

We say \(\limc{f(x)} = \infty\) if for all \(M > 0\), there exists \(\delta > 0\) such that if \(0 < \disabs{x - c} < \delta\), then \(f(x) > M\). Think of vertical asymptotes.


\begin{ntheorem}
    {Infinite Limits Cauchy Criterion} If \((x_n) \rightarrow \infty\), \((x_n)\) will not be Cauchy. It is possible to have \(x_{n+1} - x_n\) approach 0, but \((x_n)\) converges to \(\infty\).
\end{ntheorem}


% \renewcommand{\theenumi}{\alph{enumi}}
% \renewcommand{\labelenumi}{(\theenumi)}
% \subsection{Exercises}

% \begin{exercise}
%     {4.2.10 (Right and Left Limits)} Introductory calculus courses typically refer to the right-hand limit of a function as the limit obtained by “letting \(x\) approach \(a\) from the right-hand side.”
%     \begin{enumerate}
%         \item[(a)] Give a proper definition in the style of \defref{4.2.1} for the right-hand and left-hand limit statements:
%               \[
%                   \lim_{x \to a^+} f(x) = L \quad \text{and} \quad \lim_{x \to a^-} f(x) = M.
%               \]

%         \item[(b)] Prove that \(\lim_{x \to a} f(x) = L\) if and only if both the right and left-hand limits equal \(L\).
%     \end{enumerate}
% \end{exercise}

% \sol{
%     \begin{enumerate}
%         \item[(a)] \textbf{Definition of Right-Hand Limit:} \\
%         Let \(f: A \to \mathbb{R}\) and let \(c\) be a limit point of \(A\). We say that \(\lim_{x \to c^+} f(x) = L\) if for every \(\epsilon > 0\), there exists \(\delta > 0\) such that whenever \(0 < x - c < \delta\) (and \(x \in A\)), it follows that \(|f(x) - L| < \epsilon\).

%         \textbf{Definition of Left-Hand Limit:} \\
%         Let \(f: A \to \mathbb{R}\) and let \(c\) be a limit point of \(A\). We say that \(\lim_{x \to c^-} f(x) = M\) if for every \(\epsilon > 0\), there exists \(\delta > 0\) such that whenever \(-\delta < x - c < 0\) (and \(x \in A\)), it follows that \(|f(x) - M| < \epsilon\).

%         \item[(b)] 
%             \iffpf{
%                 Suppose that \(\lim_{x \to c} f(x) = L\). By definition, for every \(\epsilon > 0\), there exists \(\delta > 0\) such that whenever \(0 < |x - c| < \delta\), we have \(|f(x) - L| < \epsilon\). This inequality holds for all \(x\) approaching \(c\) from both sides, so both the right-hand and left-hand limits equal \(L\).
%                 }{
%                 Conversely, suppose that \(\lim_{x \to c^+} f(x) = L\) and \(\lim_{x \to c^-} f(x) = L\). Then, for every \(\epsilon > 0\), there exist \(\delta_1, \delta_2 > 0\) such that:
%             \begin{itemize}
%                 \item If \(0 < x - c < \delta_1\), then \(|f(x) - L| < \epsilon\).
%                 \item If \(-\delta_2 < x - c < 0\), then \(|f(x) - L| < \epsilon\).
%             \end{itemize}
%             Let \(\delta = \min\{\delta_1, \delta_2\}\). Then, for all \(x\) with \(0 < |x - c| < \delta\), we have \(|f(x) - L| < \epsilon\).
%             }{
%                 Therefore, \(\lim_{x \to c} f(x) = L\).
%                 }
%             {xgray}
%     \end{enumerate}
% }
% % \parbox{0.90\textwidth}{
% \begin{exercise}
%     {4.2.11 (Squeeze Theorem)}Let \(f\), \(g\), and \(h\) satisfy \(f(x) \leq g(x) \leq h(x)\) for all \(x\) in some common domain \(A\). If \(\lim_{x \to c} f(x) = L\) and \(\lim_{x \to c} h(x) = L\) at some limit point \(c\) of \(A\), show \(\lim_{x \to c} g(x) = L\) as well.
% \end{exercise}

% \expf{
%     Let \(\epsilon > 0\). Since \(\lim_{x \to c} f(x) = L\), there exists \(\delta_1 > 0\) such that whenever \(0 < |x - c| < \delta_1\), we have \(|f(x) - L| < \epsilon\). Similarly, there exists \(\delta_2 > 0\) such that whenever \(0 < |x - c| < \delta_2\), we have \(|h(x) - L| < \epsilon\). Let \(\delta = \min\{\delta_1, \delta_2\}\). Then, for all \(x\) with \(0 < |x - c| < \delta\), we have \(f(x) \leq g(x) \leq h(x)\). Thus, \(|g(x) - L| < \epsilon\), and \(\lim_{x \to c} g(x) = L\).
% }


% \renewcommand{\theenumi}{\arabic{enumi}}
% \renewcommand{\labelenumi}{\theenumi.}

\section{Continuous Functions}

\begin{definition}
    We say a function \(f\) is \textit{continuous} at \(c\) if
    \[
        \limc{f(x)} = f(c).
    \]

    Equivalent definition:

    For all \(\epsilon > 0\), there exists \(\delta > 0\) such that if \(\disabs{x - c} < \delta\), then
    \[
        \disabs{f(x) - f(c)} < \epsilon.
    \]
\end{definition}

Idea: When \(x\) is close to \(c\), \(f(x)\) is close to \(f(c)\). Then, for the topological definition, we can say if \(x \in V_\delta(c)\) then \(f(x) \in V_\epsilon(f(c))\).

\begin{definition}
    We say function \(f\) is \textit{continuous on a set} \(D\) if \(f\) is continuous at every point in \(D\).
\end{definition}

The following are equivalent (TFAE):

\begin{enumerate}
    \item \(\limc{f(x)} = L\)
    \item For all sequences \((x_n)\) such that \((x_n) \rightarrow c\), we have \(\limn f(x_n) = L\).
\end{enumerate}

\begin{center}
    \underline{Continuous Functions} (THM 4.3.2 in book)
\end{center}

Claim: Let \(a \in \R\). Then \(f(x) = a\) is continuous.

\begin{customframedproof}[linecolor=xgray]
    \begin{proof}
        Let \(c \in \R\). Let \(\epsilon > 0\). Set \(\delta = 8\). Now, if \(x - c < 8\), then \(\disabs{f(x) - f(c)} = a - a = 0 < \epsilon\). Thus, constant functions are continuous.
    \end{proof}
\end{customframedproof}

Claim: \(f(x) = x\) is continuous.

\begin{customframedproof}[linecolor=xgray]
    \begin{proof}
        Let \(c \in \R\). Let \(\epsilon > 0\). Set \(\delta = \epsilon\). If \(\disabs{x - c} < \delta\), then \(\disabs{f(x) - f(c)} = \disabs{x - c} < \delta = \epsilon\). Thus, the identity function is continuous.
    \end{proof}
\end{customframedproof}

% \namrefthm{Algebraic Limit Theorem} tells us that if \(f\) and \(g\) are continuous at \(c\), then \(f + g\), \(f - g\), \(fg\), and \(\frac{f}{g}\) (so long as \(g \ne 0\)) are continuous at \(c\). From this, we know that polynomials are continuous. We'll show that in the next theorem.

% \begin{ntheorem}
%     {Thm. 4.3.2 in book}
% \end{ntheorem}

Claim: \(g(x) = \sqrt{x}\) is continuous on \([0,\infty)\).

\begin{customframedproof}[linecolor=xgray]
    \begin{proof}
        \begin{itemize}
            \item \textbf{Case 1:} \(c \ne 0\) \\

                  Let \(c \in [0,\infty)\). Let \(\epsilon > 0\). Set \(\delta < \epsilon\). If \(\disabs{x - c} < \delta\), then \(\disabs{g(x) - g(c)} = \disabs{\sqrt{x} - \sqrt{c}} = \frac{\disabs{x - c}}{\sqrt{x} + \sqrt{c}} < \frac{\delta}{\sqrt{c}} < \epsilon\). Thus, \(g(x) = \sqrt{x}\) is continuous on \([0,\infty)\).

            \item \textbf{Case 2:} \(c = 0\) \\

                  Let \(c \in [0,\infty)\). Let \(\epsilon > 0\) Set \(\delta = \epsilon^2\). If \(\disabs{x - 0} < \delta\), then \(\disabs{g(x) - g(0)} = \disabs{\sqrt{x} - 0} = \sqrt{x} < \sqrt{\delta} = \epsilon\). Thus, \(g(x) = \sqrt{x}\) is continuous on \([0,\infty)\).
        \end{itemize}
        Therefore, \(g(x)\) is continuous for all \(c \in [0,\infty)\).\qedhere
    \end{proof}
\end{customframedproof}

\begin{ntheorem}
    {Compositions of Continuous Functions} Let \(f\) be continuous at \(c\). Let \(g\) be continuous at \(f(x)\). Then,
    \[
        g \circ f(x) = g(f(x)) \text{ is continuous.}
    \]
\end{ntheorem}

\tpf{
    Let \(\epsilon > 0\). Since \(g\) is continuous at \(f(c)\), there exists \(\delta_1 > 0\) such that if \(\disabs{x - f(c)} < \delta_1\), then \(\disabs{g(x) - g(f(c))} < \epsilon\). Since \(f\) is continuous at \(c\), there exists \(\delta_2 > 0\) such that if \(\disabs{x - c} < \delta_2\), then \(\disabs{f(x) - f(c)} < \delta_1\). Thus, if \(\disabs{x - c} < \delta_2\), then \(\disabs{g(f(x)) - g(f(c))} < \epsilon\). Therefore, \(g \circ f(x)\) is continuous.
}

\begin{center}
    \underline{Most Common Applications of Continuity} \\

    Is with limits.
\end{center}

If \(f\) is continuous at \(c\) and \((x_n) \to c\), then \(\limn f(x_n) = f(c)\). Hence.
\[
    \limn f(x_n) = f(\limn x_n).
\]

% \renewcommand{\theenumi}{\alph{enumi}}
% \renewcommand{\labelenumi}{(\theenumi)}
% \subsection{Exercises}

% \begin{exercise}
%     {4.3.1} Let \( g(x) = \sqrt[3]{x} \).
%     \begin{enumerate}
%         \item Prove that \( g \) is continuous at \( c = 0 \).
%         \item Prove that \( g \) is continuous at a point \( c \neq 0 \). (The identity \( a^3 - b^3 = (a - b)(a^2 + ab + b^2) \) will be helpful.)
%     \end{enumerate}
% \end{exercise}

% \sol{
%     \begin{enumerate}
%         \item \innerpf{
%             To prove that \( g \) is continuous at \( c = 0 \), we need to show that
%             \[
%                 \lim_{x \to 0} g(x) = g(0) = 0.
%             \]
%             Let \( \epsilon > 0 \). Choose \( \delta = \epsilon^3 \). Then, if \( |x - 0| = |x| < \delta \), we have
%             \[
%                 |g(x) - g(0)| = |\sqrt[3]{x} - 0| = |\sqrt[3]{x}| = |x|^{1/3} < \delta^{1/3} = (\epsilon^3)^{1/3} = \epsilon.
%             \]
%             Therefore, \( g \) is continuous at \( c = 0 \).
%         }
%         \item \innerpf{
%             Let \( c \neq 0 \) and \( \epsilon > 0 \). We need to find \( \delta > 0 \) such that whenever \( |x - c| < \delta \), it follows that \( |g(x) - g(c)| < \epsilon \).

%             Using the identity \( a^3 - b^3 = (a - b)(a^2 + ab + b^2) \), set \( a = \sqrt[3]{x} \) and \( b = \sqrt[3]{c} \). Then,
%             \[
%                 x - c = (\sqrt[3]{x} - \sqrt[3]{c})\left( (\sqrt[3]{x})^2 + \sqrt[3]{x}\sqrt[3]{c} + (\sqrt[3]{c})^2 \right).
%             \]
%             Therefore,
%             \[
%                 \sqrt[3]{x} - \sqrt[3]{c} = \frac{x - c}{ (\sqrt[3]{x})^2 + \sqrt[3]{x}\sqrt[3]{c} + (\sqrt[3]{c})^2 }.
%             \]
%             Since \( c \neq 0 \) and \( x \) is close to \( c \), we can choose \( \delta_1 > 0 \) such that \( |x - c| < \delta_1 \) implies \( |x| \geq \dfrac{|c|}{2} > 0 \). This ensures that \( \sqrt[3]{x} \) and \( \sqrt[3]{c} \) are bounded away from zero.

%             Let
%             \[
%                 M = \min_{|x - c| < \delta_1} \left| (\sqrt[3]{x})^2 + \sqrt[3]{x}\sqrt[3]{c} + (\sqrt[3]{c})^2 \right| > 0.
%             \]
%             Choose \( \delta = \min\{ \delta_1, M\epsilon \} \). Then, if \( |x - c| < \delta \), we have
%             \[
%                 |g(x) - g(c)| = \left| \frac{x - c}{ (\sqrt[3]{x})^2 + \sqrt[3]{x}\sqrt[3]{c} + (\sqrt[3]{c})^2 } \right| \leq \frac{|x - c|}{M} < \frac{M\epsilon}{M} = \epsilon.
%             \]
%             Therefore, \( g \) is continuous at \( c \neq 0 \).
%         }
%     \end{enumerate}
% }

% \begin{exercise}
%     {4.3.8} Decide if the following claims are true or false, providing either a short proof or counterexample to justify each conclusion. Assume throughout that \( g \) is defined and continuous on all of \( \mathbb{R} \).
%     \begin{enumerate}
%         \item If \( g(x) \geq 0 \) for all \( x < 1 \), then \( g(1) \geq 0 \) as well.
%         \item If \( g(r) = 0 \) for all \( r \in \mathbb{Q} \), then \( g(x) = 0 \) for all \( x \in \mathbb{R} \).
%     \end{enumerate}
% \end{exercise}

% \sol{
%     \begin{enumerate}
%         \item \textbf{True.}

%         Since \( g \) is continuous at \( x = 1 \), we have
%         \[
%             g(1) = \lim_{x \to 1} g(x).
%         \]
%         For all \( x < 1 \), \( g(x) \geq 0 \), so the left-hand limit satisfies
%         \[
%             \lim_{x \to 1^-} g(x) \geq 0.
%         \]
%         Therefore,
%         \[
%             g(1) = \lim_{x \to 1} g(x) = \lim_{x \to 1^-} g(x) \geq 0.
%         \]
%         Thus, \( g(1) \geq 0 \).

%         \item \textbf{True.}

%         Since \( g \) is continuous on \( \mathbb{R} \) and \( g(r) = 0 \) for all rational numbers \( r \), it follows that for any \( x \in \mathbb{R} \), we can find a sequence of rational numbers \( (r_n) \) converging to \( x \). By continuity,
%         \[
%             g(x) = \lim_{n \to \infty} g(r_n) = \lim_{n \to \infty} 0 = 0.
%         \]
%         Therefore, \( g(x) = 0 \) for all \( x \in \mathbb{R} \).
%     \end{enumerate}
% }

% \begin{exercise}
%     {4.3.9} Assume \( h \colon \mathbb{R} \to \mathbb{R} \) is continuous on \( \mathbb{R} \) and let \( K = \{x \colon h(x) = 0\} \). Show that \( K \) is a closed set.
% \end{exercise}

% \expf{
%     To show that \( K \) is closed, we need to show that its complement \( \mathbb{R} \setminus K \) is open. Alternatively, we can show that \( K \) contains all its limit points.

%     Let \( (x_n) \) be a sequence in \( K \) that converges to some limit \( x \in \mathbb{R} \). Since \( h \) is continuous on \( \mathbb{R} \) and \( h(x_n) = 0 \) for all \( n \), it follows that
%     \[
%         \lim_{n \to \infty} h(x_n) = h\left( \lim_{n \to \infty} x_n \right) = h(x).
%     \]
%     Therefore, \( h(x) = 0 \), which means \( x \in K \).

%     Since every limit point of \( K \) is contained in \( K \), the set \( K \) is closed.
% }



% \renewcommand{\theenumi}{\arabic{enumi}}
% \renewcommand{\labelenumi}{\theenumi.}

\section{Continuous Functions on Compact Sets}

\begin{ntheorem}
    {Extreme Value Theorem} If \(K\) is compact and \(f : K \to \R\) is continuous, then \(f\) attains a maximum and minimum value on \(K\).

    In other words, there exists \(a \in K\) such that \(f(a) = \sup\{f(x) \mid x \in K\}\). Also, there exists \(b \in K\) such that \(f(b) = \inf\{f(x) \mid x \in K\}\).
\end{ntheorem}

\tpf{
    We know \(f\) is bounded on \(K\) from \hyperref[lm:4.4.2]{Lemma 4.4.2} below. Hence,
    \[
        S = \sup\{f(x) \mid x \in K\} \text{ exists.}
    \]
    For every natural number, there exists an \(x_n \in K\) such that \(S - \frac{1}{n} < f(x_n) \leq S\). It follows that \(\limn f(x_n) = S\). So, now we have a sequence, \((x_n)\) in the compact set \(K\). Since \(K\) is compact, by the Bolzano-Weierstrass Theorem, there exists a convergent subsequence,
    \[
        (x_{n_j}) \text{ with } a = \lim_{j \to \infty} x_{n_j} \in K.
    \]
    Since \(f\) is continuous, we have,
    \[
        f(a) = \lim_{j \to \infty} f(x_{n_j}) = S.
    \]
    By a similar method, there exists \(b \in K\) such that \[
        f(b) = \inf\{f(x) \mid x \in K\}.
    \]
    \textbf{Note:} This proof hinges on the fact that \(f\) is bounded! We need to show that \(f\) is bounded on \(K\) with a proof with subcovers.
}

\begin{lemma}
    How do we know \(f\) is bounded on \(K\)? That is,
    \[
        f(K) = \{f(x) \mid x \in K\}.
    \]
    Show that \(f(K)\) is bounded.
\end{lemma}


\lemmapf{
    Let \(c \in K\). Since \(f\) is continuous at \(c\), there exists \(\delta_c > 0\) such that if \(\disabs{x - c} < \delta_c\), then
    \[
        \disabs{f(x) - f(c)} < 1.
    \]
    Do this over every \(c \in K\).  We get an open cover of \(K\).
    \[
        \mathcal{O} = \{V_{\delta_c} (c) \mid c \in K\}.
    \]
    Since \(K\) is compact, the Heine-Borel Theorem says there exists a finite subcover. We get \(c_1, c_2 \dots, c_n \in K\) such that \(K \subseteq \bigcup_{i = 1}^n V_{\delta_{c_i}} (c_i)\). Thus,
    \begin{align*}
        f(K) & \subseteq f\left(\bigcup_{i = 1}^n V_{\delta_{c_i}} (c_i)\right) \\
             & \subseteq \bigcup_{i = 1}^n f(V_{\delta_{c_i}} (c_i))            \\
             & \subseteq \bigcup_{i = 1}^n (f(c_i) - 1, f(c_i) + 1)             \\
             & \subseteq [\min \{f(c_i)\} - 1, \max \{f(c_i)\} + 1].
    \end{align*}
    Therefore, \(f(K)\) is bounded.
}

\begin{ntheorem}
    {Preservation of Compact Sets}If \(K\) is compact and \(f \colon K \to \R\) is continuous, then \(f(K)\) is compact.
\end{ntheorem}

\tpf{
    Let \((y_n)\) be a sequence in \(f(K)\). We will show \((y_n)\) has a convergent subsequence with its limit in \(f(K)\).

    For each \(n\) there exists \(x_n \in K\) such that \(f(x_n) = y_n\). So \((x_n)\) is a sequence in a compact set \(K\). There exists a convergent subsequence \((x_{n_j})\) with
    \[
        a
        = \lim_{j \to \infty} x_{n_j} \in K.
    \]
    Now consider the corresponding subsequence \((y_{n_j})\) in \(f(K)\). Since \(f\) is continuous, we have
    \begin{align*}
        \lim_{j \to \infty} y_{n_j} = \lim_{j \to \infty} f(x_{n_j}) \\
         & = f(a) \in f(K).
    \end{align*}
    So, \(x_{n_j}\) is a convergent subsequence with limit in \(f(K)\). Therefore, \(f(K)\) is compact.
}

\begin{definition}
    A function \(f \colon A \to \R\) is \textit{uniformly continuous} on \(A\) if for all \(\epsilon > 0\), there exists \(\delta > 0\) such that for all \(x, c \in A\), if \(\disabs{x - c} < \delta\), then \(\disabs{f(x) - f(c)} < \epsilon\).
\end{definition}

Compare this definition with \defref{4.3.1}. The difference is that the \(\delta\) is independent of \(x\). That is, we have to find one \(\delta\) that needs to work for every point \(x\).

% Put a table environment here with 3x2.

% 1x1: Continuous: 
% 1x2: Uniformly Continuous:
% 2x1: For all \(c \in A\), we can find a corresponding \(\delta\).
% 2x2: For all \(c \in A\), we need one \(\delta\) that works at every \(c\).
% 3x

\begin{definition}
    A function \(f \colon A \to \R\) is \textit{\underline{not} uniformly continuous} on \(A\) if there exists \(\epsilon_{0} > 0\) such that for all \(\delta > 0\), there exists \(x, c \in A\) such that \(\disabs{x - c} < \delta\) and \(\disabs{f(x) - f(c)} \geq \epsilon_{0}\).
\end{definition}

\begin{theorem}
    If \(K \subseteq \R\) is compact and \(f \colon K \to \R\) is continuous, then \(f\) is uniformly continuous on \(K\).
\end{theorem}

\tpf{
    Suppose \(f\) is not uniformly continuous on \(K\). Then, there exists \(\epsilon_{0} > 0\) such that for all \(\ninn\), there exists \(x_{n},y_{n} \in K\) such that \(\disabs{x_{n} - y_{n} < \frac{1}{n}}\) and \(\disabs{f(x_{n}) - f(y_{n})} \geq \epsilon_{0}\). 

    We now have two sequences \((x_{n})\) and \((y_{n})\) in \(K\). Since \(K\) is compact, by the \namrefthm{heine-borel theorem}, there exists a convergent subsequence \((x_{n_{i}})\) which converges to a point \(x_{0} \in K\).
    
    Since \(K\) is compact, \((y_{n_{i}})\) has a convergent subsequence \(\left(y_{n_{i_{j}}}\right)\) which converges to a point \(y_{0} \in K\). Notice that since \(\left(x_{n_{i_{j}}}\right)\) is a subsequence of \((x_{n_{i}})\), it converges to \(x_{0}\). Since \(f\) is continuous:

    \[
        \lim_{j \to \infty} f\left(x_{n_{i_{j}}}\right) = f(x_{0}) \quad \text{and} \quad \lim_{j \to \infty} f\left(y_{n_{i_{j}}}\right) = f(y_{0}).
    \]
    Because
    \[
        \disabs{f\left(x_{n_{i_{j}}}\right) - f\left(y_{n_{i_{j}}}\right)} < \frac{1}{n_{i_{j}}},
    \]
    we can see that 
    \[
        \lim_{j \to \infty} \disabs{x_{n_{i_{j}}} - y_{n_{i_{j}}}} = 0.
    \]
    It follows that \(x_{0} = y_{0}\). But this is a contradiction because \(\disabs{f(x_{0}) - f(y_{0})} \geq \epsilon_{0}\). Therefore, \(f\) is uniformly continuous on \(K\).
}

\section{The Intermediate Value Theorem}

\begin{ntheorem}
    {Intermediate Value Theorem} Let \(f \colon [a,b] \to \R\) be continuous. If \(L\) is a real number satisfying \(f(a) < L < f(b)\) or \(f(b) < L < f(a)\), then there exists \(c \in (a,b)\) such that \(f(c) = L\).
\end{ntheorem}

\textbf{Note:} IVT does not guarantee where, or how many \(c\)'s are in the interval. It only guarantees that at least one \(c\) exists.

\tpf{
    Without the loss of generality, assume \(f(a) < f(b)\) and let \(y \in (f(a),f(b))\). Let \(I_{1} = [a_{1},b_{1}]\). Bisect \(I_{1}\) into two intervals \([a_{1},d]\) and \(d,b_{1}\) where \(d = \frac{a_{1} + b_{1}}{2}\). 
\begin{itemize}
    \item If \(f(d) < y\), set \(a_{2} = d\), \(b_{2} = b_{1}\), and \(I_{2} = [a_{2},b_{2}]\). Notice that \(f(a_{2}) < y < f(b_{2})\).
    \item If \(f(d) > y\), then set \(a_{2} = a_{1}\), \(b_{2} = d\), and \(I_{2} = [a_{2},b_{2}]\). Notice that \(f(a_{2}) < y < f(b_{2})\).
\end{itemize}
Repeat this process indefinitely. We end up with a sequence of nested intervals \(I_{n} = [a_{n},b_{n}]\), where 
\begin{itemize}
    \item \(I_{n} \subseteq I_{n-1}\)
    \item \(f(a_{n}) < y < f(b_{n})\)
    \item \(\disabs{a_{n} - b_{n}} = \dfrac{a_{n} - b_{n}}{2^{n-1}}\)
\end{itemize}
By the \namrefthm{Nested Interval Property}, there exists a point \(c\) such that \(c \in \bigcap_{n=1}^{\infty} I_{n}\). In fact, there is a unique point \(c\) in the intersection. It follows that 
\[
    c = \lim_{n \to \infty} a_{n} \quad \text{and} \quad c = \lim_{n \to \infty} b_{n}.
\]
Since \(f\) is continuous, we have
\[
    f(c) = \lim_{n \to \infty} f(a_{n}) \leq y \qquad \text{and} \qquad f(c) = \lim_{n \to \infty} f(b_{n}) \geq y.
\]
Therefore, \(f(c) = y\).
}

\begin{center}
    What Is Important About Continuous Functions?
\end{center}

If \(\limn(x_{n} = x)\), then \(\limn f(x_{n)} = f(x)\).

Think about \(f(x) = 2^{x}\). Thus, \(f(x)\) makes sense if \(x \in \Q\):

\[
    2^{\frac{p}{q}} = \sqrt[q]{2^{p}}.
\]
But how do we make sense of something like \(2^{\pi}\)? 

We can find \(f \colon \Q \to \R\) is continuous.

We can define \(f \colon \R \to \R\) to be continuous 

If \((q_{n})\) is in \(\Q\) and \((a_{n} \to \pi)\), then we define 
\[
    f(\pi) = \limn f(q_{n}).
\]







% \printglossary[title=Glossary]

%-%-%-%-%-%-%-%-%-%-%-%-%-%-%-%-%-%-%-%-%-%-%-%-%-%-%-%-%-%-%-%-%-%-%-%-%-%-%
%-%-%-%-%-%-%-%-%-%-%-%-%-%-%-%-%-%-%-%-%-%-%-%-%-%-%-%-%-%-%-%-%-%-%-%-%-%-%

% \chapter*{Exercises 1}\label{Exercises 1}
% \vspace*{-0.25in}
% \begin{exercise}
    {1.2.3}Decide which of the following represent true statements about
    the nature of sets. For any that are false, provide a specific example where the
    statement in question does not hold.
    \begin{enumerate}
        \item If \(A_1 \supseteq A_2 \supseteq A_3 \supseteq A_4\dots\) are all sets containing an infinite number of
              elements, then the intersection \(\bigcap^\infty_{n=1} A_n\) is infinite as well.
        \item If \(A_1 \supseteq A_2 \supseteq A_3 \supseteq A_4\dots\) are all finite, nonempty sets of real numbers,
              then the intersection \(\bigcap^\infty_{n=1} A_n\) is finite and nonempty.
        \item \(A \cap (B \cup C) = (A \cap B) \cup C\)
        \item \(A \cap (B \cap C) = (A \cap B) \cap C\)
        \item \(A \cap (B \cup C) = (A \cap B) \cup (A \cap C)\)
    \end{enumerate}
\end{exercise}

\sol{
    \begin{enumerate}
        \item This is false. Consider the following as a counterexample: If we define \(A_1\) as \(A_n = \{n, n+1, n+2,\dots\} = \{k \in \N \ | \ k \geq n\}\), we can see why the intersection of these sets of infinite numbers are actually empty. Consider a number \(m\) that actually satisfies \(m \in A_n\) for every \(A_n\) in our collection of sets. Because \(m\) is not an element of \(A_{m + 1}\), no such \(m\) exists and the intersection is empty.
        \item This is true.
        \item False. Consider sets \(A = \{1, 2, 3\}\), \(B = \{3, 6, 7\}\) and \(C = \{5\}\). Note that
              \(A \cap (B \cup C) = \{3\}\) is not equal to \((A \cap B) \cup C = \{3, 5\}\).
        \item This is true.
        \item This is true.
    \end{enumerate}
}

\begin{exercise}
    {1.2.5 (De Morgan's Laws)}(\hyperref[def:1.2.1]{This definition} may prove useful.) Let \(A\) and \(B\) be subsets of \(\R\).
    \begin{enumerate}
        \item  If \(x \in (A \cap B)^c\), explain why \(x \in A^c \cup B^c\). This shows that \((A \cap B)^c \subseteq A^c \cup B^c\).
        \item Prove the reverse inclusion \((A \cap B)^c \supseteq A^c \cup B^c\), and conclude that \((A \cap B)^c = A^c \cup B^c\).
        \item Show \((A \cup B)^c = A^c \cap B^c\) by demonstrating inclusion both ways.
    \end{enumerate}
\end{exercise}

\sol{
    \begin{enumerate}
        \item If \(x \in (A \cap B)^c\), and we know that \(A^c = \{x \in \R \ \colon x \notin A\}\), then we know \(x\) must cannot exist in \(A^c\) and \(B^c\) because \((A \cap B)^c = \{x \in \R \ \colon x \notin (A \cap B)\}\). Thus, \(x\) is in either \(A^c\) or \(B^c\). Put another way \(x \in A^c \cup B^c\). Since we have shown that an element that started in \((A \cap B)^c\) ended up in \(A^c \cup B^c\), then we know \((A \cap B)^c \subseteq A^c \cup B^c\).
        \item Assume  \(y \in A^c \cup B^c\). Thus, it must be the case that \(y \notin A\) or \(y \notin B\). Hence, \(y\) cannot be exist in both sets at the same time, so \(y \in (A \cap B)^c\). Because we have taken an element that started in \(A^c \cup B^c\) and have shown that it exists in \((A \cap B)^c\), we have proven \(A^c \cup B^c \subseteq (A \cap B)^c\).

              Therefore, since both inclusions hold, \((A \cap B)^{c} = A^{c} \cup B^{c}\).
        \item \sbseteqpf{
                  Let \(x \in (A \cup B)^c\).

                  By definition of the complement, this means \(x \notin A \cup B\).
                  Therefore, by definition of union, \(x \notin A\) and \(x \notin B\).
                  Thus, \(x \in A^c\) and \(x \in B^c\).
                  Hence, \(x \in (A^c \cap B^c)\).
                  So, \((A \cup B)^c \subseteq A^c \cap B^c\).}
              {Let \(x \in (A^c \cap B^c)\).

                  By definition of intersection, \(x \in A^c\) and \(x \in B^c\).
                  So, by definition of the complement \(x \notin A\) and \(x \notin B\).
                  Therefore, \(x \notin A \cup B\).
                  Hence, \(x \in (A \cup B)^c\).
                  Thus, \(A^c \cap B^c \subseteq (A \cup B)^c\).}
              {Since both inclusions hold, we have:

                  \[(A \cup B)^c = A^c \cap B^c.\qedhere\] }
              {xgray}
    \end{enumerate}
}



\begin{exercise}
    {1.2.7}Given a function \(f\) and a subset \(A\) of its domain, let \(f(A)\)
    represent the range of \(f\) over the set \(A\); that is, \(f(a) = \{f(x) \colon x \in A\}\).
    \begin{enumerate}
        \item Let \(f(x) = x^2\). If \(A = [0,2]\) (the closed interval \(\{x \in \R \colon 0 \leq x \leq 2\}\)) and \(B = [1,4]\), find \(f(A)\) and \(f(B)\). Does \(f(A\cap B) = f(A) \cap f(B)\) in this case? Does \(f(A \cup B) = f(A) \cup f(B)\)?
        \item Find two sets \(A\) and \(B\) for which \(f(A\cap B) \ne f(A) \cap f(B)\).
        \item Show that, for an arbitrary function \(g \ \colon \ \R \rightarrow \R\), it is always true that \(g(A \cap B) \subseteq g(A) \cap g(B)\) for all sets \(A,B\subseteq \R\).
        \item Form and prove a conjecture about the relationship between \(g(A \cup B)\) and \(g(A) \cup g(B)\) for an arbitrary function \(g\).
    \end{enumerate}
\end{exercise}

\sol{
    \begin{enumerate}
        \item Since \(f(x) = x^2\), the intervals of \(f(A)\) would be \([0,4]\) and \(f(B)\) would be \([1,16]\). The interval of the intersection of \(A \cap B\) is \([1,2]\). Take this through our function, we get \(f(A \cap B) = [1,4]\). On the other side of the equation, we already know the intervals of \(f(A)\) and \(f(B)\), and the intersection of theirs would be \([1,4]\). So they do equal each other. We know \(f(A \cup B)\) and \(f(A) \cup f(B)\) will be equivalent because \(f(A \cup B)\) has an interval of \([0,16]\), and \(f(A) \cup f(B)\) also has an interval of \([0,16]\) because taking the union of \([0,4] \cup [1,16]\) is \([0,16]\).
        \item Two sets could be \(A = [-1,0]\) and \(B = [0,1]\). For \(f(A \cap B)\), we have \(\{0\}\) but \(f(A) \cap f(B) = [0,1]\).
        \item \innerpf{
                  Let \(x \in g(A \cap B)\). Using the definition of function, we know there exists a \(y \in A \cap B\) to which that \(y\) is mapped to as \(g(y) = x\). From the definition of intersection, we know \(y \in A\) and \(y \in B\) such that \(x = g(y) \in g(A)\) and \(x = g(y) \in g(B)\) because \(y \in A \cap B\). Putting it together, we have \(x \in g(A) \cap g(B)\) thus proving \(g(A \cap B) \subseteq g(A) \cap g(B)\)
              }
        \item Conjecture: For any function \(g\) defined as \(g \ \colon \R \rightarrow \R\) and for any subsets \(A,B \subseteq \R\), it is always that case that 
        \[
            g(A \cup B) = g(A) \cup g(B).
        \]

              \sbseteqpf{
                  Let \(y \in g(A \cup B)\). 
                  
                  By definition of function, there exists \(x \in A \cup B\) such that \(g(x) = y\). From the definition of union, \(y \in A\) or \(y \in B\) and thus, \(g(x) \in g(A)\) or \(g(x) \in g(B)\). Together, we have \(g(x) = y \in g(A) \cup g(B)\). Therefore, \(g(A \cup B) \subseteq g(A) \cup g(B)\).
              }{
                  Let \(y \in g(A) \cup g(B)\).

                  By definition of union, \(y \in g(A)\) or \(y \in g(B)\). By definition of function, this means there exists an \(x \in A\) or \(x \in B\) such that \(g(x) = y\). Hence, \(x \in A \cup B\) by the definition of union, and \(g(x) \in g(A \cup B)\). Therefore, \(g(A) \cup g(B) \subseteq g(A \cup B)\).
              }{
                  Since we have shown the inclusion for both directions, this proves that \(g(A \cup B) = g(A) \cup g(B)\).
              }%
              {xgray}
    \end{enumerate}
}

\begin{exercise}
    {1.2.8}Given a function \(f \ \colon A \rightarrow B\) can be defined as either \hyperref[def:1.2.5]{injective} or \hyperref[def:1.2.6]{surjective}, give an example of each or state that the request is impossible:
    \begin{enumerate}
        \item \(f \ \colon \N \rightarrow \N\) that is 1-1 but not onto.
        \item \(f \ \colon \N \rightarrow \N\) that is onto but not 1-1.
        \item \(f \ \colon \N \rightarrow \Z\) that is 1-1 and onto.
    \end{enumerate}
\end{exercise}

\sol{
    \begin{enumerate}
        \item \textbf{Possible.} Define \(f \colon \N \rightarrow \N\) as \(f(a) = a + 1\). This function is injective because when \(f(a_{1}) = f(a_{2})\), we have:
        \begin{align*}
            a_{1} + 1 &= a_{2} + 1 \\
            a_{1} &= a_{2}.
        \end{align*}
        However, this function is not surjective because the entire co-domain is not covered; that being 1.
        \item \textbf{Possible.} Consider the example: 
        \[
            f(1) = 1, f(2) = 1, f(3) = 2, f(4) = 3, \ldots.
        \] 
        This function is surjective because for every \(b\) in the co-domain there is an \(a\) in the domain. However, it is not injective because we have two distinct values, 1 and 2, that map to the same \(b\), 1.
        \item \textbf{Possible.} Consider the example:
        \[
            f(1) = 0, f(2) = 1, f(3) = -1, f(4) = 2, f(5) = -2, \ldots.
        \]
        Since every value \(b\) in the co-domain is mapped to an element in  the domain, this function is surjective. Additionally, because that mapping is unqiue (i.e., \(a_{1} \ne a_{2} \implies f(a_{1}) \ne f(a_{2})\)), this function is injective by definition.
    \end{enumerate}
}

% \chapter*{Exercises 2}\label{Exercises 2}
% \vspace*{-0.25in}
% 
\begin{exercise}
    {1.3.4}Let \(A_1,A_2,A_3\dots\) be a collection of nonempty sets each of which is bounded above.
    \begin{enumerate}
        \item Find a formula for \(\sup(A_1 \cup A_2)\). Extend this to \(\sup (\bigcup^n_{k=1}A_k)\).
        \item Consider \(\sup(\bigcup^\infty_{k=1} A_k)\). Does the formula in (a) extend to the infinite case?
    \end{enumerate}
\end{exercise}

\sol{\hfill
    \begin{enumerate}
        \item Let \(A_1\) and \(A_2\) be nonempty sets, each bounded above. To find the largest of the two suprema, we can use the following: \(\sup(A_1 \cap A_2) = \max\{\sup A_1, \sup A_2\}\). If we extend this notion to \(\sup (\bigcup^n_{k=1}A_k)\), we can use the same idea from before and write it as \(\sup (\bigcup^n_{k=1}A_k) = \max\{\sup A_1, \sup A_2, \dots, \sup A_n\}\).
        \item The formula does not extend to the infinite case. Consider the counterexample \(\bigcup^\infty_{k=1} A_k\) where \(A_k := [k, k+1]\). Even though these sets are bounded above, when we take the union of them, we approach infinity, which is not bounded: \(\bigcup^\infty_{k=1} A_k = [1,2] \cup [2,3] \cup \dots = [1, \infty)\).
    \end{enumerate}
}

\begin{exercise}
    {1.3.5}As in Example 1.3.7, let \(A \subseteq \R\) be nonempty and bounded above, and let \(c \in \R\). This time define the set \(cA = \{ca \ \colon a \in A\}\).
    \begin{enumerate}
        \item If \(c \geq 0\), show that \(\sup(cA) = c\sup A\).
        \item Postulate a similar type of statement for \(\sup(cA)\) for the case \(c < 0\).
    \end{enumerate}
\end{exercise}

\sol{
    \begin{enumerate}
        \item Let \(A \subseteq \R\) be nonempty and bounded above. Define the set \(cA := \{ca \ \colon a \in A\}\). From the axiom of completeness, because \(A\) is bounded above, we know there is a least upper bound, \(s = \sup A\). Following from Example 1.3.7, we see that \(a \leq s\) for all \(a \in A\) which implies \(ca \leq cs\) for all \(a \in A\). Thus, \(cs\) is an upper bound for \(cA\), and the first condition of Definition 1.3.2 is satisfied. For the second condition, we need to look at both \(c = 0\) and \(c > 0\) to avoid dividing by zero. So, we have two cases:
              \begin{itemize}
                  \item \textbf{\(c = 0\):} If \(c = 0\), then \(cA = \{0 \colon \ a \in A\} = \{0\}\). Since the only element in \(cA\) is \(0\), \(\sup(cA) = 0\). Similarly, because \(c = 0\), \(c \sup A = 0 \cdot \sup A = 0\). Therefore, \(\sup(cA) = c \sup(A)\).
                  \item \textbf{\(c > 0\):} Let \(b\) be an arbitrary upper bound for \(cA\) and \(c > 0\). In other words, \(ca \leq b\) for all \(a \in A\). This is equivalent to \(a \leq b / c\) where \(c \ne 0\), from which we can see that \(b / c\) is an upper bound for \(A\). Because \(s\) is the least upper bound of \(A\), \(s \leq b / c\), which can be rewritten as \(cs \leq b\). This verifies the second part of Definition 1.3.2, and we conclude \(\sup(cA) = c \sup A\).
              \end{itemize}
        \item Postulate: If \(c < 0\), then \(\sup(cA) = c \inf (A)\). \\
              % \expf{Assume \(c < 0\), Multiplying \(a \leq s\) by \(c\) reverses the inequality, so \(ca \geq cs\) for all \(a \in A\). Therefore, \(cs\) is a lower bound of for \(cA\). To verify this as the lowest upper bound, let \(m  = \inf A\). Using the same logic as before, for all \(a \in A\), \(m \leq a\), so \(cm \geq ca\). This shows that \(cm\) is the least upper bound of \(cA\), and thus \(\sup(cA) = cm = c \inf A\).}
    \end{enumerate}
}

\begin{exercise}
    {1.3.8}Compute, without proofs, the suprema and infima (if they exist) of the following sets:
    \begin{enumerate}
        \item \(\left\{\frac{m}{n} : m, n \in \mathbb{N} \text{ with } m < n\right\}\).
        \item \(\left\{\frac{{(-1)}^{m}}{n} : m, n \in \mathbb{N}\right\}\).
        \item \(\left\{\frac{n}{3n+ 1} : n \in \mathbb{N}\right\}\).
        \item \(\left\{\frac{m}{m+ n} : m, n \in \mathbb{N}\right\}\).
    \end{enumerate}
\end{exercise}

\sol{To avoid writing out every set definition, I am going to denote each set as \(A_n\) where \(n\) corresponds to the numerical value of the list from (a) - (d).
    \begin{enumerate}
        \item \(\sup A_1 = 1\), \(\inf A_1 = 0\)
        \item \(\sup A_2 = 1\), \(\inf A_2 = -1\)
        \item \(\sup A_3 = \frac{1}{3}\), \(\inf A_3 = \frac{1}{4}\)
        \item \(\sup A_4 = 1\), \(\inf A_3 = 0\)
    \end{enumerate}
}

\begin{exercise}
    {1.4.1}Recall that \(\mathbb{I}\) stands for the set of irrational numbers.
    \begin{enumerate}
        \item Show that if \(a,b \in \Q\), then \(ab\) and \(a + b\) are elements of \(\Q\) as well.
        \item Show that if \(a \in \Q\) and \(t \in \mathbb{I}\), then \(a + t \in \mathbb{I}\) and \(at \in \mathbb{I}\) as long as \(a \ne 0\).
        \item Part (a) can be summarized by saying that \(\Q\) is closed under addition and multiplication. Is \(\mathbb{I}\) closed under addition and multiplication? Given two irrational numbers \(s\) and \(t\), what can we say about \(s + t\) and \(st\)? In other words, are there two irrational numbers that can be added and multiplied such that you get a number \(x\) such that \(x \notin \mathbb{I}\).
    \end{enumerate}
\end{exercise}

\sol{
    \begin{enumerate}
        \item Let \(a,b \in \Q\). This means there exists some \(p,q,a,b \in \Z\) such that \[a = \frac{p}{q}\] and \[b = \frac{a}{b}\] where \(q,b \ne 0\). The product of these numbers is \[ab = \frac{p}{q} \cdot \frac{a}{b} = \frac{pa}{qb}.\] Since \(pa,qb \in \Z\), \(ab \in \Q\). The sum of these numbers is \[a + b = \frac{p}{q} + \frac{a}{b} = \frac{pb + aq}{qb}.\] Since \(pb + aq, qb \in \Z\), \(a + b \in \Q\).
        \item Let \(a \in \Q\) and \(t \in \I\). Assume, for contradiction, that \(a + t \in \Q\). This would imply \(t = (a + t) - a\) (because we can subtract \(t + a\) from the original equation and rearrange terms). Since \(a + t, a \in \Q\) their sum would be rational because the rational numbers are closed under addition. However, that would contradict the assumption that \(t \in \I\). Hence, \(a + t \in \I\).
        \item For \(\I\), it is not closed under addition and multiplication. Consider the following counterexample: \(\sqrt{2} + (-\sqrt{2}) = 0\) which is not in the irrationals. For multiplication, consider \(\sqrt{2} \cdot \sqrt{2} = 2\), which is also not in the irrationals.
    \end{enumerate}
}


% \fancyhead[R]{Exercises 3}

% \chapter*{Exercises 3}\label{Exercises 3}
% \vspace*{-0.25in}
% \begin{exercise}
    {2.2.2(b)}Verify, using \defref{2.2.3}, that
    the following sequences converge to the proposed limit.
    \begin{enumerate}
        \setcounter{enumi}{1}
        \item \(\limn{\frac{2n^2}{n^3+3}} = 0\)
    \end{enumerate}
\end{exercise}

\expf{\hfill
    \begin{enumerate}
        \setcounter{enumi}{1}
        \item Let \(\epsilon > 0\). Choose \(N = \frac{2}{\epsilon}\). Then, for \(n > N\),
              \begin{align*}
                  \disabs{\frac{2n^2}{n^3 + 3} - 0} & = \disabs{\frac{2n^2}{n^3 + 3}} \\
                                                    & = \frac{2n^2}{n^3 + 3}          \\
                                                    & < \frac{2n^2}{n^3}              \\
                                                    & = \frac{2}{n}                   \\
                                                    & < \frac{2}{N}                   \\
                                                    & = \frac{2}{2/\epsilon}          \\
                                                    & = \epsilon.
              \end{align*}
              Therefore, \(\limn{\frac{2n^2}{n^3 + 3}} = 0\). \qedhere
    \end{enumerate}
}

\begin{exercise}
    {2.2.3} Describe what we would have to demonstrate in order to disprove each of the following statements.
    \begin{enumerate}
        \item At every college in the United States, there is a student who is at least
              seven feet tall.
        \item For all colleges in the United States, there exists a professor who gives every student a grade of either A or B.
        \item There exists a college in the United States where every student is at least six feet tall.
    \end{enumerate}
\end{exercise}

\sol{
    \begin{enumerate}
        \item There is at least one college in the United States where all students are less than seven feet tall.
        \item There is at least one college in the United States where all professors give at least one student a grade of C or lower.
        \item For all colleges in the United States, there exists a student who is less than six feet tall.
    \end{enumerate}
}

\begin{exercise}
    {2.2.4} Give an example of each or state that the request is impossible.
    For any that are impossible, give a compelling argument for why that is the case.
    \begin{enumerate}
        \item A sequence with an infinite number of ones that converges to one.
        \item A sequence with an infinite number of ones that converges to a limit not equal to one.
        \item A divergent sequence such such that for every \(n \in \N\) it is possible to find \(n\) consecutive ones somewhere in the sequence.
    \end{enumerate}
\end{exercise}

\sol{
    \begin{enumerate}
        \item Possible. Example: \(a_n = 1 \text{ for all } n \in \N\).
        \item Impossible. We can use the \namrefthm{Uniqueness of Limits} to show that this is impossible. Essentially, because we found in (a) that the limit of an infinite number of ones converges to 1, it must be the case that if the limit is not equal to one, then the sequence cannot converge to one.
        \item Possible. Example: \(\{2,2,2,2,4,4,4,4,1,1,1,1,1,3,3,\dots\}\)
    \end{enumerate}
}


\begin{exercise}
    {2.3.2}Using only \defref{2.2.3}, prove that if \((x_n) \rightarrow 2\), then
    \begin{enumerate}
        \item \(\displaystyle \left(\frac{2x_n - 1}{3}\right) \rightarrow 1\);
        \item \(\displaystyle \left(1 / x_n\right) \rightarrow 1 / 2\).
    \end{enumerate}
    (For this exercise the Algebraic Limit Theorem is off-limits, so to speak.)
\end{exercise}

\sol{ \hfill
    \begin{enumerate}
        \item \textit{Proof.} Let \(\epsilon > 0\). Since \((x_n)\) converges to 2, there exists \(N \in \N\) such that for all \(n \geq N\), \(\disabs{x_n - 2} < \epsilon\). Now, for any \(n \geq N\),
              \begin{align*}
                  \disabs{\frac{2x_n - 1}{3} - 1} & = \disabs{\frac{2x_n - 1 - 3}{3}} \\
                                                  & = \disabs{\frac{2x_n - 4}{3}}     \\
                                                  & = \frac{2}{3}\disabs{x_n - 2}     \\
                                                  & < \disabs{x_n - 2}                \\
                                                  & < \epsilon
              \end{align*}
              Therefore, \(\frac{2x_n - 1}{3} \rightarrow 1\) \qed
        \item \textit{Proof.} Let \(\epsilon > 0\). Since \((x_n)\) converges to 2, there exists \(N \in \N\) such that for all \(n \geq N\), \(\disabs{x_n - 2} < \epsilon\). Now, for any \(n \geq N\),
              \begin{align*}
                  \disabs{\frac{1}{x_n} - \frac{1}{2}} & = \disabs{\frac{2 - x_n}{2x_n}}  \\
                                                       & = \frac{1}{2x_n}\disabs{2 - x_n} \\
              \end{align*} \qed
    \end{enumerate}
}

% \chapter*{Exercises 4}\label{ex 4}
% \vspace*{-0.25in}
% \begin{exercise}
    {2.5.1}Give an example of each of the following, or argue that such a request is impossible.
    \begin{enumerate}
        \item A sequence that has a subsequence that is bounded but contains no subsequence that converges.
        \item A sequence that does not contain 0 or 1 as a term but contains subsequences converging to each of these values.
    \end{enumerate}
\end{exercise}

\sol{
    \begin{enumerate}
        \item Impossible. This violates the \namrefthm{Bolzano-Weierstrass Theorem}. It assures us that every bounded sequence has a convergent subsequence. If a subsequence is bounded, then it must have a convergent subsequence.
        \item Consider the sequence \((\frac{1}{2},\frac{1}{2},\frac{1}{3},\frac{2}{3},\frac{1}{4},\frac{3}{4}, \cdots \frac{1}{n}, \frac{(n-1)}{n})\). From this, you can have a subsequence \((\frac{1}{2}, \frac{1}{3}, \dots, \frac{1}{n})\) which converges to 0, and also a subsequence \((\frac{1}{2}, \frac{2}{3}, \dots \frac{n - 1}{n})\), which converges to 1.
    \end{enumerate}
}

\begin{exercise}
    {2.5.2} Decide whether the following propositions are true or false, providing a short justification for each conclusion.
    \begin{enumerate}
        \item If every proper subsequence of \((x_n)\) converges, then \((x_n)\) converges as well.
              \setcounter{enumi}{2}
        \item If \((x_n)\) is bounded and diverges, then there exist two subsequences of \((x_n)\) that converge to different limits.
    \end{enumerate}
\end{exercise}

\sol{
    \begin{enumerate}
        \item False. As shown in Example 2.5.4, if we have a sequence like \((1, -\frac{1}{2}, \frac{1}{3}, -\frac{1}{4}, \frac{1}{5}, -\frac{1}{5}, \frac{1}{5}, \cdots)\), we see that it is divergent, and has two proper subsequences \((\frac{1}{5},\frac{1}{5}, \cdots)\) and \((-\frac{1}{5}, -\frac{1}{5}, \cdots)\) that both converge to values \(\frac{1}{5}\) and \(-\frac{1}{5}\), respectively.
              \setcounter{enumi}{2}
        \item True. If we assume that \((x_n)\) is bounded and diverges, then by the \namrefthm{Bolzano-Weierstrass Theorem}, every bounded sequence has a convergent subsequence. Let \((x_{n_k})\) be that convergent subsequence of \((x_n)\), and let \(L_1\) be its limit. Since \((x_n)\) diverges, it cannot converge to \(L_1\). Thus, there exists an \(\epsilon_0 > 0\) such that for all terms of \((x_n)\), they stay outside the \(\epsilon_0\)-neighborhood of \(L_1\). Assume this sub-subsequence \((x_{m_k})\) contains these terms. When we apply the same logic to this sub-subsequence, we see that by the \namrefthm{Bolzano-Weierstrass Theorem}, \((x_{m_k})\) has a convergent subsequence with limit \(L_2\), where \(L_2 \ne L_1\) because the terms of \((x_{m_k})\) stay outside the \(\epsilon_0\)-neighborhood of \(L_1\). Thus \((x_n)\) contains two subsequences that converge to different limits, \(L_1\) and \(L_2\).
    \end{enumerate}
}

\begin{exercise}
    {2.5.5} Assume \((a_n)\) is a bounded sequence with the property that every convergent subsequence of \((a_n)\) converges to the same limit \(a \in \mathbb{R}\). Show that \((a_n)\) must converge to \(a\).
\end{exercise}

\expf{
Suppose that \((a_n)\) does not converge to \(a \in \R\). By the definition of convergence, this means there is a positive real number \(\epsilon_0\) such that no matter how large we choose \(N \in \N\), there will always exist some \(n > N\) where \(|a_n - a| \geq \epsilon_0\). In a formal way, this shows that \((a_n)\) does not converge to \(a\) within the \(\epsilon_0\)-neighborhood.

We aim to demonstrate that this leads to a contradiction by constructing a subsequence of \((a_n)\) that stays outside this neighborhood. Begin by selecting \(n_1\) such that \(|a_{n_1} - a| \geq \epsilon_0\). Next, since the condition holds for all \(N \in \N\), we can find another index \(n_2 > n_1\) such that \(|a_{n_2} - a| \geq \epsilon_0\). Continuing this process, we generate an increasing sequence of indices \(n_1 < n_2 < n_3 < \dots\) such that for each \(i \in \N\), \(|a_{n_i} - a| \geq \epsilon_0\).

Now consider the subsequence \((a_{n_i})\) we have built. Since \((a_n)\) is bounded by assumption, its subsequence \((a_{n_i})\) is also bounded. By the \namrefthm{Bolzano-Weierstrass Theorem}, every bounded sequence has a convergent subsequence. Let \((a_{n_{i_k}})\) denote a convergent subsequence of \((a_{n_i})\). According to our assumption, any convergent subsequence of \((a_n)\) must converge to \(a\).

However, each term of \((a_{n_{i_k}})\) remains outside the \(\epsilon_0\)-neighborhood of \(a\). Thus, it is impossible for \((a_{n_{i_k}})\) to converge to \(a\). This contradiction implies that our initial assumption—that \((a_n)\) does not converge to \(a\)—is false. Therefore, the sequence \((a_n)\) must converge to \(a\).
}



\begin{exercise}
    {2.5.6} Use a similar strategy to the one in \numrefthm{2.5.5} to show
    \[\lim b^{1/n} \text{ exists for all } b \geq 0\] and find the value of the limit. (The results in \hyperref[ex:2.3.1]{Exercise 2.3.1} may be assumed.)
\end{exercise}

\expf{
    For this proof, we will be examining different cases that \(b\) may fall under. These cases are \(b = 0\) and \(b > 0\).
    \begin{enumerate}[label=\arabic*.]
        \item If \(b = 0\), then the sequence \((b^{1/n})\) is simply \((0, 0, 0, \dots)\). Thus, the limit of this sequence is 0.
        \item For \(b > 0\), it is more complicated. We see that this sequence is decreasing because as \(n\) increases, the exponent \(\frac{1}{n}\) trends toward a higher denominator, so \(b^{1/n}\) decreases. We can see this by observing the property that for any \(m < n\), \(\frac{1}{n} < \frac{1}{m}\). This implies that \(b^{1/n} < b^{1/m}\).

              Further, we see that this sequence is bounded below by \(0\) when \(b > 1\) and by 1 when \(b = 1\) (because \(1\) to any exponent is just 1). Therefore, this sequence is monotone and bounded below, so it must converge to some \(L \geq 0\) by the \namrefthm{Monotone Convergence Theorem}.

              To find \(L\), we need to consider another set of cases for \(b\). \begin{enumerate}[label = \roman*.]
                  \item If \(b = 1\), then \(b^{1/n} = 1\) for all \(n\), so its limit is 1.
                  \item If \(b > 1\), then as \(n\) increases, the exponent \(\frac{1}{n}\) approaches 0, so \(b^{1/n}\) approaches \(b^0 = 1\). So, its limit is also 1.
                  \item If \(0 < b < 1\), then we know \(\limn b^{1/n} = 1\) because \(b^{1/n}\) is a decreasing sequence bounded by 0, and when we have smaller and smaller powers of a number less than 1, it pushes it closer to 1.
              \end{enumerate}
    \end{enumerate}
    Through these cases, we see that for all \(b \geq 0\), \(\limn b^{1/n} = 1\).
}


% \chapter*{Exercises 5}\label{ex 5}
% \vspace*{-0.25in}
% \fancyhead[R]{Exercises 5}
% \begin{exercise}
    {3.2.4}Let \(A\) be a nonempty and bounded above set so that \(s = \sup(A)\) exists. 
    \begin{enumerate}
        \item Show that \(s \in \bar{A}\).
        \item Can an open set contain its supremum?
    \end{enumerate}
\end{exercise}

\sol{
    \begin{enumerate}
        \item e
        \item e
    \end{enumerate}
}

\begin{exercise}
    {3.2.6}Decide whether the following statements are true or false. Provide counter examples for those that are false, and supply proofs for those that are true. (Choose 2.)
    \begin{enumerate}
        \item An open set that contains every rational number must necessarily be all of \(\R\).
        \item The \namrefthm{Nested Interval Property} remains true if the term ``closed interval'' is replaced by ``closed set.''
        \item Every nonempty open set contains a rational number.
        \item Every bounded infinite closed set contains a rational number.
        \item The Cantor set (\defref{3.1.1}) is closed. 
    \end{enumerate}
\end{exercise}

\sol{
    \begin{enumerate}
        \item e
        \item e
    \end{enumerate}
}

\begin{exercise}
    {3.2.8} Assume \(A\) is an open set and \(B\) is a closed set. Determine if the following sets are definitely open, definitely closed, both or neither.
    \begin{enumerate}
        \item \(\overline{A \cup B}\)
        \item \(A \setminus B = \{x \in A \mid x \notin B\}\)
        \item \((A^c \cup B)^c\)
        \item \((A \cap B) \cup (A^c \cap B)\)
        \item \(\overline{A}^c \cap \overline{A^c}\)
    \end{enumerate}
\end{exercise}

\sol{
    \begin{enumerate}
        \item e
        \item e
        \item e
        \item e
        \item ee
    \end{enumerate}
}

\begin{exercise}
    {3.2.11} \begin{enumerate}
        \item Prove that \(\overline{A \cup B} = \overline{A} \cup \overline{B}\).
        \item Does this result about closures extend to infinite unions of sets?
    \end{enumerate}
\end{exercise}

\sol{
    \begin{enumerate}
        \item \innerpf{}
        \item NO (Refuses to elaborate).
    \end{enumerate}
}

\begin{exercise}
    {3.3.4}Assume \(K\) is compact and \(F\) is closed. Decide if the following sets are definitely compact, definitely closed, both, or neither. 
    \begin{enumerate}
        \item \(K \cap F\)
        \item \(\overline{F^c \cup K^c}\)
        \item \(K \setminus F = \{x \in K \mid x \notin F\}\)
        \item \(\overline{K \cap F^c}\)
    \end{enumerate}
\end{exercise}

% \chapter*{Presentation Problem 1}\label{pres 1}
% \vspace*{-0.25in}
% 
\begin{exercise}
    {1.2.13}For this exercise, assume \hyperref[ex:1.2.5]{Exercise 1.2.5} has been successfully completed.
    \begin{enumerate}[label=(\alph*)]
        \item Show how induction can be used to conclude that
        \[
        (A_1 \cup A_2 \cup \cdots \cup A_n)^c = A_1^c \cap A_2^c \cap \cdots \cap A_n^c
        \]
        for any finite \(n \in \mathbb{N}\).
    
        \item It is tempting to appeal to induction to conclude
        \[
        \left( \bigcup_{i=1}^{\infty} A_i \right)^c = \bigcap_{i=1}^{\infty} A_i^c,
        \]
        but induction does not apply here. Induction is used to prove that a particular statement holds for every value of \(n \in \mathbb{N}\), but this does not imply the validity of the infinite case. To illustrate this point, find an example of a collection of sets \(B_1, B_2, B_3, \dots\) where
        \[
        \bigcap_{i=1}^{n} B_i \neq \emptyset \quad \text{is true for every } n \in \mathbb{N},
        \]
        but
        \[
        \bigcap_{i=1}^{\infty} B_i = \emptyset
        \]
        fails.
    
        \item Nevertheless, the infinite version of De Morgan’s Law stated in (b) is a valid statement. Provide a proof that does not use induction.
    \end{enumerate}
\end{exercise}
\newpage
\sbseteqpf%
%
[In this proof, we plan to prove (c). Thus, we need to show that:
\[
\left( \bigcup_{i=1}^{\infty} A_i \right)^c \subseteq \bigcap_{i=1}^{\infty} A_i^c
\]
and
\[
\left( \bigcup_{i=1}^{\infty} A_i \right)^c \supseteq \bigcap_{i=1}^{\infty} A_i^c.
\]]%
%
{Let \(x \in \left( \bigcup_{i=1}^{\infty} A_i \right)^c\). This means \(x\) is in the union set of \(A_i\) for all \(i \in \N\). Then, because we are taking the complement of \(\left( \bigcup_{i=1}^{\infty} A_i \right)\), that means \(x \notin A_i\) for all \(i \in \mathbb{N}\). Hence, \(x\) is in the complement of each \(A_i\). Thus, we can use the definition of intersection to assert \(x \in \bigcap_{i=1}^{\infty} A_i^c\).

Therefore, we have shown:
\[
\left( \bigcup_{i=1}^{\infty} A_i \right)^c \subseteq \ \bigcap_{i=1}^{\infty} A_i^c.
\]}%
%
%
%
{Similar to before, let \(x \in \bigcap_{i=1}^{\infty} A_i^c\). Because \(x \in A^c_i\) for all \(i \in \N\) we know \(x \notin A_i\). Hence, \(x \notin \left(\bigcup_{i=1}^{\infty} A_i\right)\), which means \(x \in \left( \bigcup_{i=1}^{\infty} A_i \right)^c\).

Therefore, we have shown:
\[
\left( \bigcup_{i=1}^{\infty} A_i \right)^c \supseteq \ \bigcap_{i=1}^{\infty} A_i^c.
\]}%
%
%
{By showing both inclusions, we see that:
\[
\left( \bigcup_{i=1}^{\infty} A_i \right)^c = \bigcap_{i=1}^{\infty} A_i^c.
\] \qedhere
}%
%
%
{xgray}


% \addcontentsline{toc}{chapter}{Presentation Problem 1}

% \chapter*{Presentation Problem 2}\label{pres 2}
% \vspace*{-0.25in}
% 

\begin{exercise}
    {2.2.7} {Here are two useful definitions:}
    \begin{enumerate}[label=(\roman*)]
        \item A sequence \((a_n)\) is \emph{eventually} in a set \(A \subseteq \mathbb{R}\) if there exists an \(N \in \mathbb{N}\) such that \(a_n \in A\) for all \(n \geq N\).
        \item A sequence \((a_n)\) is \emph{frequently} in a set \(A \subseteq \mathbb{R}\) if, for every \(N \in \mathbb{N}\), there exists an \(n \geq N\) such that \(a_n \in A\).
        \begin{enumerate}[label=(\alph*)]
            \item Is the sequence \((-1)^n\) eventually or frequently in the set \(\{1\}\)?
            \item Which definition is stronger? Does frequently imply eventually, or does eventually imply frequently?
            \item Give an alternate rephrasing of \defref{2.2.3B} using either frequently or eventually. Which is the term we want?
            \item Suppose an infinite number of terms of a sequence \((x_n)\) are equal to \(2\). Is \((x_n)\) necessarily eventually in the interval \((1.9, 2.1)\)? Is it frequently in \((1.9, 2.1)\)?
        \end{enumerate}
    \end{enumerate}

\end{exercise}

\sol{
    \begin{enumerate}[label=(\alph*)]
        \item The sequence \((-1)^n\) is \textit{frequently} in the set \(\{1\}\) because for every \(N \in \mathbb{N}\), we can find an \(n \geq N\) such that \((-1)^n = 1\).
        \item The definition of \textit{eventually} is stronger because \textit{eventually} implies \textit{frequently}, but \textit{frequently} does not imply \textit{eventually}.
        \item An alternate rephrasing of Definition 2.2.3B using \textit{eventually} is: A sequence \((a_n)\) converges to \(a\) if, given any \(\epsilon\)-neighborhood---\(V_\epsilon(a)\) of \(a\)---\(a_n\) is \textit{eventually} in \(V_\epsilon(a)\). The term we want is eventually.
        \item If an infinite number of terms of a sequence \((x_n)\) are equal to \(2\),  \((x_n)\) is not \textit{eventually} in \((1.9, 2.1)\) because we can have a sequence \((a_n)\) that will not settle in \((1.9, 2.1)\). For example, \((a_n) = (0,2,0,2,\cdots)\) does not settle in \((1.9, 2.1)\). Whereas, \((x_n)\) is \textit{frequently} in the interval \((1.9, 2.1)\) because for every \(N \in \mathbb{N}\) there exists an \(n \geq N\) such that \(x_n \in (1.9, 2.1)\) for all \(n \geq N\). We can see an instance of this being true by examining the previous example. 
    \end{enumerate}
}

\begin{tcolorbox}%
    [rounded corners, colframe=xred, colback=xred!15, coltitle=white, label=def:2.2.3B, title={\large \textbf{Definition 2.2.3B}}]
    A sequence \( (a_n) \) converges to \( a \) if, given any \( \epsilon \)-neighborhood \( V_\epsilon(a) \) of \( a \), there exists a point in the sequence after which all of the terms are in \( V_\epsilon(a) \). In other words, every \( \epsilon \)-neighborhood contains all but a finite number of the terms of \( (a_n) \).
\end{tcolorbox}
% \addcontentsline{toc}{chapter}{Presentation Problem 2}

%-%-%-%-%-%-%-%-%-%-%-%-%-%-%-%-%-%-%-%-%-%-%-%-%-%-%-%-%-%-%-%-%-%-%-%-%-%-%
\end{document}
%-%-%-%-%-%-%-%-%-%-%-%-%-%-%-%-%-%-%-%-%-%-%-%-%-%-%-%-%-%-%-%-%-%-%-%-%-%-%

%-%-%-%-%-%-%-%-%-%-%-%-%-%-%-%-%-%-%-%-%-%-%-%-%-%-%-%-%-%-%-%-%-%-%-%-%-%-%